% !TeX root = ../大学生数学竞赛概率论与数理统计试题选编与专题训练.tex
% 由主文档通过 \input 载入；请编译主文档。
\setcounter{chapter}{5} % 保留原书章号。
\chapter{解题方法与公式速查}
\section*{方法一：先定义试验，再计算概率}
\addcontentsline{toc}{section}{方法一：先定义试验，再计算概率}
概率空间由样本空间、事件集合和概率规则组成。本科计算题常把这三者写在自然语言中，第一步应把它们翻译为明确的数学对象。例如“任取两件产品”必须区分有放回和无放回，“随机取两点”必须区分每点均匀、彼此独立以及相对于长度、面积还是体积均匀。若这些条件尚未确定，仅凭“随机”二字不能选定分布。

古典概率中的分子和分母必须在同一种等可能基本结果上计数。抽取有先后次序时，可以把每个有序序列视为一个基本结果；也可以转成无序子集，但分子与分母应同时转化。若某些无序结果对应的序列数不同，无序结果并不等可能。例如两次掷骰子的和有十一种可能值，却不能将每个和的概率都记为十一分之一。

组合数 $\binom nk$ 适合无序选择，排列数 $n!/(n-k)!$ 适合无放回的有序选择。有放回的 $k$ 次独立选择通常产生 $n^k$ 个有序结果。处理“恰好”时特别注意不要把“至少”计数带入；手套题中选中一副完整手套后，剩余两只不能再来自同一副。

遇到“至少一个”时先试补事件。若 $A_i$ 表示第 $i$ 个对象满足要求，通常 $P(\bigcup_iA_i)$ 比 $P(\bigcap_iA_i^c)$ 难计算。独立性成立时后者是乘积；独立性不成立时，仍可使用条件概率连乘，但不能直接相乘边缘概率。容斥公式适合许多对称事件的交集容易计数的情形，如帽子、信封与占位问题。
\begin{equation*}
P\left(\bigcup_{i=1}^nA_i\right)
=\sum_{\varnothing\ne J\subset\{1,\ldots,n\}}(-1)^{|J|-1}P\left(\bigcap_{j\in J}A_j\right).
\end{equation*}
计算后至少检查三个边界：概率是否在 $[0,1]$；样本容量为零或一时是否合乎题意；参数趋近零或一时结果是否具有合理极限。这些检查不能替代推导，却常能发现漏项、错位和重复计数。对应训练为习题~\ref{q:Y1.9}、\ref{q:Y2.3}、\ref{q:Y2.5}。

\section*{方法二：条件概率与信息产生机制}
\addcontentsline{toc}{section}{方法二：条件概率与信息产生机制}
条件概率（conditional probability）在 $P(B)>0$ 时定义为 $P(A\mid B)=P(A\cap B)/P(B)$。分母改变了参考总体：我们只在已知 $B$ 发生的情形中比较 $A$。因此 $P(A\mid B)$ 与 $P(B\mid A)$ 通常不同。检测、证人和分类问题中，题目常给后者却要求前者，此时需要贝叶斯公式。

若 $H_1,\ldots,H_k$ 构成互斥完备假设，先验为 $\pi_j$，观察数据 $D$ 的似然为 $L_j=P(D\mid H_j)$，则后验为
\[
P(H_j\mid D)=\frac{\pi_jL_j}{\sum_i\pi_iL_i}.
\]
实际解题可先计算未归一化权重 $w_j=\pi_jL_j$，最后统一除以权重和。这样避免对每个假设重复计算全概率分母，也便于发现哪些假设的似然为零。两假设情形可使用赔率：后验赔率等于先验赔率乘以似然比。

连续观测没有非零单点概率，但有密度似然。给定观测 $x$ 后，$L(\theta;x)=f_\theta(x)$ 是参数的函数，不要求对参数积分等于一。贝叶斯后验则还需乘先验密度，再对参数归一化。将似然本身视为参数的概率分布，只有在明确采用相应先验且归一化有意义的特定情形下才可能成立。

连续检测需要检查的是给定真实状态后的条件独立，而不是检测结果的无条件独立。若两个结果由同一个未知状态共同影响，即使它们条件独立，边际上也往往相关。三门问题进一步说明：观察到“某门被打开且里面没有奖品”，其似然取决于主持人如何选门。信息并不只在看到的结果里，也在产生结果的规则里。对应训练为习题~\ref{q:Y2.23}、\ref{q:M-P7-1}、\ref{q:M-P9-2}。

\section*{方法三：分布函数比密度更基本}
\addcontentsline{toc}{section}{方法三：分布函数比密度更基本}
随机变量的分布函数（cumulative distribution function）为 $F(x)=P(X\le x)$，对所有实值随机变量都存在。它单调不减、右连续，在正负无穷处的极限分别为一和零。单点概率等于跳跃量 $P(X=x)=F(x)-F(x-)$。若分布绝对连续，存在非负可积函数 $f$ 使 $F(x)=\int_{-\infty}^xf(t)dt$；此时单点概率均为零，$f$ 只是几乎处处定义的密度。

验证候选密度，先查非负性和支撑，再查积分为一。积分上下限由支撑决定，而不是看到二元密度便自动写成正方形。对 $0<x<y<1$ 的三角区域，固定 $y$ 时 $x$ 从零到 $y$；固定 $x$ 时 $y$ 从 $x$ 到一。换积分次序之前画出不等式对应的区域，能避免把条件约束漏掉。

若求 $Y=g(X)$ 的分布，最稳妥的方法是先写
\[
F_Y(y)=P(g(X)\le y).
\]
例如平方变换通常把事件变成 $-\sqrt y\le X\le\sqrt y$，自然揭示两条分支。严格单调可微变换可直接使用密度公式，但需要逆函数和导数绝对值。分段单调变换必须对所有逆像求和：
\[
f_Y(y)=\sum_{x:g(x)=y}\frac{f_X(x)}{|g'(x)|},
\]
在常规可微、导数非零的分支上成立，并按几乎处处意义理解。若变换把一段区间压成一个点，结果可能出现原子，此时单纯求导可能漏掉点质量。

独立变量的最大值和最小值常用分布函数最方便：最大值不超过 $t$ 等价于所有变量不超过 $t$；最小值大于 $t$ 等价于所有变量大于 $t$。这些事件关系对离散分布同样适用，但求单点质量时要使用分布函数差。对应训练为习题~\ref{q:Y3.20}、\ref{q:Y3.25}、\ref{q:B8.7}。

\section*{方法四：联合分布、支撑与独立性}
\addcontentsline{toc}{section}{方法四：联合分布、支撑与独立性}
联合密度描述的是随机向量在区域内的概率。求边缘分布应对另一个坐标积分，且积分区间可能依赖保留的坐标。求和变换 $Z=X+Y$ 时，若引入 $W=X$，原支撑 $x,y>0$ 变成 $z>0,0<w<z$。联合密度里即使不再显含 $w$，求边缘仍要积分区间长度，不能直接把它当成 $Z$ 的密度。

独立性（independence）的定义是联合概率对可测集合分解。密度存在时，等价于 $f_{X,Y}(x,y)=f_X(x)f_Y(y)$ 几乎处处成立。这里包括支撑的指示函数。内部表达式能够分离，但支撑不是合适的直积时，不能草率认定独立。相反，在整个直积支撑上分解为两个已经归一化的边缘密度，可以直接推出独立。

二维变换 $u=g(x,y),v=h(x,y)$ 的常规公式要求在所考虑区域内一一可逆、可微且雅可比非零。若逆变换为 $(x,y)=T^{-1}(u,v)$，则
\[
f_{U,V}(u,v)=f_{X,Y}(T^{-1}(u,v))\left|\det D T^{-1}(u,v)\right|.
\]
题目若有多个逆像，需将各分支贡献相加。极坐标中的面积元素是 $r\,dr\,d\theta$，因此从极坐标密度回到平面时除以 $r$。把这个方向搞反常使密度无法归一化。

零协方差只是二阶线性关系的信息。联合正态情形中零相关推出独立；两个边缘分别正态并不足够。判断独立时可用确定性函数关系找反例，例如 $Y^2=X^2$ 会使相应事件重复，破坏概率乘积。对应训练为习题~\ref{q:A4.23}、\ref{q:Y4.25}、\ref{q:B17.15}、\ref{q:B17.16}。

\section*{方法五：先用结构求期望，再求分布}
\addcontentsline{toc}{section}{方法五：先用结构求期望，再求分布}
期望的线性性不需要独立。若目标量能写成指示变量之和，通常可以不求整个分布就得到均值。例如成功数量 $T=\sum_iI_i$ 的期望为 $\sum_iP(I_i=1)$；即使事件相关，这个公式仍成立。求方差时才出现交叉协方差：
\[
\operatorname{Var}\left(\sum_iX_i\right)=\sum_i\operatorname{Var}(X_i)+2\sum_{i<j}\operatorname{Cov}(X_i,X_j).
\]
独立且二阶矩有限可使协方差为零，反过来零协方差一般不意味着独立。

非负变量可使用尾积分公式 $EX=\int_0^\infty P(X>t)dt$；非负整数变量对应尾和 $EX=\sum_{k\ge1}P(X\ge k)$。证明来自将变量写成指示函数积分并使用非负积分交换定理，因此即使期望为无穷也成立。这个方法特别适合最大值、最小值、停时和剩余寿命。

条件期望（conditional expectation）允许按一个容易控制的状态分层：先固定随机样本数、类别或强度，计算条件均值，再取外层期望。全方差公式为
\[
\operatorname{Var}(Y)=E[\operatorname{Var}(Y\mid X)]+\operatorname{Var}(E[Y\mid X]).
\]
前项是固定状态后的内部波动，后项是状态变化造成的均值波动。随机和中常见错误是只保留前项，漏掉随机个数本身导致的方差。

计算结果应保留量纲：均值的单位与变量相同，方差的单位是变量单位平方，标准差恢复原单位，相关系数没有单位。对协方差或方差出现负数，应检查减法和混合模型分解；但协方差本身可以为负。对应训练为习题~\ref{q:A5.9}、\ref{q:A5.25}、\ref{q:A6.4}、\ref{q:B17.5}。

\section*{方法六：生成函数与极限定理}
\addcontentsline{toc}{section}{方法六：生成函数与极限定理}
对非负整数变量，概率生成函数（probability generating function）为 $G_X(s)=Es^X$。在 $|s|\le1$ 范围通常有良好定义，系数给出概率。独立和的生成函数相乘，随机和则产生复合函数。在相应矩有限或以单侧极限解释时，$G'_X(1)=EX$，$G''_X(1)=E[X(X-1)]$，注意二阶导数不是直接给出 $EX^2$。

特征函数（characteristic function）为 $\varphi_X(t)=Ee^{itX}$，对任意实值随机变量都存在。独立性给出 $\varphi_{X+Y}=\varphi_X\varphi_Y$，仿射变换给出 $\varphi_{aX+b}(t)=e^{itb}\varphi_X(at)$。莱维连续性定理将特征函数的逐点极限与依分布收敛联系起来；使用时应确认极限在零点连续并确实对应一个概率分布。

大数定律（law of large numbers）关注平均值稳定到总体均值；中心极限定理（central limit theorem）关注偏差乘以 $\sqrt n$ 后的波动形状。有限正方差的独立同分布版本是
\[
\sqrt n\frac{\overline X-\mu}{\sigma}\xrightarrow dN(0,1).
\]
“样本均值近似正态”并不等于每个原始观测近似正态，也不等于未乘 $\sqrt n$ 的误差服从标准正态。两组样本均值之差的方差要分别除以各自容量。

依分布收敛的分布函数判据通常只对固定连续点成立。若阈值随 $n$ 移动，应使用夹逼或紧性论证。加入一个依概率趋于零的扰动不改变分布极限，这是斯卢茨基思想的常见形式，不要求扰动与原量独立。对应训练为习题~\ref{q:A6.29}、\ref{q:A6.30}、\ref{q:A7.21}、\ref{q:A7.25}。

\section*{方法七：抽样分布是统计推断的桥梁}
\addcontentsline{toc}{section}{方法七：抽样分布是统计推断的桥梁}
统计量（statistic）是计算公式不含未知参数的样本函数；它的分布可以依赖未知参数。抽样分布是把样本看成随机变量后，统计量在重复抽样中的分布。观测到具体数据后，统计量变成数值，但用于评价误差的参考分布来自观测前的模型。

正态样本中三条结构最重要：样本均值为正态；中心化残差平方和为卡方；均值与残差向量独立。令 $n\ge2$，则
\[
\overline X\sim N(\mu,\sigma^2/n),\qquad\frac{(n-1)S^2}{\sigma^2}\sim\chi^2_{n-1}.
\]
两者独立使学生化比值服从 $t_{n-1}$。两组独立正态样本的标准化方差之比服从 $F$ 分布。这里“独立”和“正态”各有作用，不能只因为公式里出现样本方差就自动使用 $t$ 或 $F$。

自由度是可独立变化方向的数目。用样本均值估计总体均值后，残差满足和为零的线性约束，故平方和只剩 $n-1$ 个正交正态方向。在回归中估计 $p$ 个满秩线性系数，残差自由度是 $n-p$；在分类计数中，总数和参数拟合也会消耗自由度，但适用条件需要另外检查。

顺序统计量从另一条途径提供抽样分布。最大值的分布来自“所有观测都不超过阈值”，最小值的分布来自“所有观测都超过阈值”。端点参数影响支撑时，极值往往比均值包含更多相关信息。对应训练为习题~\ref{q:B8.3}、\ref{q:B8.4}、\ref{q:B8.7}、\ref{q:B8.8}。

\section*{方法八：区分估计方法与估计性质}
\addcontentsline{toc}{section}{方法八：区分估计方法与估计性质}
矩估计（method of moments）通过总体矩等于样本矩构造参数估计。极大似然估计（maximum likelihood estimation）固定观察数据，在参数空间中最大化联合密度或质量函数。这两者是构造方法；无偏、一致、效率和均方误差则是评价性质。一个矩估计可以有偏，一个极大似然估计也可能没有有限期望，但仍在其他意义下有用。

均方误差分解为方差与偏差平方：
\[
E_\theta(T-\theta)^2=\operatorname{Var}_\theta(T)+[E_\theta T-\theta]^2.
\]
比较无偏估计时可以只比较方差；比较有偏与无偏估计时必须同时考虑偏差。参数空间受约束时，把估计投影回允许集合可能引入偏差，却减小每个允许真值下的平方误差。

证明一致性常先对样本矩用大数定律，再用连续映射定理。对比值要检查分母极限非零；对数与倒数要处理有限样本中未定义的事件。渐近正态性常由中心极限定理加德尔塔方法得到：若 $\sqrt n(T_n-\theta)\xrightarrow dN(0,v)$ 且 $g$ 在 $\theta$ 处可微，则
\[
\sqrt n[g(T_n)-g(\theta)]\xrightarrow dN(0,[g'(\theta)]^2v).
\]
若一阶导数为零，不能套用这个式子声称得到非退化极限，应继续分析更高阶项。

极大似然解题应写清参数空间和支撑指标，随后才求导。驻点还需判断是最大值，并检查边界和极大值是否达到。均匀分布端点估计正是最大值位于约束边界的典型例子。对应训练为习题~\ref{q:B9.5}、\ref{q:B10.2}、\ref{q:B10.4}、\ref{q:B17.22}。

\section*{方法九：充分性、完全性与信息下界}
\addcontentsline{toc}{section}{方法九：充分性、完全性与信息下界}
充分统计量（sufficient statistic）压缩了模型中关于参数的全部样本信息。因子分解定理要求联合密度写成参数只通过统计量进入的部分与不含参数部分之积。检查支撑时也应把指示函数包括进去。伯努利样本的成功总数、泊松样本的总计数、正态样本的两个平方和型统计量都是常见例子。

拉奥—布莱克化以 $E(U\mid T)$ 替代原无偏估计 $U$，保留期望并使方差不增大。条件期望首先需要作为样本函数可计算；充分性使有关条件分布不依赖未知参数。如果统计量进一步完全，即其函数在所有参数下期望为零只能几乎处处为零，则莱曼—谢费定理给出唯一的均匀最小方差无偏估计。

费希尔信息量（Fisher information）衡量似然对参数的局部敏感度。在正规可微模型中，得分期望为零，信息可写成得分平方的期望，也可写成负的对数似然二阶导数期望。克拉美—拉奥不等式给出无偏估计 $g(\theta)$ 的方差下界
\[
\operatorname{Var}_\theta(U)\ge\frac{[g'(\theta)]^2}{I_n(\theta)}.
\]
使用它需核对支撑不随参数产生未处理的边界项、微分积分交换等条件。它只是一个下界，不保证存在达到下界的估计。

充分性、完全性与信息量回答不同问题：前者关乎数据压缩，中者关乎无偏估计唯一性，后者关乎局部精度限制。它们之间不能仅凭名称相似互相替代。对应训练为习题~\ref{q:B11.3}、\ref{q:B11.4}、\ref{q:B11.6}、\ref{q:B11.8}。

\section*{方法十：区间估计从枢轴反演}
\addcontentsline{toc}{section}{方法十：区间估计从枢轴反演}
枢轴量（pivotal quantity）是样本与未知参数的函数，但其分布不含未知参数。构造区间的一般路线是先找到枢轴，再在其已知分布上选择概率为 $1-\alpha$ 的区域，最后解不等式把参数放在中间。例如正态均值的标准化量、正态方差的卡方量、指数样本和的伽马量都可以这样处理。

反演时注意单调方向。若 $a\le C/\theta\le b$ 且所有量正，则 $C/b\le\theta\le C/a$，分位数顺序应交换。很多方差区间错误正是把较小的卡方分位数放在下端点分母，从而得到端点颠倒。比值参数还需检查分子和分母的自由度。

置信水平（confidence level）是重复抽样中随机区间覆盖固定参数的概率。它不自动等于给定数据以后参数位于已实现区间的后验概率。贝叶斯可信区间（credible interval）则基于后验分布积分。两者数值上有时相同，但逻辑解释不同；一个区间可以有 $95\%$ 的频率覆盖，而在某个特定先验与数据下的后验概率远小于 $95\%$。

样本量规划通常解误差半径不等式。若标准差已知，半径为 $z\sigma/\sqrt n$，所以将半径减半需四倍样本量。未知比例可以用方差上界 $1/4$ 规划，但这仍通常基于正态近似，不能直接声称对所有有限样本精确覆盖。对应训练为习题~\ref{q:M-C22-1}、\ref{q:B13.1}、\ref{q:M-P10-5}、\ref{q:B17.11}。

\section*{方法十一：检验的四个概率方向}
\addcontentsline{toc}{section}{方法十一：检验的四个概率方向}
假设检验（hypothesis testing）由零假设、备择、统计量和拒绝域构成。第一类错误是零假设真实时拒绝，第二类错误是某个备择真实时不拒绝。功效（power）等于备择下拒绝概率，是参数的函数而非通常意义上的一个固定数。对于复合零假设，检验水平要求其中所有参数的误拒概率都不超过 $\alpha$。

$p$ 值是在零假设参考分布下，出现至少同样极端结果的概率。它不是零假设为真的概率，也不是给定这一次拒绝后判断错误的概率。后一类逆向概率需要零假设与备择的先验权重和相应似然。显著性阈值应预先规定；看到数据后更改单侧、双侧方向，会改变原定错误控制。

构造拒绝域后，再在备择分布下积分得到功效。不能把零假设临界值重新以备择均值为中心设定，否则就改变了检验规则。中央区间通常是双侧检验的接受区域，其补集才是拒绝区域。离散检验中拒绝域大小常无法恰好等于 $\alpha$，除非允许边界随机化。

不拒绝零假设不等于证明其为真。若要证明实际差异足够小，应预先规定等效界并使用等效检验；若要比较多个假设，还要控制多重检验产生的至少一次误报概率。置信区间能同时显示差异大小与不确定性，通常比仅报“显著／不显著”更有解释力。对应训练为习题~\ref{q:M-P8-3}、\ref{q:M-P9-1}、\ref{q:B14.3}、\ref{q:B14.4}。

\section*{方法十二：分类计数检验的建模检查}
\addcontentsline{toc}{section}{方法十二：分类计数检验的建模检查}
皮尔逊拟合优度统计量为 $Q=\sum_j(O_j-E_j)^2/E_j$。它衡量观察计数与零假设期望计数的偏离。多项独立抽样、正类别概率和足够大的样本保证常用卡方近似有依据。若零假设完全指定 $k$ 个概率，自由度为 $k-1$；若从同一分组多项似然适当估计 $r$ 个参数，常规正则条件下自由度为 $k-1-r$。

“每个期望计数至少五”是常见经验检查，并不是覆盖所有情形的精确数学边界。极稀疏类别可事先合理合并，也可用精确条件法或模拟零分布。合并类别会改变自由度，且不应在看见哪个类别偏离最大后任意调整规则并仍沿用未经调整的显著性水平。

列联表独立性检验中，期望计数为行和乘列和除以总数。对于 $r\times c$ 表，自由度通常为 $(r-1)(c-1)$。两组类别分布同质性与联合表独立性在数值形式上相同，但抽样设计和零假设叙述不同。固定边际和的二乘二表，可用超几何分布作费希尔精确检验。

实际解释应区分统计偏离与产生偏离的原因。拟合优度拒绝只能说明指定模型与数据不相容，不能单凭它断言欺骗、因果效应或某种机制成立。包装内相关、重复测量和选择性收集都会使简单独立计数模型失效。对应训练为习题~\ref{q:M-E2-II5}、\ref{q:B15.2}、\ref{q:B15.4}、\ref{q:M-F-8}。

\section*{方法十三：回归中的线性与随机性}
\addcontentsline{toc}{section}{方法十三：回归中的线性与随机性}
线性回归（linear regression）中的“线性”指对未知系数线性。模型可以包含 $x^2,x^3,\log x,|x-c|$ 等已知基函数。最小二乘（least squares）本身是一个确定性的最优化问题，不需概率模型就能求出拟合系数；要解释其无偏性、方差或检验分布，才需要对误差建立随机模型。

满列秩设计矩阵保证最小二乘解唯一。正规方程 $X^\mathsf T(Y-X\widehat\beta)=0$ 表明残差与设计空间正交。包含截距时，残差和为零，拟合线经过均值点。高斯—马尔可夫定理只需零均值、同方差和不相关误差，给出在线性无偏类中的最优性；精确 $t$、$F$ 分布则进一步使用正态误差。

对均值响应的区间与新个体预测区间不同。前者衡量回归曲线估计误差，后者还包含新观测的不可约噪声。两者的中心可相同，半径却差一个额外的误差方差项。观测横坐标越远离样本中心，预测标准误通常越大；在样本范围外还存在模型形式不确定性，这不是单纯增加公式中的标准误就能完全解决的。

变换响应会改变拟合目标。最小化对数残差平方与最小化原尺度残差平方是不同的问题；对数正态模型中把拟合的对数均值直接取指数得到条件中位数，均值还需相应校正。回归相关不能自动解释为因果，遗漏变量可能改变斜率。对应训练为习题~\ref{q:M-C26-1c}、\ref{q:B16.2}、\ref{q:B16.4}、\ref{q:B17.23}。

\section*{方法十四：重抽样与模拟误差}
\addcontentsline{toc}{section}{方法十四：重抽样与模拟误差}
非参数自助法（nonparametric bootstrap）从经验分布有放回重抽，参数自助法（parametric bootstrap）从拟合的参数模型重抽。每次重抽样容量通常与原样本相同，从而模拟同一统计量在同样样本规模下的波动。增加重抽样次数 $B$ 降低的是模拟误差，不会把原始的 $n$ 个观测变成 $nB$ 个独立总体观测。

百分位区间直接取自助统计量的两个经验分位数；基本区间先估计统计量相对原估计的误差分布，再作反演。学生化、偏差校正等方法又有不同构造。报告区间时应说明方法、重抽样方式、重抽样次数和分位数约定。对相关数据，普通逐点重抽会破坏依赖结构，需要与抽样设计相适应的方法。

自助法并非“完全无假设”。它依赖经验分布能在所关心的统计量层面近似真实抽样分布。边界参数、不光滑统计量、极值、重尾或小样本，都可能使普通方法失效。尤其参数自助法继承拟合模型的结构；若模型本身错误，模拟得再多也只是精确地模拟了错误模型。

蒙特卡洛积分把积分写成期望，再用样本平均估计。独立样本通常使标准误按平方根速率下降。方差缩减可利用负相关配对、控制变量或重要性抽样，但必须重新核对无偏性和方差存在性。不能仅凭模拟曲线“看起来符合”就把随机证据当作严格证明。对应训练为习题~\ref{q:B13.6}、\ref{q:M-PF-29}、\ref{q:B17.17}、\ref{q:B17.18}。

\section*{方法十五：解答的自检与复习路线}
\addcontentsline{toc}{section}{方法十五：解答的自检与复习路线}
概率题完成后，先核查事件或积分区域与文字是否一致，再检查归一化、支持范围与边界。当密度不积分为一时，不应先把常数“修到正确”，而应回查是不是漏掉雅可比、混淆圆周与圆内均匀、或者在联合变换后漏了边缘积分。条件概率若大于一，多半是分母或条件事件写错；所谓独立性若由一句“显然”带过，应回到定义或联合密度分解核实。

统计题完成后，核查总体模型、抽样独立性、已知与未知参数、统计量分布、自由度和尾方向。标准差与方差单位不同，标准误还要考虑样本容量。极大似然估计的数值并不能直接代替未知参数出现在精确枢轴中；一旦这样替代，常常需要说明近似依据。所有小数先由精确式计算，在最终结果处统一舍入，避免多步后验更新中的舍入传播。

建议第一轮先做事件、条件概率、分布与期望的基本计算；第二轮做联合变换、极限定理和抽样分布；第三轮做估计与检验，并把每个区间都追溯到所用枢轴；最后做综合题，练习从自然语言挑选合适模型。证明题与计算题交替进行，能避免只会套用符号而不知道条件何时失效。

错误记录宜包括“原判断、出错原因、修正后的条件或公式、一个反例”。例如把零相关当作独立，应补记联合正态的额外条件和 $Y=SX$ 的反例；把不拒绝当作证明相等，应补记功效和区间宽度；把均匀分布的极大似然用内点求导处理，应补记参数相关支撑。这样形成的复习记录比只抄最终答案更能迁移到新题。


\clearpage
\section*{常用分布：定义、矩与适用情形}
\addcontentsline{toc}{section}{常用分布：定义、矩与适用情形}
\subsection*{伯努利分布与二项分布}
伯努利变量 $X\sim\operatorname{Bern}(p)$ 取零或一，$P(X=1)=p$。它表示一次二分类试验，不要求实际结果的文字名称为“成功”或“失败”。
\[
EX=p,\quad\operatorname{Var}(X)=p(1-p),\quad G_X(s)=1-p+ps.
\]
若 $n$ 次试验独立且成功概率相同，则成功数 $S$ 服从二项分布：
\[
P(S=k)=\binom nkp^k(1-p)^{n-k},\quad k=0,\ldots,n,
\]
\[
ES=np,\quad\operatorname{Var}(S)=np(1-p),\quad G_S(s)=(1-p+ps)^n.
\]
二项模型有三个独立的检查点：次数固定、试验独立、单次概率相同。不同成功概率的独立和仍可用生成函数乘积，但一般不是单参数二项分布；无放回抽样通常具有负相关，常应使用超几何分布。
\subsection*{超几何分布}
总体规模 $N$，其中 $M$ 个特征对象，无放回抽取 $n$ 个，特征数量 $X$ 的质量函数为
\[
P(X=k)=\frac{\binom Mk\binom{N-M}{n-k}}{\binom Nn},
\]
\[
\max(0,n+M-N)\le k\le\min(n,M).
\]
令 $p=M/N$，当 $N>1$ 时，
\[
EX=np,\qquad\operatorname{Var}(X)=np(1-p)\frac{N-n}{N-1}.
\]
支撑范围来自不能抽到超过总体已有数量的对象。若 $n=N$，方差为零，体现普查没有抽样不确定性。若总体本身的数值被当作随机生成，另一个层次的总体随机性不包括在这个条件抽样方差中。
\subsection*{几何分布与负二项分布}
本文几何分布的基本约定是记录首次成功的试验编号：
\[
P(T=k)=p(1-p)^{k-1},\quad k\ge1,
\]
\[
ET=1/p,\qquad\operatorname{Var}(T)=(1-p)/p^2.
\]
若记录成功前的失败次数 $F=T-1$，则均值变为 $(1-p)/p$，方差不变。记录第 $r$ 次成功前的失败次数 $K$ 时，
\[
P(K=k)=\binom{k+r-1}{r-1}p^r(1-p)^k,\quad k\ge0,
\]
\[
EK=\frac{r(1-p)}p,\qquad\operatorname{Var}(K)=\frac{r(1-p)}{p^2}.
\]
记录总次数时应再加 $r$。分布名称不足以唯一确定参数约定；解题时应先写出质量函数或明确计数对象。
\subsection*{泊松分布}
\[
P(X=k)=e^{-\lambda}\frac{\lambda^k}{k!},\quad k=0,1,\ldots,
\]
\[
EX=\operatorname{Var}(X)=\lambda,\qquad G_X(s)=e^{\lambda(s-1)}.
\]
独立泊松变量之和仍为泊松，参数相加。泊松总数经过相互独立的分类后，各类别计数仍是相互独立泊松变量。这两条性质分别对应事件来源叠加和事件流分裂。均值等于方差是泊松模型的性质，不是所有计数变量的性质；混合强度通常使方差大于均值。

\subsection*{均匀分布与指数分布}
$X\sim U(a,b)$，$a<b$ 时，
\[
f(x)=\frac{\mathbf1_{(a,b)}(x)}{b-a},\qquad EX=\frac{a+b}2,
\]
\[
\operatorname{Var}(X)=\frac{(b-a)^2}{12}.
\]
端点是否包含不影响连续分布，但在参数相关支撑的极大似然问题中，应明确采用的密度版本以讨论最大值是否达到。

$T\sim\operatorname{Exp}(\lambda)$，$\lambda>0$ 时，
\[
f(t)=\lambda e^{-\lambda t}\mathbf1_{(0,\infty)}(t),\quad P(T>t)=e^{-\lambda t}\quad(t\ge0),
\]
\[
ET=1/\lambda,\quad\operatorname{Var}(T)=1/\lambda^2.
\]
指数分布满足无记忆性 $P(T>s+t\mid T>s)=P(T>t)$，这里 $s,t\ge0$。恒定风险率模型可导出指数寿命，但不能从“等待时间”这个词自动推断指数分布。

\subsection*{伽马分布与贝塔分布}
本文伽马采用形状 $a>0$、率 $b>0$：
\[
f(x)=\frac{b^a}{\Gamma(a)}x^{a-1}e^{-bx}\mathbf1_{(0,\infty)}(x),
\]
\[
EX=a/b,\qquad\operatorname{Var}(X)=a/b^2.
\]
相同率的独立伽马变量相加，形状相加；不同率时一般不保持伽马形式。$n$ 个同率指数变量之和是形状 $n$ 的伽马变量。

贝塔分布采用两个正形状参数：
\[
f(x)=\frac{x^{a-1}(1-x)^{b-1}}{B(a,b)},\quad0<x<1,
\]
\[
B(a,b)=\frac{\Gamma(a)\Gamma(b)}{\Gamma(a+b)},\quad EX=\frac a{a+b},
\]
\[
\operatorname{Var}(X)=\frac{ab}{(a+b)^2(a+b+1)}.
\]
它既用于比例的先验，也出现在均匀顺序统计量和同率伽马变量份额的分布中。分布参数为 $a,b$ 时，密度幂次为 $a-1,b-1$，更新共轭参数时要避免差一。

\subsection*{正态分布与三种抽样分布}
$X\sim N(\mu,\sigma^2)$ 时，标准化 $Z=(X-\mu)/\sigma\sim N(0,1)$。正态密度为
\[
f(x)=\frac1{\sqrt{2\pi}\sigma}\exp\left[-\frac{(x-\mu)^2}{2\sigma^2}\right].
\]
独立正态变量线性组合仍为正态；联合正态向量的任意线性组合也为正态。仅有各边缘正态不够。

若 $Z_1,\ldots,Z_\nu$ 独立标准正态，则 $V=\sum_iZ_i^2\sim\chi^2_\nu$，其均值为 $\nu$，方差为 $2\nu$。若 $Z\sim N(0,1)$ 与 $V\sim\chi^2_\nu$ 独立，则 $T=Z/\sqrt{V/\nu}\sim t_\nu$。若两个独立卡方变量自由度为 $a,b$，则标准化比值 $(V_1/a)/(V_2/b)\sim F_{a,b}$。比值中每个卡方变量都必须先除以自己的自由度。

\clearpage
\section*{估计与检验公式的条件清单}
\addcontentsline{toc}{section}{估计与检验公式的条件清单}
\subsection*{一个正态总体的均值}
方差已知时使用 $Z=\sqrt n(\overline X-\mu_0)/\sigma$。双侧水平 $\alpha$ 拒绝域为 $|Z|>z_{1-\alpha/2}$，相应均值区间为
\[
\overline X\pm z_{1-\alpha/2}\frac\sigma{\sqrt n}.
\]
方差未知时，正态总体下用 $T=\sqrt n(\overline X-\mu_0)/S\sim t_{n-1}$，相应区间将 $z,\sigma$ 替为 $t_{n-1;1-\alpha/2},S$。在非正态有限方差总体中，学生化均值通常只是渐近标准正态，不能直接声称精确 $t$ 分布。
\subsection*{一个正态总体的方差}
未知均值时用 $Q=(n-1)S^2/\sigma_0^2\sim\chi^2_{n-1}$。右侧方差检验拒绝大 $Q$，左侧检验拒绝小 $Q$。等尾双侧区间为
\[
\left[\frac{(n-1)S^2}{\chi^2_{n-1;1-\alpha/2}},\frac{(n-1)S^2}{\chi^2_{n-1;\alpha/2}}\right].
\]
若均值为已知常数，改用相对该常数的平方和，自由度为 $n$。不能先用样本均值中心化又声称保留了 $n$ 个自由度。
\subsection*{两组独立正态均值}
共同方差未知时使用合并方差
\[
S_p^2=\frac{(n-1)S_1^2+(m-1)S_2^2}{n+m-2}.
\]
均值差的标准误为 $S_p\sqrt{1/n+1/m}$，自由度为 $n+m-2$。若不假设共同方差，使用两组各自估计方差构成标准误，并用韦尔奇近似自由度。配对数据先构造差值样本，再作单样本分析；两独立样本公式不适用于任意前后测量。
\subsection*{贝塔—二项与伽马—泊松更新}
$\theta\sim\operatorname{Beta}(a,b)$，观测 $n$ 次伯努利中的 $s$ 次成功后，后验为 $\operatorname{Beta}(a+s,b+n-s)$，下一次成功概率为 $(a+s)/(a+b+n)$。

$\lambda\sim\operatorname{Gamma}(a,b)$，观测独立的 $X_i\sim\operatorname{Pois}(t_i\lambda)$ 后，后验为
\[
\operatorname{Gamma}\left(a+\sum_iX_i,\ b+\sum_it_i\right).
\]
指数寿命 $X_i\sim\operatorname{Exp}(\lambda)$ 配伽马先验时，形状增加样本容量，率增加总寿命：$\operatorname{Gamma}(a+n,b+\sum_iX_i)$。泊松计数与指数等待对参数的作用不同，不要把两条更新公式交换。
\subsection*{一元回归的常用量}
令 $S_{xx}=\sum_i(x_i-\overline x)^2$、$S_{xy}=\sum_i(x_i-\overline x)(y_i-\overline y)$。含截距且 $S_{xx}>0$ 时，
\[
\widehat a=S_{xy}/S_{xx},\qquad\widehat b=\overline y-\widehat a\overline x.
\]
独立零均值同方差误差下，斜率方差为 $\sigma^2/S_{xx}$，截距方差为 $\sigma^2(1/n+\overline x^2/S_{xx})$。正态误差且 $n>2$ 时，残差方差估计为 $s^2=SSE/(n-2)$，再用 $t_{n-2}$ 分位数构造精确区间。误差异方差或相关时，上述方差表达式需要修改。

\clearpage
\section*{常用分位数表与计算约定}
\addcontentsline{toc}{section}{常用分位数表与计算约定}
以下所有分位数都采用下侧概率约定 $P(X\le q_p)=p$，保留四位小数。表格用于手工练习，正式计算可直接按分布函数求逆。表中舍入可能使临界附近的结果略有差异，应以未舍入表达式为准。


\subsection*{标准正态分位数}

\begin{center}\begin{tabular}{cc}\toprule 下侧概率 $p$ & $z_p$\\\midrule

0.800 & 0.8416\\

0.850 & 1.0364\\

0.900 & 1.2816\\

0.950 & 1.6449\\

0.960 & 1.7507\\

0.975 & 1.9600\\

0.990 & 2.3263\\

0.995 & 2.5758\\

0.999 & 3.0902\\

\bottomrule\end{tabular}\end{center}

\subsection*{学生 t 分位数}

\begin{center}\begin{tabular}{rrrr}\toprule $\nu$ & $t_{\nu;0.95}$ & $t_{\nu;0.975}$ & $t_{\nu;0.99}$\\\midrule

1 & 6.3138 & 12.7062 & 31.8205\\

2 & 2.9200 & 4.3027 & 6.9646\\

3 & 2.3534 & 3.1824 & 4.5407\\

4 & 2.1318 & 2.7764 & 3.7469\\

5 & 2.0150 & 2.5706 & 3.3649\\

6 & 1.9432 & 2.4469 & 3.1427\\

7 & 1.8946 & 2.3646 & 2.9980\\

8 & 1.8595 & 2.3060 & 2.8965\\

9 & 1.8331 & 2.2622 & 2.8214\\

10 & 1.8125 & 2.2281 & 2.7638\\

12 & 1.7823 & 2.1788 & 2.6810\\

15 & 1.7531 & 2.1314 & 2.6025\\

19 & 1.7291 & 2.0930 & 2.5395\\

20 & 1.7247 & 2.0860 & 2.5280\\

25 & 1.7081 & 2.0595 & 2.4851\\

30 & 1.6973 & 2.0423 & 2.4573\\

40 & 1.6839 & 2.0211 & 2.4233\\

60 & 1.6706 & 2.0003 & 2.3901\\

100 & 1.6602 & 1.9840 & 2.3642\\

\bottomrule\end{tabular}\end{center}

\subsection*{卡方下侧分位数}

\begin{center}\begin{tabular}{rrrrr}\toprule $\nu$ & $0.025$ & $0.05$ & $0.95$ & $0.975$\\\midrule

1 & 0.0010 & 0.0039 & 3.8415 & 5.0239\\

2 & 0.0506 & 0.1026 & 5.9915 & 7.3778\\

3 & 0.2158 & 0.3518 & 7.8147 & 9.3484\\

4 & 0.4844 & 0.7107 & 9.4877 & 11.1433\\

5 & 0.8312 & 1.1455 & 11.0705 & 12.8325\\

6 & 1.2373 & 1.6354 & 12.5916 & 14.4494\\

7 & 1.6899 & 2.1673 & 14.0671 & 16.0128\\

8 & 2.1797 & 2.7326 & 15.5073 & 17.5345\\

9 & 2.7004 & 3.3251 & 16.9190 & 19.0228\\

10 & 3.2470 & 3.9403 & 18.3070 & 20.4832\\

11 & 3.8157 & 4.5748 & 19.6751 & 21.9200\\

12 & 4.4038 & 5.2260 & 21.0261 & 23.3367\\

15 & 6.2621 & 7.2609 & 24.9958 & 27.4884\\

19 & 8.9065 & 10.1170 & 30.1435 & 32.8523\\

20 & 9.5908 & 10.8508 & 31.4104 & 34.1696\\

30 & 16.7908 & 18.4927 & 43.7730 & 46.9792\\

40 & 24.4330 & 26.5093 & 55.7585 & 59.3417\\

60 & 40.4817 & 43.1880 & 79.0819 & 83.2977\\

\bottomrule\end{tabular}\end{center}

\clearpage\subsection*{F 分布的 $0.95$ 下侧分位数}

行表示分子自由度，列表示分母自由度。

\begin{center}\begin{tabular}{rrrrr}\toprule $a\backslash b$ & $5$ & $10$ & $20$ & $60$\\\midrule

1 & 6.6079 & 4.9646 & 4.3512 & 4.0012\\

2 & 5.7861 & 4.1028 & 3.4928 & 3.1504\\

3 & 5.4095 & 3.7083 & 3.0984 & 2.7581\\

4 & 5.1922 & 3.4780 & 2.8661 & 2.5252\\

5 & 5.0503 & 3.3258 & 2.7109 & 2.3683\\

10 & 4.7351 & 2.9782 & 2.3479 & 1.9926\\

20 & 4.5581 & 2.7740 & 2.1242 & 1.7480\\

\bottomrule\end{tabular}\end{center}

\clearpage\section*{八组图解：把公式与图形对应起来}\addcontentsline{toc}{section}{八组图解：把公式与图形对应起来}

\subsection*{图解一：圆周与圆内的均匀取点}

\begin{center}% 使用共享样式已加载的 TikZ，保留原图数据。
\begin{tikzpicture}[x=.87\linewidth,y=.69\linewidth]
% 直接设置数据坐标单位，避免先形成大尺寸再缩放造成溢出。
\begin{scope}[x=0.435\linewidth,y=0.383333333364\linewidth,shift={(1,0)}]
\clip (-1,0) rectangle (1,1.8);
\draw[black,thick] plot coordinates {(-0.97000,1.309352) (-0.95916,1.125337) (-0.94832,1.003171) (-0.93749,0.914623) (-0.92665,0.846726) (-0.91581,0.792581) (-0.90497,0.748135) (-0.89413,0.710832) (-0.88330,0.678966) (-0.87246,0.651355) (-0.86162,0.627145) (-0.85078,0.605706) (-0.83994,0.586559) (-0.82911,0.569334) (-0.81827,0.553740) (-0.80743,0.539542) (-0.79659,0.526552) (-0.78575,0.514616) (-0.77492,0.503604) (-0.76408,0.493409) (-0.75324,0.483941) (-0.74240,0.475122) (-0.73156,0.466886) (-0.72073,0.459176) (-0.70989,0.451943) (-0.69905,0.445144) (-0.68821,0.438740) (-0.67737,0.432699) (-0.66654,0.426991) (-0.65570,0.421590) (-0.64486,0.416472) (-0.63402,0.411618) (-0.62318,0.407007) (-0.61235,0.402623) (-0.60151,0.398452) (-0.59067,0.394479) (-0.57983,0.390691) (-0.56899,0.387078) (-0.55816,0.383628) (-0.54732,0.380333) (-0.53648,0.377183) (-0.52564,0.374172) (-0.51480,0.371290) (-0.50397,0.368532) (-0.49313,0.365892) (-0.48229,0.363363) (-0.47145,0.360940) (-0.46061,0.358619) (-0.44978,0.356394) (-0.43894,0.354261) (-0.42810,0.352217) (-0.41726,0.350258) (-0.40642,0.348381) (-0.39559,0.346581) (-0.38475,0.344856) (-0.37391,0.343204) (-0.36307,0.341622) (-0.35223,0.340107) (-0.34140,0.338657) (-0.33056,0.337269) (-0.31972,0.335943) (-0.30888,0.334675) (-0.29804,0.333465) (-0.28721,0.332311) (-0.27637,0.331210) (-0.26553,0.330162) (-0.25469,0.329165) (-0.24385,0.328218) (-0.23302,0.327320) (-0.22218,0.326470) (-0.21134,0.325666) (-0.20050,0.324908) (-0.18966,0.324194) (-0.17883,0.323525) (-0.16799,0.322899) (-0.15715,0.322315) (-0.14631,0.321773) (-0.13547,0.321272) (-0.12464,0.320811) (-0.11380,0.320391) (-0.10296,0.320011) (-0.09212,0.319669) (-0.08128,0.319367) (-0.07045,0.319103) (-0.05961,0.318877) (-0.04877,0.318689) (-0.03793,0.318539) (-0.02709,0.318427) (-0.01626,0.318352) (-0.00542,0.318315) (0.00542,0.318315) (0.01626,0.318352) (0.02709,0.318427) (0.03793,0.318539) (0.04877,0.318689) (0.05961,0.318877) (0.07045,0.319103) (0.08128,0.319367) (0.09212,0.319669) (0.10296,0.320011) (0.11380,0.320391) (0.12464,0.320811) (0.13547,0.321272) (0.14631,0.321773) (0.15715,0.322315) (0.16799,0.322899) (0.17883,0.323525) (0.18966,0.324194) (0.20050,0.324908) (0.21134,0.325666) (0.22218,0.326470) (0.23302,0.327320) (0.24385,0.328218) (0.25469,0.329165) (0.26553,0.330162) (0.27637,0.331210) (0.28721,0.332311) (0.29804,0.333465) (0.30888,0.334675) (0.31972,0.335943) (0.33056,0.337269) (0.34140,0.338657) (0.35223,0.340107) (0.36307,0.341622) (0.37391,0.343204) (0.38475,0.344856) (0.39559,0.346581) (0.40642,0.348381) (0.41726,0.350258) (0.42810,0.352217) (0.43894,0.354261) (0.44978,0.356394) (0.46061,0.358619) (0.47145,0.360940) (0.48229,0.363363) (0.49313,0.365892) (0.50397,0.368532) (0.51480,0.371290) (0.52564,0.374172) (0.53648,0.377183) (0.54732,0.380333) (0.55816,0.383628) (0.56899,0.387078) (0.57983,0.390691) (0.59067,0.394479) (0.60151,0.398452) (0.61235,0.402623) (0.62318,0.407007) (0.63402,0.411618) (0.64486,0.416472) (0.65570,0.421590) (0.66654,0.426991) (0.67737,0.432699) (0.68821,0.438740) (0.69905,0.445144) (0.70989,0.451943) (0.72073,0.459176) (0.73156,0.466886) (0.74240,0.475122) (0.75324,0.483941) (0.76408,0.493409) (0.77492,0.503604) (0.78575,0.514616) (0.79659,0.526552) (0.80743,0.539542) (0.81827,0.553740) (0.82911,0.569334) (0.83994,0.586559) (0.85078,0.605706) (0.86162,0.627145) (0.87246,0.651355) (0.88330,0.678966) (0.89413,0.710832) (0.90497,0.748135) (0.91581,0.792581) (0.92665,0.846726) (0.93749,0.914623) (0.94832,1.003171) (0.95916,1.125337) (0.97000,1.309352)};
\draw[black,dashed,thick] plot coordinates {(-0.97000,0.154765) (-0.95916,0.180073) (-0.94832,0.202002) (-0.93749,0.221558) (-0.92665,0.239325) (-0.91581,0.255674) (-0.90497,0.270863) (-0.89413,0.285078) (-0.88330,0.298457) (-0.87246,0.311109) (-0.86162,0.323119) (-0.85078,0.334556) (-0.83994,0.345476) (-0.82911,0.355929) (-0.81827,0.365952) (-0.80743,0.375582) (-0.79659,0.384848) (-0.78575,0.393774) (-0.77492,0.402384) (-0.76408,0.410698) (-0.75324,0.418734) (-0.74240,0.426506) (-0.73156,0.434030) (-0.72073,0.441317) (-0.70989,0.448380) (-0.69905,0.455229) (-0.68821,0.461874) (-0.67737,0.468322) (-0.66654,0.474583) (-0.65570,0.480663) (-0.64486,0.486569) (-0.63402,0.492307) (-0.62318,0.497884) (-0.61235,0.503305) (-0.60151,0.508574) (-0.59067,0.513697) (-0.57983,0.518677) (-0.56899,0.523519) (-0.55816,0.528226) (-0.54732,0.532803) (-0.53648,0.537252) (-0.52564,0.541576) (-0.51480,0.545779) (-0.50397,0.549863) (-0.49313,0.553832) (-0.48229,0.557686) (-0.47145,0.561430) (-0.46061,0.565064) (-0.44978,0.568591) (-0.43894,0.572014) (-0.42810,0.575333) (-0.41726,0.578551) (-0.40642,0.581670) (-0.39559,0.584690) (-0.38475,0.587614) (-0.37391,0.590442) (-0.36307,0.593177) (-0.35223,0.595820) (-0.34140,0.598371) (-0.33056,0.600832) (-0.31972,0.603205) (-0.30888,0.605489) (-0.29804,0.607687) (-0.28721,0.609798) (-0.27637,0.611825) (-0.26553,0.613767) (-0.25469,0.615625) (-0.24385,0.617401) (-0.23302,0.619095) (-0.22218,0.620708) (-0.21134,0.622240) (-0.20050,0.623692) (-0.18966,0.625064) (-0.17883,0.626358) (-0.16799,0.627573) (-0.15715,0.628710) (-0.14631,0.629769) (-0.13547,0.630751) (-0.12464,0.631656) (-0.11380,0.632484) (-0.10296,0.633236) (-0.09212,0.633913) (-0.08128,0.634513) (-0.07045,0.635038) (-0.05961,0.635488) (-0.04877,0.635862) (-0.03793,0.636162) (-0.02709,0.636386) (-0.01626,0.636536) (-0.00542,0.636610) (0.00542,0.636610) (0.01626,0.636536) (0.02709,0.636386) (0.03793,0.636162) (0.04877,0.635862) (0.05961,0.635488) (0.07045,0.635038) (0.08128,0.634513) (0.09212,0.633913) (0.10296,0.633236) (0.11380,0.632484) (0.12464,0.631656) (0.13547,0.630751) (0.14631,0.629769) (0.15715,0.628710) (0.16799,0.627573) (0.17883,0.626358) (0.18966,0.625064) (0.20050,0.623692) (0.21134,0.622240) (0.22218,0.620708) (0.23302,0.619095) (0.24385,0.617401) (0.25469,0.615625) (0.26553,0.613767) (0.27637,0.611825) (0.28721,0.609798) (0.29804,0.607687) (0.30888,0.605489) (0.31972,0.603205) (0.33056,0.600832) (0.34140,0.598371) (0.35223,0.595820) (0.36307,0.593177) (0.37391,0.590442) (0.38475,0.587614) (0.39559,0.584690) (0.40642,0.581670) (0.41726,0.578551) (0.42810,0.575333) (0.43894,0.572014) (0.44978,0.568591) (0.46061,0.565064) (0.47145,0.561430) (0.48229,0.557686) (0.49313,0.553832) (0.50397,0.549863) (0.51480,0.545779) (0.52564,0.541576) (0.53648,0.537252) (0.54732,0.532803) (0.55816,0.528226) (0.56899,0.523519) (0.57983,0.518677) (0.59067,0.513697) (0.60151,0.508574) (0.61235,0.503305) (0.62318,0.497884) (0.63402,0.492307) (0.64486,0.486569) (0.65570,0.480663) (0.66654,0.474583) (0.67737,0.468322) (0.68821,0.461874) (0.69905,0.455229) (0.70989,0.448380) (0.72073,0.441317) (0.73156,0.434030) (0.74240,0.426506) (0.75324,0.418734) (0.76408,0.410698) (0.77492,0.402384) (0.78575,0.393774) (0.79659,0.384848) (0.80743,0.375582) (0.81827,0.365952) (0.82911,0.355929) (0.83994,0.345476) (0.85078,0.334556) (0.86162,0.323119) (0.87246,0.311109) (0.88330,0.298457) (0.89413,0.285078) (0.90497,0.270863) (0.91581,0.255674) (0.92665,0.239325) (0.93749,0.221558) (0.94832,0.202002) (0.95916,0.180073) (0.97000,0.154765)};
\end{scope}
\draw[->] (0,0) -- (1.02,0);
\draw[->] (0,0) -- (0,1.02);
\draw (0,0) -- (0,-.012) node[below,font=\small] {$ -1 $};
\draw (0.25,0) -- (0.25,-.012) node[below,font=\small] {$ -0.5 $};
\draw (0.5,0) -- (0.5,-.012) node[below,font=\small] {$ 0 $};
\draw (0.75,0) -- (0.75,-.012) node[below,font=\small] {$ 0.5 $};
\draw (1,0) -- (1,-.012) node[below,font=\small] {$ 1 $};
\draw (0,0) -- (-.012,0) node[left,font=\small] {$ 0 $};
\draw (0,0.2777777778) -- (-.012,0.2777777778) node[left,font=\small] {$ 0.5 $};
\draw (0,0.5555555556) -- (-.012,0.5555555556) node[left,font=\small] {$ 1 $};
\draw (0,0.8333333333) -- (-.012,0.8333333333) node[left,font=\small] {$ 1.5 $};
\node[below,font=\small] at (.5,-.09) {$x$};
\node[rotate=90,font=\small] at (-.13,.5) {横坐标密度};
\node[anchor=north west,fill=white,inner sep=3pt,font=\footnotesize] at (.03,.98) {%
\begin{tabular}{@{}ll@{}}
\tikz[baseline=-.5ex]{\draw[black,thick] (0,0) -- (.5,0);} & 圆周均匀\\
\tikz[baseline=-.5ex]{\draw[black,dashed,thick] (0,0) -- (.5,0);} & 圆内均匀\\
\end{tabular}};
\end{tikzpicture}\end{center}

取单位圆。圆周均匀指角度按弧长均匀；圆内均匀指平面面积均匀。前者横坐标密度在两端增大，后者在中心最大。两者均值都为零，但方差分别为 $1/2$ 和 $1/4$。图中为显示内部形状，仅绘制 $|x|\le0.97$；圆周密度在端点的可积奇性没有被改成有限值。

用截条面积可理解圆内密度：横坐标附近宽 $dx$ 的竖条长度约为 $2\sqrt{1-x^2}$，除以圆面积 $\pi$ 得到密度。圆周投影则需要对角度变换，两条对称弧贡献相加。参见习题~\ref{q:Y3.26}、\ref{q:A5.35}。

\clearpage\subsection*{图解二：二项计数与连续性修正}

\begin{center}% 使用共享样式已加载的 TikZ，保留原图数据。
\begin{tikzpicture}[x=.87\linewidth,y=.69\linewidth]
% 直接设置数据坐标单位，避免先形成大尺寸再缩放造成溢出。
\begin{scope}[x=0.0348\linewidth,y=5.09225092269\linewidth,shift={(-3,0)}]
\clip (3,0) rectangle (28,0.1355);
\filldraw[fill=gray!25,draw=gray] ([xshift=-2pt]3.00000,0) rectangle ([xshift=2pt]3.00000,0.000028);
\filldraw[fill=gray!25,draw=gray] ([xshift=-2pt]4.00000,0) rectangle ([xshift=2pt]4.00000,0.000140);
\filldraw[fill=gray!25,draw=gray] ([xshift=-2pt]5.00000,0) rectangle ([xshift=2pt]5.00000,0.000551);
\filldraw[fill=gray!25,draw=gray] ([xshift=-2pt]6.00000,0) rectangle ([xshift=2pt]6.00000,0.001771);
\filldraw[fill=gray!25,draw=gray] ([xshift=-2pt]7.00000,0) rectangle ([xshift=2pt]7.00000,0.004770);
\filldraw[fill=gray!25,draw=gray] ([xshift=-2pt]8.00000,0) rectangle ([xshift=2pt]8.00000,0.010989);
\filldraw[fill=gray!25,draw=gray] ([xshift=-2pt]9.00000,0) rectangle ([xshift=2pt]9.00000,0.021978);
\filldraw[fill=gray!25,draw=gray] ([xshift=-2pt]10.00000,0) rectangle ([xshift=2pt]10.00000,0.038619);
\filldraw[fill=gray!25,draw=gray] ([xshift=-2pt]11.00000,0) rectangle ([xshift=2pt]11.00000,0.060185);
\filldraw[fill=gray!25,draw=gray] ([xshift=-2pt]12.00000,0) rectangle ([xshift=2pt]12.00000,0.083830);
\filldraw[fill=gray!25,draw=gray] ([xshift=-2pt]13.00000,0) rectangle ([xshift=2pt]13.00000,0.105017);
\filldraw[fill=gray!25,draw=gray] ([xshift=-2pt]14.00000,0) rectangle ([xshift=2pt]14.00000,0.118948);
\filldraw[fill=gray!25,draw=gray] ([xshift=-2pt]15.00000,0) rectangle ([xshift=2pt]15.00000,0.122347);
\filldraw[fill=gray!25,draw=gray] ([xshift=-2pt]16.00000,0) rectangle ([xshift=2pt]16.00000,0.114700);
\filldraw[fill=gray!25,draw=gray] ([xshift=-2pt]17.00000,0) rectangle ([xshift=2pt]17.00000,0.098314);
\filldraw[fill=gray!25,draw=gray] ([xshift=-2pt]18.00000,0) rectangle ([xshift=2pt]18.00000,0.077247);
\filldraw[fill=gray!25,draw=gray] ([xshift=-2pt]19.00000,0) rectangle ([xshift=2pt]19.00000,0.055757);
\filldraw[fill=gray!25,draw=gray] ([xshift=-2pt]20.00000,0) rectangle ([xshift=2pt]20.00000,0.037039);
\filldraw[fill=gray!25,draw=gray] ([xshift=-2pt]21.00000,0) rectangle ([xshift=2pt]21.00000,0.022677);
\filldraw[fill=gray!25,draw=gray] ([xshift=-2pt]22.00000,0) rectangle ([xshift=2pt]22.00000,0.012811);
\filldraw[fill=gray!25,draw=gray] ([xshift=-2pt]23.00000,0) rectangle ([xshift=2pt]23.00000,0.006684);
\filldraw[fill=gray!25,draw=gray] ([xshift=-2pt]24.00000,0) rectangle ([xshift=2pt]24.00000,0.003223);
\filldraw[fill=gray!25,draw=gray] ([xshift=-2pt]25.00000,0) rectangle ([xshift=2pt]25.00000,0.001436);
\filldraw[fill=gray!25,draw=gray] ([xshift=-2pt]26.00000,0) rectangle ([xshift=2pt]26.00000,0.000592);
\filldraw[fill=gray!25,draw=gray] ([xshift=-2pt]27.00000,0) rectangle ([xshift=2pt]27.00000,0.000225);
\filldraw[fill=gray!25,draw=gray] ([xshift=-2pt]28.00000,0) rectangle ([xshift=2pt]28.00000,0.000079);
\draw[black,thick] plot coordinates {(3.00000,0.000130) (3.12563,0.000149) (3.25126,0.000172) (3.37688,0.000198) (3.50251,0.000227) (3.62814,0.000261) (3.75377,0.000298) (3.87940,0.000341) (4.00503,0.000389) (4.13065,0.000444) (4.25628,0.000505) (4.38191,0.000574) (4.50754,0.000651) (4.63317,0.000737) (4.75879,0.000834) (4.88442,0.000942) (5.01005,0.001063) (5.13568,0.001197) (5.26131,0.001346) (5.38693,0.001511) (5.51256,0.001694) (5.63819,0.001896) (5.76382,0.002119) (5.88945,0.002365) (6.01508,0.002635) (6.14070,0.002932) (6.26633,0.003257) (6.39196,0.003613) (6.51759,0.004002) (6.64322,0.004427) (6.76884,0.004888) (6.89447,0.005390) (7.02010,0.005935) (7.14573,0.006524) (7.27136,0.007162) (7.39698,0.007850) (7.52261,0.008591) (7.64824,0.009388) (7.77387,0.010243) (7.89950,0.011160) (8.02513,0.012140) (8.15075,0.013187) (8.27638,0.014302) (8.40201,0.015489) (8.52764,0.016748) (8.65327,0.018083) (8.77889,0.019495) (8.90452,0.020986) (9.03015,0.022557) (9.15578,0.024209) (9.28141,0.025942) (9.40704,0.027759) (9.53266,0.029657) (9.65829,0.031638) (9.78392,0.033701) (9.90955,0.035845) (10.03518,0.038067) (10.16080,0.040366) (10.28643,0.042740) (10.41206,0.045186) (10.53769,0.047700) (10.66332,0.050278) (10.78894,0.052916) (10.91457,0.055608) (11.04020,0.058350) (11.16583,0.061135) (11.29146,0.063957) (11.41709,0.066808) (11.54271,0.069682) (11.66834,0.072571) (11.79397,0.075465) (11.91960,0.078357) (12.04523,0.081238) (12.17085,0.084098) (12.29648,0.086928) (12.42211,0.089718) (12.54774,0.092459) (12.67337,0.095141) (12.79899,0.097753) (12.92462,0.100286) (13.05025,0.102730) (13.17588,0.105076) (13.30151,0.107313) (13.42714,0.109434) (13.55276,0.111429) (13.67839,0.113291) (13.80402,0.115010) (13.92965,0.116580) (14.05528,0.117993) (14.18090,0.119245) (14.30653,0.120329) (14.43216,0.121240) (14.55779,0.121975) (14.68342,0.122530) (14.80905,0.122903) (14.93467,0.123091) (15.06030,0.123095) (15.18593,0.122914) (15.31156,0.122548) (15.43719,0.122001) (15.56281,0.121273) (15.68844,0.120369) (15.81407,0.119292) (15.93970,0.118047) (16.06533,0.116639) (16.19095,0.115075) (16.31658,0.113362) (16.44221,0.111506) (16.56784,0.109517) (16.69347,0.107401) (16.81910,0.105167) (16.94472,0.102826) (17.07035,0.100385) (17.19598,0.097856) (17.32161,0.095247) (17.44724,0.092568) (17.57286,0.089829) (17.69849,0.087041) (17.82412,0.084212) (17.94975,0.081353) (18.07538,0.078473) (18.20101,0.075581) (18.32663,0.072686) (18.45226,0.069798) (18.57789,0.066923) (18.70352,0.064071) (18.82915,0.061247) (18.95477,0.058461) (19.08040,0.055717) (19.20603,0.053022) (19.33166,0.050382) (19.45729,0.047802) (19.58291,0.045285) (19.70854,0.042837) (19.83417,0.040460) (19.95980,0.038157) (20.08543,0.035932) (20.21106,0.033785) (20.33668,0.031719) (20.46231,0.029735) (20.58794,0.027833) (20.71357,0.026013) (20.83920,0.024276) (20.96482,0.022621) (21.09045,0.021047) (21.21608,0.019553) (21.34171,0.018138) (21.46734,0.016800) (21.59296,0.015538) (21.71859,0.014348) (21.84422,0.013230) (21.96985,0.012181) (22.09548,0.011198) (22.22111,0.010279) (22.34673,0.009421) (22.47236,0.008622) (22.59799,0.007878) (22.72362,0.007188) (22.84925,0.006549) (22.97487,0.005957) (23.10050,0.005411) (23.22613,0.004908) (23.35176,0.004444) (23.47739,0.004019) (23.60302,0.003628) (23.72864,0.003271) (23.85427,0.002944) (23.97990,0.002646) (24.10553,0.002375) (24.23116,0.002128) (24.35678,0.001904) (24.48241,0.001701) (24.60804,0.001518) (24.73367,0.001352) (24.85930,0.001202) (24.98492,0.001068) (25.11055,0.000947) (25.23618,0.000838) (25.36181,0.000741) (25.48744,0.000654) (25.61307,0.000577) (25.73869,0.000507) (25.86432,0.000446) (25.98995,0.000391) (26.11558,0.000343) (26.24121,0.000300) (26.36683,0.000262) (26.49246,0.000228) (26.61809,0.000199) (26.74372,0.000173) (26.86935,0.000150) (26.99497,0.000130) (27.12060,0.000113) (27.24623,0.000097) (27.37186,0.000084) (27.49749,0.000072) (27.62312,0.000062) (27.74874,0.000054) (27.87437,0.000046) (28.00000,0.000039)};
\end{scope}
\draw[->] (0,0) -- (1.02,0);
\draw[->] (0,0) -- (0,1.02);
\draw (0.08,0) -- (0.08,-.012) node[below,font=\small] {$ 5 $};
\draw (0.28,0) -- (0.28,-.012) node[below,font=\small] {$ 10 $};
\draw (0.48,0) -- (0.48,-.012) node[below,font=\small] {$ 15 $};
\draw (0.68,0) -- (0.68,-.012) node[below,font=\small] {$ 20 $};
\draw (0.88,0) -- (0.88,-.012) node[below,font=\small] {$ 25 $};
\draw (0,0) -- (-.012,0) node[left,font=\small] {$ 0 $};
\draw (0,0.295202952) -- (-.012,0.295202952) node[left,font=\small] {$ 0.04 $};
\draw (0,0.5904059041) -- (-.012,0.5904059041) node[left,font=\small] {$ 0.08 $};
\draw (0,0.8856088561) -- (-.012,0.8856088561) node[left,font=\small] {$ 0.12 $};
\node[below,font=\small] at (.5,-.09) {成功次数};
\node[rotate=90,font=\small] at (-.13,.5) {单位宽度概率或密度};
\node[anchor=north east,fill=white,inner sep=3pt,font=\footnotesize] at (.98,.98) {%
\begin{tabular}{@{}ll@{}}
\tikz[baseline=-.5ex]{\filldraw[fill=gray!25,draw=gray] (0,-.05) rectangle (.5,.05);} & 二项概率\\
\tikz[baseline=-.5ex]{\draw[black,thick] (0,0) -- (.5,0);} & 正态近似\\
\end{tabular}};
\end{tikzpicture}\end{center}

本图采用 $n=50,p=0.3$。柱子表示各整数的概率，可理解为宽为一的直方图；曲线表示均值 $np$、方差 $np(1-p)$ 的正态密度。直接把一般离散单点概率与任意宽度的密度高度比较没有意义，这里以单位宽度建立可比尺度。

近似 $P(X\ge k)$ 时，将连续阈值放在 $k-1/2$，使连续区域覆盖完整整数柱。近似 $P(a\le X\le b)$ 时使用区间 $[a-1/2,b+1/2]$。修正改善许多中等容量计算，但不提供统一的误差保证。参见习题~\ref{q:M-P8-1ab}。

\clearpage\subsection*{图解三：零相关与依赖形状}

\begin{center}% 使用共享样式已加载的 TikZ，保留原图数据。
\begin{tikzpicture}[x=.69\linewidth,y=.69\linewidth]
% 直接设置数据坐标单位，避免先形成大尺寸再缩放造成溢出。
\begin{scope}[x=0.287500000023\linewidth,y=0.287500000023\linewidth,shift={(1.2,1.2)}]
\clip (-1.2,-1.2) rectangle (1.2,1.2);
\draw[only marks,mark=*,mark size=.65pt,black] plot coordinates {(1.00000,0.000000) (0.99939,0.034899) (0.99756,0.069756) (0.99452,0.104528) (0.99027,0.139173) (0.98481,0.173648) (0.97815,0.207912) (0.97030,0.241922) (0.96126,0.275637) (0.95106,0.309017) (0.93969,0.342020) (0.92718,0.374607) (0.91355,0.406737) (0.89879,0.438371) (0.88295,0.469472) (0.86603,0.500000) (0.84805,0.529919) (0.82904,0.559193) (0.80902,0.587785) (0.78801,0.615661) (0.76604,0.642788) (0.74314,0.669131) (0.71934,0.694658) (0.69466,0.719340) (0.66913,0.743145) (0.64279,0.766044) (0.61566,0.788011) (0.58779,0.809017) (0.55919,0.829038) (0.52992,0.848048) (0.50000,0.866025) (0.46947,0.882948) (0.43837,0.898794) (0.40674,0.913545) (0.37461,0.927184) (0.34202,0.939693) (0.30902,0.951057) (0.27564,0.961262) (0.24192,0.970296) (0.20791,0.978148) (0.17365,0.984808) (0.13917,0.990268) (0.10453,0.994522) (0.06976,0.997564) (0.03490,0.999391) (0.00000,1.000000) (-0.03490,0.999391) (-0.06976,0.997564) (-0.10453,0.994522) (-0.13917,0.990268) (-0.17365,0.984808) (-0.20791,0.978148) (-0.24192,0.970296) (-0.27564,0.961262) (-0.30902,0.951057) (-0.34202,0.939693) (-0.37461,0.927184) (-0.40674,0.913545) (-0.43837,0.898794) (-0.46947,0.882948) (-0.50000,0.866025) (-0.52992,0.848048) (-0.55919,0.829038) (-0.58779,0.809017) (-0.61566,0.788011) (-0.64279,0.766044) (-0.66913,0.743145) (-0.69466,0.719340) (-0.71934,0.694658) (-0.74314,0.669131) (-0.76604,0.642788) (-0.78801,0.615661) (-0.80902,0.587785) (-0.82904,0.559193) (-0.84805,0.529919) (-0.86603,0.500000) (-0.88295,0.469472) (-0.89879,0.438371) (-0.91355,0.406737) (-0.92718,0.374607) (-0.93969,0.342020) (-0.95106,0.309017) (-0.96126,0.275637) (-0.97030,0.241922) (-0.97815,0.207912) (-0.98481,0.173648) (-0.99027,0.139173) (-0.99452,0.104528) (-0.99756,0.069756) (-0.99939,0.034899) (-1.00000,0.000000) (-0.99939,-0.034899) (-0.99756,-0.069756) (-0.99452,-0.104528) (-0.99027,-0.139173) (-0.98481,-0.173648) (-0.97815,-0.207912) (-0.97030,-0.241922) (-0.96126,-0.275637) (-0.95106,-0.309017) (-0.93969,-0.342020) (-0.92718,-0.374607) (-0.91355,-0.406737) (-0.89879,-0.438371) (-0.88295,-0.469472) (-0.86603,-0.500000) (-0.84805,-0.529919) (-0.82904,-0.559193) (-0.80902,-0.587785) (-0.78801,-0.615661) (-0.76604,-0.642788) (-0.74314,-0.669131) (-0.71934,-0.694658) (-0.69466,-0.719340) (-0.66913,-0.743145) (-0.64279,-0.766044) (-0.61566,-0.788011) (-0.58779,-0.809017) (-0.55919,-0.829038) (-0.52992,-0.848048) (-0.50000,-0.866025) (-0.46947,-0.882948) (-0.43837,-0.898794) (-0.40674,-0.913545) (-0.37461,-0.927184) (-0.34202,-0.939693) (-0.30902,-0.951057) (-0.27564,-0.961262) (-0.24192,-0.970296) (-0.20791,-0.978148) (-0.17365,-0.984808) (-0.13917,-0.990268) (-0.10453,-0.994522) (-0.06976,-0.997564) (-0.03490,-0.999391) (-0.00000,-1.000000) (0.03490,-0.999391) (0.06976,-0.997564) (0.10453,-0.994522) (0.13917,-0.990268) (0.17365,-0.984808) (0.20791,-0.978148) (0.24192,-0.970296) (0.27564,-0.961262) (0.30902,-0.951057) (0.34202,-0.939693) (0.37461,-0.927184) (0.40674,-0.913545) (0.43837,-0.898794) (0.46947,-0.882948) (0.50000,-0.866025) (0.52992,-0.848048) (0.55919,-0.829038) (0.58779,-0.809017) (0.61566,-0.788011) (0.64279,-0.766044) (0.66913,-0.743145) (0.69466,-0.719340) (0.71934,-0.694658) (0.74314,-0.669131) (0.76604,-0.642788) (0.78801,-0.615661) (0.80902,-0.587785) (0.82904,-0.559193) (0.84805,-0.529919) (0.86603,-0.500000) (0.88295,-0.469472) (0.89879,-0.438371) (0.91355,-0.406737) (0.92718,-0.374607) (0.93969,-0.342020) (0.95106,-0.309017) (0.96126,-0.275637) (0.97030,-0.241922) (0.97815,-0.207912) (0.98481,-0.173648) (0.99027,-0.139173) (0.99452,-0.104528) (0.99756,-0.069756) (0.99939,-0.034899)};
\end{scope}
\draw[->] (0,0) -- (1.02,0);
\draw[->] (0,0) -- (0,1.02);
\draw (0.08333333333,0) -- (0.08333333333,-.012) node[below,font=\small] {$ -1 $};
\draw (0.5,0) -- (0.5,-.012) node[below,font=\small] {$ 0 $};
\draw (0.9166666667,0) -- (0.9166666667,-.012) node[below,font=\small] {$ 1 $};
\draw (0,0.08333333333) -- (-.012,0.08333333333) node[left,font=\small] {$ -1 $};
\draw (0,0.5) -- (-.012,0.5) node[left,font=\small] {$ 0 $};
\draw (0,0.9166666667) -- (-.012,0.9166666667) node[left,font=\small] {$ 1 $};
\node[below,font=\small] at (.5,-.09) {$X$};
\node[rotate=90,font=\small] at (-.13,.5) {$Y$};
\end{tikzpicture}\end{center}

图示单位圆周上的等角度点，用来展示支撑形状；它不是一份用于估计相关系数的独立随机样本。圆周均匀模型中两个坐标的协方差为零，但满足 $X^2+Y^2=1$ 的确定关系，因此不独立。给定横坐标的绝对值，纵坐标的绝对值已经确定。

相关系数只概括线性共变，不能识别所有非线性依赖。把散点图与模型支撑一起看，比单报一个接近零的相关系数更有信息。联合正态是零相关推出独立的特殊结构，圆周分布不具有该结构。参见习题~\ref{q:A5.35}、\ref{q:B17.15}。

\clearpage\subsection*{图解四：指数样本平均的正态化}

\begin{center}% 使用共享样式已加载的 TikZ，保留原图数据。
\begin{tikzpicture}[x=.87\linewidth,y=.69\linewidth]
% 直接设置数据坐标单位，避免先形成大尺寸再缩放造成溢出。
\begin{scope}[x=0.115999999971\linewidth,y=0.627272727279\linewidth,shift={(3.5,0)}]
\clip (-3.5,0) rectangle (4,1.1);
\draw[black,dotted,thick] plot coordinates {(-3.50000,0.000000) (-3.46575,0.000000) (-3.43151,0.000000) (-3.39726,0.000000) (-3.36301,0.000000) (-3.32877,0.000000) (-3.29452,0.000000) (-3.26027,0.000000) (-3.22603,0.000000) (-3.19178,0.000000) (-3.15753,0.000000) (-3.12329,0.000000) (-3.08904,0.000000) (-3.05479,0.000000) (-3.02055,0.000000) (-2.98630,0.000000) (-2.95205,0.000000) (-2.91781,0.000000) (-2.88356,0.000000) (-2.84932,0.000000) (-2.81507,0.000000) (-2.78082,0.000000) (-2.74658,0.000000) (-2.71233,0.000000) (-2.67808,0.000000) (-2.64384,0.000000) (-2.60959,0.000000) (-2.57534,0.000000) (-2.54110,0.000000) (-2.50685,0.000000) (-2.47260,0.000000) (-2.43836,0.000000) (-2.40411,0.000000) (-2.36986,0.000000) (-2.33562,0.000000) (-2.30137,0.000000) (-2.26712,0.000000) (-2.23288,0.000000) (-2.19863,0.000000) (-2.16438,0.000000) (-2.13014,0.000000) (-2.09589,0.000000) (-2.06164,0.000000) (-2.02740,0.000000) (-1.99315,0.000000) (-1.95890,0.000000) (-1.92466,0.000000) (-1.89041,0.000000) (-1.85616,0.000000) (-1.82192,0.000000) (-1.78767,0.000000) (-1.75342,0.000000) (-1.71918,0.000000) (-1.68493,0.000000) (-1.65068,0.000000) (-1.61644,0.000000) (-1.58219,0.000000) (-1.54795,0.000000) (-1.51370,0.000000) (-1.47945,0.000000) (-1.44521,0.000000) (-1.41096,0.000000) (-1.37671,0.000000) (-1.34247,0.000000) (-1.30822,0.000000) (-1.27397,0.000000) (-1.23973,0.000000) (-1.20548,0.000000) (-1.17123,0.000000) (-1.13699,0.000000) (-1.10274,0.000000) (-1.06849,0.000000) (-1.03425,0.000000) (-1.00000,0.000000) (-0.96575,0.966333) (-0.93151,0.933800) (-0.89726,0.902362) (-0.86301,0.871982) (-0.82877,0.842625) (-0.79452,0.814257) (-0.76027,0.786843) (-0.72603,0.760353) (-0.69178,0.734754) (-0.65753,0.710017) (-0.62329,0.686113) (-0.58904,0.663014) (-0.55479,0.640693) (-0.52055,0.619123) (-0.48630,0.598279) (-0.45205,0.578137) (-0.41781,0.558673) (-0.38356,0.539864) (-0.34932,0.521688) (-0.31507,0.504125) (-0.28082,0.487152) (-0.24658,0.470752) (-0.21233,0.454903) (-0.17808,0.439588) (-0.14384,0.424788) (-0.10959,0.410487) (-0.07534,0.396667) (-0.04110,0.383313) (-0.00685,0.370408) (0.02740,0.357937) (0.06164,0.345887) (0.09589,0.334242) (0.13014,0.322989) (0.16438,0.312115) (0.19863,0.301607) (0.23288,0.291453) (0.26712,0.281641) (0.30137,0.272159) (0.33562,0.262996) (0.36986,0.254142) (0.40411,0.245586) (0.43836,0.237318) (0.47260,0.229328) (0.50685,0.221607) (0.54110,0.214146) (0.57534,0.206937) (0.60959,0.199970) (0.64384,0.193237) (0.67808,0.186732) (0.71233,0.180445) (0.74658,0.174370) (0.78082,0.168500) (0.81507,0.162827) (0.84932,0.157345) (0.88356,0.152048) (0.91781,0.146929) (0.95205,0.141982) (0.98630,0.137202) (1.02055,0.132583) (1.05479,0.128119) (1.08904,0.123806) (1.12329,0.119638) (1.15753,0.115610) (1.19178,0.111718) (1.22603,0.107956) (1.26027,0.104322) (1.29452,0.100810) (1.32877,0.097416) (1.36301,0.094136) (1.39726,0.090967) (1.43151,0.087904) (1.46575,0.084945) (1.50000,0.082085) (1.53425,0.079321) (1.56849,0.076651) (1.60274,0.074070) (1.63699,0.071577) (1.67123,0.069167) (1.70548,0.066838) (1.73973,0.064588) (1.77397,0.062414) (1.80822,0.060312) (1.84247,0.058282) (1.87671,0.056320) (1.91096,0.054424) (1.94521,0.052591) (1.97945,0.050821) (2.01370,0.049110) (2.04795,0.047456) (2.08219,0.045859) (2.11644,0.044315) (2.15068,0.042823) (2.18493,0.041381) (2.21918,0.039988) (2.25342,0.038642) (2.28767,0.037341) (2.32192,0.036084) (2.35616,0.034869) (2.39041,0.033695) (2.42466,0.032560) (2.45890,0.031464) (2.49315,0.030405) (2.52740,0.029381) (2.56164,0.028392) (2.59589,0.027436) (2.63014,0.026513) (2.66438,0.025620) (2.69863,0.024757) (2.73288,0.023924) (2.76712,0.023118) (2.80137,0.022340) (2.83562,0.021588) (2.86986,0.020861) (2.90411,0.020159) (2.93836,0.019480) (2.97260,0.018824) (3.00685,0.018191) (3.04110,0.017578) (3.07534,0.016986) (3.10959,0.016415) (3.14384,0.015862) (3.17808,0.015328) (3.21233,0.014812) (3.24658,0.014313) (3.28082,0.013831) (3.31507,0.013366) (3.34932,0.012916) (3.38356,0.012481) (3.41781,0.012061) (3.45205,0.011655) (3.48630,0.011262) (3.52055,0.010883) (3.55479,0.010517) (3.58904,0.010163) (3.62329,0.009820) (3.65753,0.009490) (3.69178,0.009170) (3.72603,0.008862) (3.76027,0.008563) (3.79452,0.008275) (3.82877,0.007996) (3.86301,0.007727) (3.89726,0.007467) (3.93151,0.007216) (3.96575,0.006973) (4.00000,0.006738)};
\draw[black,dashed,thick] plot coordinates {(-3.50000,0.000000) (-3.46575,0.000000) (-3.43151,0.000000) (-3.39726,0.000000) (-3.36301,0.000000) (-3.32877,0.000000) (-3.29452,0.000000) (-3.26027,0.000000) (-3.22603,0.000000) (-3.19178,0.000000) (-3.15753,0.000000) (-3.12329,0.000000) (-3.08904,0.000000) (-3.05479,0.000000) (-3.02055,0.000000) (-2.98630,0.000000) (-2.95205,0.000000) (-2.91781,0.000000) (-2.88356,0.000000) (-2.84932,0.000000) (-2.81507,0.000000) (-2.78082,0.000000) (-2.74658,0.000000) (-2.71233,0.000000) (-2.67808,0.000000) (-2.64384,0.000000) (-2.60959,0.000000) (-2.57534,0.000000) (-2.54110,0.000000) (-2.50685,0.000000) (-2.47260,0.000000) (-2.43836,0.000000) (-2.40411,0.000000) (-2.36986,0.000000) (-2.33562,0.000000) (-2.30137,0.000000) (-2.26712,0.000000) (-2.23288,0.000000) (-2.19863,0.000004) (-2.16438,0.000052) (-2.13014,0.000231) (-2.09589,0.000657) (-2.06164,0.001460) (-2.02740,0.002770) (-1.99315,0.004711) (-1.95890,0.007396) (-1.92466,0.010918) (-1.89041,0.015351) (-1.85616,0.020748) (-1.82192,0.027143) (-1.78767,0.034548) (-1.75342,0.042955) (-1.71918,0.052342) (-1.68493,0.062666) (-1.65068,0.073876) (-1.61644,0.085904) (-1.58219,0.098675) (-1.54795,0.112107) (-1.51370,0.126110) (-1.47945,0.140591) (-1.44521,0.155454) (-1.41096,0.170603) (-1.37671,0.185941) (-1.34247,0.201374) (-1.30822,0.216809) (-1.27397,0.232158) (-1.23973,0.247336) (-1.20548,0.262264) (-1.17123,0.276866) (-1.13699,0.291073) (-1.10274,0.304823) (-1.06849,0.318058) (-1.03425,0.330728) (-1.00000,0.342786) (-0.96575,0.354194) (-0.93151,0.364919) (-0.89726,0.374934) (-0.86301,0.384216) (-0.82877,0.392749) (-0.79452,0.400521) (-0.76027,0.407526) (-0.72603,0.413761) (-0.69178,0.419228) (-0.65753,0.423932) (-0.62329,0.427882) (-0.58904,0.431091) (-0.55479,0.433573) (-0.52055,0.435346) (-0.48630,0.436430) (-0.45205,0.436847) (-0.41781,0.436620) (-0.38356,0.435774) (-0.34932,0.434336) (-0.31507,0.432332) (-0.28082,0.429790) (-0.24658,0.426740) (-0.21233,0.423209) (-0.17808,0.419226) (-0.14384,0.414822) (-0.10959,0.410024) (-0.07534,0.404861) (-0.04110,0.399361) (-0.00685,0.393553) (0.02740,0.387464) (0.06164,0.381119) (0.09589,0.374545) (0.13014,0.367767) (0.16438,0.360810) (0.19863,0.353696) (0.23288,0.346448) (0.26712,0.339087) (0.30137,0.331634) (0.33562,0.324109) (0.36986,0.316530) (0.40411,0.308914) (0.43836,0.301279) (0.47260,0.293641) (0.50685,0.286012) (0.54110,0.278409) (0.57534,0.270842) (0.60959,0.263325) (0.64384,0.255868) (0.67808,0.248482) (0.71233,0.241176) (0.74658,0.233958) (0.78082,0.226837) (0.81507,0.219819) (0.84932,0.212911) (0.88356,0.206118) (0.91781,0.199446) (0.95205,0.192899) (0.98630,0.186481) (1.02055,0.180194) (1.05479,0.174043) (1.08904,0.168029) (1.12329,0.162154) (1.15753,0.156419) (1.19178,0.150826) (1.22603,0.145374) (1.26027,0.140065) (1.29452,0.134898) (1.32877,0.129873) (1.36301,0.124988) (1.39726,0.120244) (1.43151,0.115639) (1.46575,0.111171) (1.50000,0.106840) (1.53425,0.102642) (1.56849,0.098577) (1.60274,0.094643) (1.63699,0.090836) (1.67123,0.087155) (1.70548,0.083598) (1.73973,0.080162) (1.77397,0.076844) (1.80822,0.073642) (1.84247,0.070553) (1.87671,0.067575) (1.91096,0.064704) (1.94521,0.061939) (1.97945,0.059276) (2.01370,0.056712) (2.04795,0.054245) (2.08219,0.051872) (2.11644,0.049591) (2.15068,0.047398) (2.18493,0.045291) (2.21918,0.043267) (2.25342,0.041324) (2.28767,0.039459) (2.32192,0.037670) (2.35616,0.035953) (2.39041,0.034307) (2.42466,0.032730) (2.45890,0.031218) (2.49315,0.029769) (2.52740,0.028382) (2.56164,0.027054) (2.59589,0.025783) (2.63014,0.024567) (2.66438,0.023403) (2.69863,0.022290) (2.73288,0.021226) (2.76712,0.020209) (2.80137,0.019237) (2.83562,0.018308) (2.86986,0.017421) (2.90411,0.016574) (2.93836,0.015766) (2.97260,0.014994) (3.00685,0.014258) (3.04110,0.013555) (3.07534,0.012885) (3.10959,0.012246) (3.14384,0.011637) (3.17808,0.011056) (3.21233,0.010502) (3.24658,0.009975) (3.28082,0.009473) (3.31507,0.008994) (3.34932,0.008539) (3.38356,0.008105) (3.41781,0.007692) (3.45205,0.007299) (3.48630,0.006926) (3.52055,0.006570) (3.55479,0.006232) (3.58904,0.005910) (3.62329,0.005604) (3.65753,0.005314) (3.69178,0.005037) (3.72603,0.004775) (3.76027,0.004525) (3.79452,0.004288) (3.82877,0.004063) (3.86301,0.003849) (3.89726,0.003646) (3.93151,0.003454) (3.96575,0.003271) (4.00000,0.003097)};
\draw[black,thick] plot coordinates {(-3.50000,0.000012) (-3.46575,0.000017) (-3.43151,0.000023) (-3.39726,0.000031) (-3.36301,0.000041) (-3.32877,0.000054) (-3.29452,0.000071) (-3.26027,0.000092) (-3.22603,0.000119) (-3.19178,0.000153) (-3.15753,0.000195) (-3.12329,0.000247) (-3.08904,0.000311) (-3.05479,0.000390) (-3.02055,0.000486) (-2.98630,0.000602) (-2.95205,0.000741) (-2.91781,0.000908) (-2.88356,0.001107) (-2.84932,0.001342) (-2.81507,0.001620) (-2.78082,0.001945) (-2.74658,0.002326) (-2.71233,0.002767) (-2.67808,0.003278) (-2.64384,0.003866) (-2.60959,0.004541) (-2.57534,0.005311) (-2.54110,0.006187) (-2.50685,0.007179) (-2.47260,0.008299) (-2.43836,0.009556) (-2.40411,0.010964) (-2.36986,0.012534) (-2.33562,0.014278) (-2.30137,0.016208) (-2.26712,0.018339) (-2.23288,0.020680) (-2.19863,0.023245) (-2.16438,0.026046) (-2.13014,0.029095) (-2.09589,0.032401) (-2.06164,0.035976) (-2.02740,0.039829) (-1.99315,0.043969) (-1.95890,0.048404) (-1.92466,0.053139) (-1.89041,0.058181) (-1.85616,0.063533) (-1.82192,0.069197) (-1.78767,0.075175) (-1.75342,0.081466) (-1.71918,0.088066) (-1.68493,0.094973) (-1.65068,0.102179) (-1.61644,0.109677) (-1.58219,0.117456) (-1.54795,0.125506) (-1.51370,0.133813) (-1.47945,0.142360) (-1.44521,0.151132) (-1.41096,0.160108) (-1.37671,0.169270) (-1.34247,0.178594) (-1.30822,0.188057) (-1.27397,0.197634) (-1.23973,0.207299) (-1.20548,0.217026) (-1.17123,0.226786) (-1.13699,0.236551) (-1.10274,0.246291) (-1.06849,0.255976) (-1.03425,0.265577) (-1.00000,0.275065) (-0.96575,0.284408) (-0.93151,0.293577) (-0.89726,0.302544) (-0.86301,0.311279) (-0.82877,0.319755) (-0.79452,0.327946) (-0.76027,0.335824) (-0.72603,0.343367) (-0.69178,0.350551) (-0.65753,0.357354) (-0.62329,0.363756) (-0.58904,0.369738) (-0.55479,0.375285) (-0.52055,0.380380) (-0.48630,0.385011) (-0.45205,0.389166) (-0.41781,0.392837) (-0.38356,0.396016) (-0.34932,0.398698) (-0.31507,0.400879) (-0.28082,0.402558) (-0.24658,0.403735) (-0.21233,0.404412) (-0.17808,0.404594) (-0.14384,0.404286) (-0.10959,0.403495) (-0.07534,0.402231) (-0.04110,0.400503) (-0.00685,0.398324) (0.02740,0.395708) (0.06164,0.392667) (0.09589,0.389218) (0.13014,0.385377) (0.16438,0.381161) (0.19863,0.376589) (0.23288,0.371679) (0.26712,0.366450) (0.30137,0.360923) (0.33562,0.355118) (0.36986,0.349054) (0.40411,0.342752) (0.43836,0.336234) (0.47260,0.329519) (0.50685,0.322628) (0.54110,0.315581) (0.57534,0.308399) (0.60959,0.301100) (0.64384,0.293704) (0.67808,0.286230) (0.71233,0.278696) (0.74658,0.271119) (0.78082,0.263517) (0.81507,0.255906) (0.84932,0.248301) (0.88356,0.240717) (0.91781,0.233169) (0.95205,0.225670) (0.98630,0.218232) (1.02055,0.210868) (1.05479,0.203589) (1.08904,0.196403) (1.12329,0.189323) (1.15753,0.182354) (1.19178,0.175507) (1.22603,0.168788) (1.26027,0.162203) (1.29452,0.155758) (1.32877,0.149459) (1.36301,0.143309) (1.39726,0.137312) (1.43151,0.131472) (1.46575,0.125790) (1.50000,0.120269) (1.53425,0.114910) (1.56849,0.109713) (1.60274,0.104680) (1.63699,0.099810) (1.67123,0.095103) (1.70548,0.090558) (1.73973,0.086173) (1.77397,0.081946) (1.80822,0.077877) (1.84247,0.073962) (1.87671,0.070200) (1.91096,0.066587) (1.94521,0.063120) (1.97945,0.059797) (2.01370,0.056615) (2.04795,0.053569) (2.08219,0.050657) (2.11644,0.047874) (2.15068,0.045218) (2.18493,0.042684) (2.21918,0.040269) (2.25342,0.037969) (2.28767,0.035779) (2.32192,0.033697) (2.35616,0.031718) (2.39041,0.029839) (2.42466,0.028056) (2.45890,0.026365) (2.49315,0.024763) (2.52740,0.023245) (2.56164,0.021809) (2.59589,0.020451) (2.63014,0.019168) (2.66438,0.017955) (2.69863,0.016811) (2.73288,0.015732) (2.76712,0.014714) (2.80137,0.013756) (2.83562,0.012853) (2.86986,0.012004) (2.90411,0.011206) (2.93836,0.010455) (2.97260,0.009750) (3.00685,0.009089) (3.04110,0.008468) (3.07534,0.007886) (3.10959,0.007340) (3.14384,0.006829) (3.17808,0.006351) (3.21233,0.005904) (3.24658,0.005485) (3.28082,0.005094) (3.31507,0.004729) (3.34932,0.004388) (3.38356,0.004070) (3.41781,0.003773) (3.45205,0.003496) (3.48630,0.003239) (3.52055,0.002999) (3.55479,0.002775) (3.58904,0.002568) (3.62329,0.002374) (3.65753,0.002195) (3.69178,0.002028) (3.72603,0.001873) (3.76027,0.001729) (3.79452,0.001596) (3.82877,0.001472) (3.86301,0.001358) (3.89726,0.001251) (3.93151,0.001153) (3.96575,0.001062) (4.00000,0.000978)};
\draw[gray,very thick] plot coordinates {(-3.50000,0.000873) (-3.46575,0.000983) (-3.43151,0.001106) (-3.39726,0.001244) (-3.36301,0.001396) (-3.32877,0.001566) (-3.29452,0.001754) (-3.26027,0.001962) (-3.22603,0.002193) (-3.19178,0.002448) (-3.15753,0.002729) (-3.12329,0.003038) (-3.08904,0.003380) (-3.05479,0.003754) (-3.02055,0.004166) (-2.98630,0.004617) (-2.95205,0.005112) (-2.91781,0.005652) (-2.88356,0.006242) (-2.84932,0.006886) (-2.81507,0.007588) (-2.78082,0.008351) (-2.74658,0.009180) (-2.71233,0.010079) (-2.67808,0.011054) (-2.64384,0.012108) (-2.60959,0.013248) (-2.57534,0.014478) (-2.54110,0.015804) (-2.50685,0.017230) (-2.47260,0.018764) (-2.43836,0.020410) (-2.40411,0.022175) (-2.36986,0.024063) (-2.33562,0.026083) (-2.30137,0.028238) (-2.26712,0.030536) (-2.23288,0.032982) (-2.19863,0.035582) (-2.16438,0.038342) (-2.13014,0.041267) (-2.09589,0.044364) (-2.06164,0.047638) (-2.02740,0.051093) (-1.99315,0.054734) (-1.95890,0.058567) (-1.92466,0.062594) (-1.89041,0.066819) (-1.85616,0.071246) (-1.82192,0.075878) (-1.78767,0.080716) (-1.75342,0.085761) (-1.71918,0.091016) (-1.68493,0.096478) (-1.65068,0.102149) (-1.61644,0.108027) (-1.58219,0.114109) (-1.54795,0.120392) (-1.51370,0.126872) (-1.47945,0.133544) (-1.44521,0.140402) (-1.41096,0.147439) (-1.37671,0.154647) (-1.34247,0.162018) (-1.30822,0.169542) (-1.27397,0.177206) (-1.23973,0.185000) (-1.20548,0.192911) (-1.17123,0.200923) (-1.13699,0.209024) (-1.10274,0.217196) (-1.06849,0.225423) (-1.03425,0.233687) (-1.00000,0.241971) (-0.96575,0.250254) (-0.93151,0.258518) (-0.89726,0.266741) (-0.86301,0.274904) (-0.82877,0.282984) (-0.79452,0.290960) (-0.76027,0.298810) (-0.72603,0.306513) (-0.69178,0.314045) (-0.65753,0.321385) (-0.62329,0.328512) (-0.58904,0.335403) (-0.55479,0.342037) (-0.52055,0.348393) (-0.48630,0.354452) (-0.45205,0.360193) (-0.41781,0.365598) (-0.38356,0.370650) (-0.34932,0.375330) (-0.31507,0.379625) (-0.28082,0.383518) (-0.24658,0.386997) (-0.21233,0.390050) (-0.17808,0.392666) (-0.14384,0.394837) (-0.10959,0.396554) (-0.07534,0.397812) (-0.04110,0.398606) (-0.00685,0.398933) (0.02740,0.398793) (0.06164,0.398185) (0.09589,0.397112) (0.13014,0.395578) (0.16438,0.393588) (0.19863,0.391149) (0.23288,0.388270) (0.26712,0.384960) (0.30137,0.381231) (0.33562,0.377095) (0.36986,0.372567) (0.40411,0.367662) (0.43836,0.362396) (0.47260,0.356787) (0.50685,0.350853) (0.54110,0.344614) (0.57534,0.338088) (0.60959,0.331298) (0.64384,0.324263) (0.67808,0.317005) (0.71233,0.309547) (0.74658,0.301910) (0.78082,0.294116) (0.81507,0.286188) (0.84932,0.278147) (0.88356,0.270015) (0.91781,0.261813) (0.95205,0.253563) (0.98630,0.245285) (1.02055,0.236999) (1.05479,0.228725) (1.08904,0.220481) (1.12329,0.212285) (1.15753,0.204154) (1.19178,0.196104) (1.22603,0.188151) (1.26027,0.180309) (1.29452,0.172591) (1.32877,0.165010) (1.36301,0.157577) (1.39726,0.150302) (1.43151,0.143195) (1.46575,0.136265) (1.50000,0.129518) (1.53425,0.122960) (1.56849,0.116598) (1.60274,0.110435) (1.63699,0.104476) (1.67123,0.098722) (1.70548,0.093176) (1.73973,0.087838) (1.77397,0.082709) (1.80822,0.077788) (1.84247,0.073074) (1.87671,0.068566) (1.91096,0.064260) (1.94521,0.060154) (1.97945,0.056244) (2.01370,0.052527) (2.04795,0.048998) (2.08219,0.045652) (2.11644,0.042485) (2.15068,0.039492) (2.18493,0.036666) (2.21918,0.034003) (2.25342,0.031496) (2.28767,0.029140) (2.32192,0.026928) (2.35616,0.024855) (2.39041,0.022915) (2.42466,0.021101) (2.45890,0.019409) (2.49315,0.017831) (2.52740,0.016362) (2.56164,0.014996) (2.59589,0.013729) (2.63014,0.012554) (2.66438,0.011466) (2.69863,0.010460) (2.73288,0.009531) (2.76712,0.008674) (2.80137,0.007885) (2.83562,0.007160) (2.86986,0.006493) (2.90411,0.005882) (2.93836,0.005322) (2.97260,0.004810) (3.00685,0.004342) (3.04110,0.003914) (3.07534,0.003525) (3.10959,0.003171) (3.14384,0.002849) (3.17808,0.002557) (3.21233,0.002292) (3.24658,0.002052) (3.28082,0.001835) (3.31507,0.001639) (3.34932,0.001462) (3.38356,0.001303) (3.41781,0.001160) (3.45205,0.001031) (3.48630,0.000915) (3.52055,0.000812) (3.55479,0.000719) (3.58904,0.000636) (3.62329,0.000563) (3.65753,0.000497) (3.69178,0.000438) (3.72603,0.000386) (3.76027,0.000339) (3.79452,0.000298) (3.82877,0.000262) (3.86301,0.000229) (3.89726,0.000201) (3.93151,0.000176) (3.96575,0.000153) (4.00000,0.000134)};
\end{scope}
\draw[->] (0,0) -- (1.02,0);
\draw[->] (0,0) -- (0,1.02);
\draw (0.06666666667,0) -- (0.06666666667,-.012) node[below,font=\small] {$ -3 $};
\draw (0.2,0) -- (0.2,-.012) node[below,font=\small] {$ -2 $};
\draw (0.3333333333,0) -- (0.3333333333,-.012) node[below,font=\small] {$ -1 $};
\draw (0.4666666667,0) -- (0.4666666667,-.012) node[below,font=\small] {$ 0 $};
\draw (0.6,0) -- (0.6,-.012) node[below,font=\small] {$ 1 $};
\draw (0.7333333333,0) -- (0.7333333333,-.012) node[below,font=\small] {$ 2 $};
\draw (0.8666666667,0) -- (0.8666666667,-.012) node[below,font=\small] {$ 3 $};
\draw (1,0) -- (1,-.012) node[below,font=\small] {$ 4 $};
\draw (0,0) -- (-.012,0) node[left,font=\small] {$ 0 $};
\draw (0,0.1818181818) -- (-.012,0.1818181818) node[left,font=\small] {$ 0.2 $};
\draw (0,0.3636363636) -- (-.012,0.3636363636) node[left,font=\small] {$ 0.4 $};
\draw (0,0.5454545455) -- (-.012,0.5454545455) node[left,font=\small] {$ 0.6 $};
\draw (0,0.7272727273) -- (-.012,0.7272727273) node[left,font=\small] {$ 0.8 $};
\draw (0,0.9090909091) -- (-.012,0.9090909091) node[left,font=\small] {$ 1 $};
\node[below,font=\small] at (.5,-.09) {标准化均值};
\node[rotate=90,font=\small] at (-.13,.5) {密度};
\node[anchor=north east,fill=white,inner sep=3pt,font=\footnotesize] at (.98,.98) {%
\begin{tabular}{@{}ll@{}}
\tikz[baseline=-.5ex]{\draw[black,dotted,thick] (0,0) -- (.5,0);} & $n=1$\\
\tikz[baseline=-.5ex]{\draw[black,dashed,thick] (0,0) -- (.5,0);} & $n=5$\\
\tikz[baseline=-.5ex]{\draw[black,thick] (0,0) -- (.5,0);} & $n=30$\\
\tikz[baseline=-.5ex]{\draw[gray,very thick] (0,0) -- (.5,0);} & 标准正态\\
\end{tabular}};
\end{tikzpicture}\end{center}

令原始观测服从率为一的指数分布，图中画的是 $\sqrt n(\overline X-1)$ 的精确伽马变换密度，及标准正态密度。样本容量增大后偏斜减弱，但有限样本的分布仍有其自身支撑和尾部特征。

横轴已经中心化并乘以 $\sqrt n$。若画未缩放的样本均值，曲线会同时向均值收缩，无法直接看清中心极限定理描述的波动形状。图中 $n=1$ 的左端跳跃来自指数变量非负，而非绘图误差。参见习题~\ref{q:A6.30}、\ref{q:B17.20}。

\clearpage\subsection*{图解五：方差统计量的偏斜}

\begin{center}% 使用共享样式已加载的 TikZ，保留原图数据。
\begin{tikzpicture}[x=.87\linewidth,y=.69\linewidth]
% 直接设置数据坐标单位，避免先形成大尺寸再缩放造成溢出。
\begin{scope}[x=0.248571428559\linewidth,y=0.394285714266\linewidth,shift={(0,0)}]
\clip (0,0) rectangle (3.5,1.75);
\draw[black,dotted,thick] plot coordinates {(0.00100,0.065455) (0.01698,0.263309) (0.03295,0.358163) (0.04893,0.426099) (0.06491,0.479136) (0.08089,0.522199) (0.09686,0.557918) (0.11284,0.587916) (0.12882,0.613285) (0.14479,0.634809) (0.16077,0.653075) (0.17675,0.668542) (0.19273,0.681573) (0.20870,0.692466) (0.22468,0.701469) (0.24066,0.708790) (0.25663,0.714607) (0.27261,0.719074) (0.28859,0.722325) (0.30457,0.724479) (0.32054,0.725638) (0.33652,0.725896) (0.35250,0.725335) (0.36847,0.724029) (0.38445,0.722046) (0.40043,0.719447) (0.41641,0.716285) (0.43238,0.712613) (0.44836,0.708476) (0.46434,0.703915) (0.48032,0.698969) (0.49629,0.693674) (0.51227,0.688062) (0.52825,0.682163) (0.54422,0.676006) (0.56020,0.669616) (0.57618,0.663016) (0.59216,0.656229) (0.60813,0.649274) (0.62411,0.642172) (0.64009,0.634939) (0.65606,0.627592) (0.67204,0.620146) (0.68802,0.612616) (0.70400,0.605013) (0.71997,0.597351) (0.73595,0.589641) (0.75193,0.581893) (0.76790,0.574117) (0.78388,0.566323) (0.79986,0.558518) (0.81584,0.550711) (0.83181,0.542909) (0.84779,0.535119) (0.86377,0.527347) (0.87974,0.519599) (0.89572,0.511880) (0.91170,0.504196) (0.92768,0.496550) (0.94365,0.488949) (0.95963,0.481394) (0.97561,0.473891) (0.99158,0.466442) (1.00756,0.459050) (1.02354,0.451719) (1.03952,0.444451) (1.05549,0.437248) (1.07147,0.430112) (1.08745,0.423046) (1.10342,0.416051) (1.11940,0.409129) (1.13538,0.402281) (1.15136,0.395508) (1.16733,0.388812) (1.18331,0.382194) (1.19929,0.375654) (1.21526,0.369193) (1.23124,0.362812) (1.24722,0.356511) (1.26320,0.350291) (1.27917,0.344152) (1.29515,0.338094) (1.31113,0.332117) (1.32711,0.326222) (1.34308,0.320408) (1.35906,0.314676) (1.37504,0.309024) (1.39101,0.303454) (1.40699,0.297965) (1.42297,0.292556) (1.43895,0.287227) (1.45492,0.281978) (1.47090,0.276808) (1.48688,0.271716) (1.50285,0.266703) (1.51883,0.261768) (1.53481,0.256910) (1.55079,0.252128) (1.56676,0.247422) (1.58274,0.242792) (1.59872,0.238236) (1.61469,0.233753) (1.63067,0.229344) (1.64665,0.225007) (1.66263,0.220742) (1.67860,0.216548) (1.69458,0.212424) (1.71056,0.208369) (1.72653,0.204382) (1.74251,0.200463) (1.75849,0.196611) (1.77447,0.192826) (1.79044,0.189105) (1.80642,0.185449) (1.82240,0.181856) (1.83837,0.178326) (1.85435,0.174858) (1.87033,0.171451) (1.88631,0.168105) (1.90228,0.164817) (1.91826,0.161589) (1.93424,0.158418) (1.95021,0.155304) (1.96619,0.152246) (1.98217,0.149243) (1.99815,0.146295) (2.01412,0.143401) (2.03010,0.140559) (2.04608,0.137769) (2.06205,0.135031) (2.07803,0.132343) (2.09401,0.129705) (2.10999,0.127116) (2.12596,0.124574) (2.14194,0.122080) (2.15792,0.119633) (2.17389,0.117232) (2.18987,0.114875) (2.20585,0.112564) (2.22183,0.110295) (2.23780,0.108070) (2.25378,0.105887) (2.26976,0.103745) (2.28574,0.101644) (2.30171,0.099583) (2.31769,0.097562) (2.33367,0.095579) (2.34964,0.093635) (2.36562,0.091728) (2.38160,0.089857) (2.39758,0.088023) (2.41355,0.086225) (2.42953,0.084461) (2.44551,0.082732) (2.46148,0.081036) (2.47746,0.079373) (2.49344,0.077743) (2.50942,0.076145) (2.52539,0.074578) (2.54137,0.073042) (2.55735,0.071536) (2.57332,0.070060) (2.58930,0.068613) (2.60528,0.067194) (2.62126,0.065804) (2.63723,0.064441) (2.65321,0.063105) (2.66919,0.061796) (2.68516,0.060513) (2.70114,0.059256) (2.71712,0.058023) (2.73310,0.056815) (2.74907,0.055632) (2.76505,0.054472) (2.78103,0.053336) (2.79700,0.052222) (2.81298,0.051131) (2.82896,0.050061) (2.84494,0.049014) (2.86091,0.047987) (2.87689,0.046982) (2.89287,0.045996) (2.90884,0.045031) (2.92482,0.044085) (2.94080,0.043158) (2.95678,0.042251) (2.97275,0.041361) (2.98873,0.040490) (3.00471,0.039637) (3.02068,0.038801) (3.03666,0.037982) (3.05264,0.037180) (3.06862,0.036395) (3.08459,0.035625) (3.10057,0.034872) (3.11655,0.034133) (3.13253,0.033410) (3.14850,0.032702) (3.16448,0.032009) (3.18046,0.031330) (3.19643,0.030664) (3.21241,0.030013) (3.22839,0.029375) (3.24437,0.028750) (3.26034,0.028138) (3.27632,0.027539) (3.29230,0.026953) (3.30827,0.026378) (3.32425,0.025816) (3.34023,0.025265) (3.35621,0.024725) (3.37218,0.024197) (3.38816,0.023680) (3.40414,0.023174) (3.42011,0.022678) (3.43609,0.022193) (3.45207,0.021718) (3.46805,0.021252) (3.48402,0.020797) (3.50000,0.020351)};
\draw[black,dashed,thick] plot coordinates {(0.00100,0.000000) (0.01698,0.000010) (0.03295,0.000130) (0.04893,0.000584) (0.06491,0.001671) (0.08089,0.003720) (0.09686,0.007062) (0.11284,0.012008) (0.12882,0.018828) (0.14479,0.027748) (0.16077,0.038937) (0.17675,0.052512) (0.19273,0.068533) (0.20870,0.087009) (0.22468,0.107898) (0.24066,0.131116) (0.25663,0.156540) (0.27261,0.184013) (0.28859,0.213351) (0.30457,0.244348) (0.32054,0.276780) (0.33652,0.310413) (0.35250,0.345005) (0.36847,0.380310) (0.38445,0.416084) (0.40043,0.452086) (0.41641,0.488083) (0.43238,0.523849) (0.44836,0.559173) (0.46434,0.593852) (0.48032,0.627701) (0.49629,0.660547) (0.51227,0.692235) (0.52825,0.722624) (0.54422,0.751592) (0.56020,0.779028) (0.57618,0.804843) (0.59216,0.828958) (0.60813,0.851312) (0.62411,0.871857) (0.64009,0.890560) (0.65606,0.907400) (0.67204,0.922366) (0.68802,0.935462) (0.70400,0.946698) (0.71997,0.956097) (0.73595,0.963689) (0.75193,0.969511) (0.76790,0.973608) (0.78388,0.976030) (0.79986,0.976834) (0.81584,0.976079) (0.83181,0.973829) (0.84779,0.970151) (0.86377,0.965115) (0.87974,0.958791) (0.89572,0.951253) (0.91170,0.942574) (0.92768,0.932827) (0.94365,0.922086) (0.95963,0.910424) (0.97561,0.897914) (0.99158,0.884625) (1.00756,0.870629) (1.02354,0.855992) (1.03952,0.840782) (1.05549,0.825061) (1.07147,0.808893) (1.08745,0.792337) (1.10342,0.775449) (1.11940,0.758285) (1.13538,0.740897) (1.15136,0.723335) (1.16733,0.705645) (1.18331,0.687872) (1.19929,0.670058) (1.21526,0.652242) (1.23124,0.634460) (1.24722,0.616748) (1.26320,0.599137) (1.27917,0.581657) (1.29515,0.564334) (1.31113,0.547194) (1.32711,0.530259) (1.34308,0.513551) (1.35906,0.497087) (1.37504,0.480886) (1.39101,0.464961) (1.40699,0.449327) (1.42297,0.433994) (1.43895,0.418973) (1.45492,0.404273) (1.47090,0.389900) (1.48688,0.375861) (1.50285,0.362160) (1.51883,0.348800) (1.53481,0.335785) (1.55079,0.323116) (1.56676,0.310792) (1.58274,0.298814) (1.59872,0.287181) (1.61469,0.275890) (1.63067,0.264940) (1.64665,0.254326) (1.66263,0.244046) (1.67860,0.234094) (1.69458,0.224468) (1.71056,0.215160) (1.72653,0.206166) (1.74251,0.197481) (1.75849,0.189098) (1.77447,0.181011) (1.79044,0.173214) (1.80642,0.165700) (1.82240,0.158462) (1.83837,0.151494) (1.85435,0.144788) (1.87033,0.138339) (1.88631,0.132137) (1.90228,0.126178) (1.91826,0.120454) (1.93424,0.114957) (1.95021,0.109681) (1.96619,0.104619) (1.98217,0.099765) (1.99815,0.095111) (2.01412,0.090650) (2.03010,0.086378) (2.04608,0.082286) (2.06205,0.078369) (2.07803,0.074621) (2.09401,0.071035) (2.10999,0.067605) (2.12596,0.064327) (2.14194,0.061193) (2.15792,0.058200) (2.17389,0.055340) (2.18987,0.052610) (2.20585,0.050004) (2.22183,0.047516) (2.23780,0.045144) (2.25378,0.042881) (2.26976,0.040723) (2.28574,0.038666) (2.30171,0.036706) (2.31769,0.034839) (2.33367,0.033060) (2.34964,0.031366) (2.36562,0.029754) (2.38160,0.028219) (2.39758,0.026758) (2.41355,0.025369) (2.42953,0.024048) (2.44551,0.022791) (2.46148,0.021596) (2.47746,0.020461) (2.49344,0.019382) (2.50942,0.018357) (2.52539,0.017383) (2.54137,0.016459) (2.55735,0.015581) (2.57332,0.014747) (2.58930,0.013956) (2.60528,0.013206) (2.62126,0.012494) (2.63723,0.011818) (2.65321,0.011178) (2.66919,0.010570) (2.68516,0.009995) (2.70114,0.009449) (2.71712,0.008932) (2.73310,0.008441) (2.74907,0.007977) (2.76505,0.007537) (2.78103,0.007121) (2.79700,0.006727) (2.81298,0.006353) (2.82896,0.006000) (2.84494,0.005665) (2.86091,0.005349) (2.87689,0.005049) (2.89287,0.004766) (2.90884,0.004498) (2.92482,0.004245) (2.94080,0.004005) (2.95678,0.003779) (2.97275,0.003565) (2.98873,0.003362) (3.00471,0.003171) (3.02068,0.002990) (3.03666,0.002820) (3.05264,0.002658) (3.06862,0.002506) (3.08459,0.002362) (3.10057,0.002226) (3.11655,0.002098) (3.13253,0.001977) (3.14850,0.001863) (3.16448,0.001755) (3.18046,0.001653) (3.19643,0.001557) (3.21241,0.001467) (3.22839,0.001381) (3.24437,0.001300) (3.26034,0.001224) (3.27632,0.001153) (3.29230,0.001085) (3.30827,0.001021) (3.32425,0.000961) (3.34023,0.000905) (3.35621,0.000851) (3.37218,0.000801) (3.38816,0.000754) (3.40414,0.000709) (3.42011,0.000667) (3.43609,0.000627) (3.45207,0.000590) (3.46805,0.000555) (3.48402,0.000522) (3.50000,0.000491)};
\draw[black,thick] plot coordinates {(0.00100,0.000000) (0.01698,0.000000) (0.03295,0.000000) (0.04893,0.000000) (0.06491,0.000000) (0.08089,0.000000) (0.09686,0.000000) (0.11284,0.000000) (0.12882,0.000000) (0.14479,0.000001) (0.16077,0.000003) (0.17675,0.000010) (0.19273,0.000027) (0.20870,0.000065) (0.22468,0.000144) (0.24066,0.000297) (0.25663,0.000575) (0.27261,0.001054) (0.28859,0.001840) (0.30457,0.003079) (0.32054,0.004957) (0.33652,0.007707) (0.35250,0.011610) (0.36847,0.016993) (0.38445,0.024225) (0.40043,0.033710) (0.41641,0.045873) (0.43238,0.061151) (0.44836,0.079972) (0.46434,0.102742) (0.48032,0.129824) (0.49629,0.161522) (0.51227,0.198062) (0.52825,0.239583) (0.54422,0.286118) (0.56020,0.337591) (0.57618,0.393813) (0.59216,0.454477) (0.60813,0.519166) (0.62411,0.587357) (0.64009,0.658434) (0.65606,0.731699) (0.67204,0.806392) (0.68802,0.881702) (0.70400,0.956794) (0.71997,1.030818) (0.73595,1.102938) (0.75193,1.172340) (0.76790,1.238257) (0.78388,1.299975) (0.79986,1.356853) (0.81584,1.408329) (0.83181,1.453929) (0.84779,1.493271) (0.86377,1.526070) (0.87974,1.552135) (0.89572,1.571372) (0.91170,1.583775) (0.92768,1.589427) (0.94365,1.588487) (0.95963,1.581189) (0.97561,1.567827) (0.99158,1.548751) (1.00756,1.524356) (1.02354,1.495072) (1.03952,1.461356) (1.05549,1.423684) (1.07147,1.382540) (1.08745,1.338413) (1.10342,1.291783) (1.11940,1.243123) (1.13538,1.192889) (1.15136,1.141515) (1.16733,1.089413) (1.18331,1.036965) (1.19929,0.984527) (1.21526,0.932420) (1.23124,0.880936) (1.24722,0.830335) (1.26320,0.780844) (1.27917,0.732661) (1.29515,0.685951) (1.31113,0.640851) (1.32711,0.597474) (1.34308,0.555903) (1.35906,0.516201) (1.37504,0.478408) (1.39101,0.442544) (1.40699,0.408614) (1.42297,0.376604) (1.43895,0.346490) (1.45492,0.318234) (1.47090,0.291789) (1.48688,0.267100) (1.50285,0.244106) (1.51883,0.222738) (1.53481,0.202926) (1.55079,0.184596) (1.56676,0.167672) (1.58274,0.152078) (1.59872,0.137738) (1.61469,0.124576) (1.63067,0.112517) (1.64665,0.101489) (1.66263,0.091421) (1.67860,0.082246) (1.69458,0.073897) (1.71056,0.066314) (1.72653,0.059436) (1.74251,0.053208) (1.75849,0.047576) (1.77447,0.042491) (1.79044,0.037907) (1.80642,0.033779) (1.82240,0.030068) (1.83837,0.026736) (1.85435,0.023748) (1.87033,0.021072) (1.88631,0.018679) (1.90228,0.016540) (1.91826,0.014632) (1.93424,0.012932) (1.95021,0.011418) (1.96619,0.010072) (1.98217,0.008877) (1.99815,0.007816) (2.01412,0.006876) (2.03010,0.006043) (2.04608,0.005307) (2.06205,0.004657) (2.07803,0.004082) (2.09401,0.003576) (2.10999,0.003130) (2.12596,0.002737) (2.14194,0.002392) (2.15792,0.002089) (2.17389,0.001822) (2.18987,0.001589) (2.20585,0.001384) (2.22183,0.001205) (2.23780,0.001048) (2.25378,0.000911) (2.26976,0.000792) (2.28574,0.000687) (2.30171,0.000596) (2.31769,0.000517) (2.33367,0.000448) (2.34964,0.000388) (2.36562,0.000335) (2.38160,0.000290) (2.39758,0.000251) (2.41355,0.000216) (2.42953,0.000187) (2.44551,0.000161) (2.46148,0.000139) (2.47746,0.000120) (2.49344,0.000103) (2.50942,0.000089) (2.52539,0.000076) (2.54137,0.000066) (2.55735,0.000056) (2.57332,0.000048) (2.58930,0.000041) (2.60528,0.000036) (2.62126,0.000030) (2.63723,0.000026) (2.65321,0.000022) (2.66919,0.000019) (2.68516,0.000016) (2.70114,0.000014) (2.71712,0.000012) (2.73310,0.000010) (2.74907,0.000009) (2.76505,0.000007) (2.78103,0.000006) (2.79700,0.000005) (2.81298,0.000005) (2.82896,0.000004) (2.84494,0.000003) (2.86091,0.000003) (2.87689,0.000002) (2.89287,0.000002) (2.90884,0.000002) (2.92482,0.000001) (2.94080,0.000001) (2.95678,0.000001) (2.97275,0.000001) (2.98873,0.000001) (3.00471,0.000001) (3.02068,0.000001) (3.03666,0.000000) (3.05264,0.000000) (3.06862,0.000000) (3.08459,0.000000) (3.10057,0.000000) (3.11655,0.000000) (3.13253,0.000000) (3.14850,0.000000) (3.16448,0.000000) (3.18046,0.000000) (3.19643,0.000000) (3.21241,0.000000) (3.22839,0.000000) (3.24437,0.000000) (3.26034,0.000000) (3.27632,0.000000) (3.29230,0.000000) (3.30827,0.000000) (3.32425,0.000000) (3.34023,0.000000) (3.35621,0.000000) (3.37218,0.000000) (3.38816,0.000000) (3.40414,0.000000) (3.42011,0.000000) (3.43609,0.000000) (3.45207,0.000000) (3.46805,0.000000) (3.48402,0.000000) (3.50000,0.000000)};
\end{scope}
\draw[->] (0,0) -- (1.02,0);
\draw[->] (0,0) -- (0,1.02);
\draw (0,0) -- (0,-.012) node[below,font=\small] {$ 0 $};
\draw (0.2857142857,0) -- (0.2857142857,-.012) node[below,font=\small] {$ 1 $};
\draw (0.5714285714,0) -- (0.5714285714,-.012) node[below,font=\small] {$ 2 $};
\draw (0.8571428571,0) -- (0.8571428571,-.012) node[below,font=\small] {$ 3 $};
\draw (0,0) -- (-.012,0) node[left,font=\small] {$ 0 $};
\draw (0,0.2857142857) -- (-.012,0.2857142857) node[left,font=\small] {$ 0.5 $};
\draw (0,0.5714285714) -- (-.012,0.5714285714) node[left,font=\small] {$ 1 $};
\draw (0,0.8571428571) -- (-.012,0.8571428571) node[left,font=\small] {$ 1.5 $};
\node[below,font=\small] at (.5,-.09) {$Q/\nu$};
\node[rotate=90,font=\small] at (-.13,.5) {密度};
\node[anchor=north east,fill=white,inner sep=3pt,font=\footnotesize] at (.98,.98) {%
\begin{tabular}{@{}ll@{}}
\tikz[baseline=-.5ex]{\draw[black,dotted,thick] (0,0) -- (.5,0);} & $\nu=3$\\
\tikz[baseline=-.5ex]{\draw[black,dashed,thick] (0,0) -- (.5,0);} & $\nu=10$\\
\tikz[baseline=-.5ex]{\draw[black,thick] (0,0) -- (.5,0);} & $\nu=30$\\
\end{tabular}};
\end{tikzpicture}\end{center}

若 $Q\sim\chi^2_\nu$，则 $Q/\nu$ 的均值为一，方差为 $2/\nu$。用相同均值尺度比较，可见自由度较小时明显右偏，自由度增大后集中程度提高。正态样本的 $S^2/\sigma^2$ 正具有这种分布，自由度为 $n-1$。

方差区间通常不对称，并非因为计算不精确，而是参考分布偏斜和倒数反演共同造成。对方差区间端点取平方根可得到标准差区间，但不能只对半径取平方根后保留原中心。参见习题~\ref{q:M-C22-1}、\ref{q:M-P10-3}。

\clearpage\subsection*{图解六：置信水平的重复抽样含义}

\begin{center}% 使用共享样式已加载的 TikZ，保留原图数据。
\begin{tikzpicture}[x=.87\linewidth,y=.69\linewidth]
% 直接设置数据坐标单位，避免先形成大尺寸再缩放造成溢出。
\begin{scope}[x=0.021219512193\linewidth,y=0.081176470572\linewidth,shift={(0,4.2)}]
\clip (0,-4.2) rectangle (41,4.3);
\draw[black,thick] plot coordinates {(1.00000,-2.827854) (1.00000,1.092074)};
\draw[black,thick] plot coordinates {(2.00000,-3.473417) (2.00000,0.446511)};
\draw[black,thick] plot coordinates {(3.00000,-1.673673) (3.00000,2.246255)};
\draw[black,thick] plot coordinates {(4.00000,-2.426365) (4.00000,1.493563)};
\draw[black,thick] plot coordinates {(5.00000,-0.682072) (5.00000,3.237856)};
\draw[black,thick] plot coordinates {(6.00000,-1.216195) (6.00000,2.703733)};
\draw[black,thick] plot coordinates {(7.00000,-3.777209) (7.00000,0.142719)};
\draw[black,thick] plot coordinates {(8.00000,-2.377894) (8.00000,1.542034)};
\draw[black,thick] plot coordinates {(9.00000,-1.772638) (9.00000,2.147290)};
\draw[black,thick] plot coordinates {(10.00000,-2.660165) (10.00000,1.259763)};
\draw[black,thick] plot coordinates {(11.00000,-3.422815) (11.00000,0.497113)};
\draw[black,thick] plot coordinates {(12.00000,-1.601206) (12.00000,2.318722)};
\draw[black,thick] plot coordinates {(13.00000,-1.766438) (13.00000,2.153490)};
\draw[black,thick] plot coordinates {(14.00000,-2.707712) (14.00000,1.212216)};
\draw[black,thick] plot coordinates {(15.00000,-3.196906) (15.00000,0.723022)};
\draw[black,thick] plot coordinates {(16.00000,-0.057448) (16.00000,3.862480)};
\draw[black,thick] plot coordinates {(17.00000,-2.927236) (17.00000,0.992692)};
\draw[black,thick] plot coordinates {(18.00000,-2.192602) (18.00000,1.727326)};
\draw[black,thick] plot coordinates {(19.00000,-0.855217) (19.00000,3.064711)};
\draw[black,thick] plot coordinates {(20.00000,-3.760312) (20.00000,0.159616)};
\draw[black,thick] plot coordinates {(21.00000,-0.760053) (21.00000,3.159875)};
\draw[black,thick] plot coordinates {(22.00000,-1.705289) (22.00000,2.214639)};
\draw[black,thick] plot coordinates {(23.00000,-2.861649) (23.00000,1.058278)};
\draw[black,thick] plot coordinates {(24.00000,-1.339421) (24.00000,2.580507)};
\draw[black,thick] plot coordinates {(25.00000,-1.516632) (25.00000,2.403296)};
\draw[black,thick] plot coordinates {(26.00000,-0.821718) (26.00000,3.098210)};
\draw[black,thick] plot coordinates {(27.00000,-2.156287) (27.00000,1.763641)};
\draw[black,thick] plot coordinates {(28.00000,-0.514062) (28.00000,3.405866)};
\draw[black,thick] plot coordinates {(29.00000,-2.637694) (29.00000,1.282234)};
\draw[black,thick] plot coordinates {(30.00000,-0.879699) (30.00000,3.040229)};
\draw[black,thick] plot coordinates {(31.00000,-2.377251) (31.00000,1.542677)};
\draw[black,thick] plot coordinates {(32.00000,-1.464345) (32.00000,2.455583)};
\draw[black,thick] plot coordinates {(33.00000,-2.332668) (33.00000,1.587260)};
\draw[black,thick] plot coordinates {(34.00000,-1.730235) (34.00000,2.189693)};
\draw[black,thick] plot coordinates {(35.00000,-1.699929) (35.00000,2.219999)};
\draw[black,thick] plot coordinates {(36.00000,-0.849247) (36.00000,3.070681)};
\draw[black,thick] plot coordinates {(37.00000,-2.769283) (37.00000,1.150645)};
\draw[black,thick] plot coordinates {(38.00000,-2.431050) (38.00000,1.488878)};
\draw[black,thick] plot coordinates {(39.00000,-0.980838) (39.00000,2.939090)};
\draw[black,thick] plot coordinates {(40.00000,-2.125817) (40.00000,1.794111)};
\draw[gray] plot coordinates {(0.00000,0.000000) (41.00000,0.000000)};
\end{scope}
\draw[->] (0,0) -- (1.02,0);
\draw[->] (0,0) -- (0,1.02);
\draw (0,0) -- (0,-.012) node[below,font=\small] {$ 0 $};
\draw (0.243902439,0) -- (0.243902439,-.012) node[below,font=\small] {$ 10 $};
\draw (0.487804878,0) -- (0.487804878,-.012) node[below,font=\small] {$ 20 $};
\draw (0.7317073171,0) -- (0.7317073171,-.012) node[below,font=\small] {$ 30 $};
\draw (0.9756097561,0) -- (0.9756097561,-.012) node[below,font=\small] {$ 40 $};
\draw (0,0.02352941176) -- (-.012,0.02352941176) node[left,font=\small] {$ -4 $};
\draw (0,0.2588235294) -- (-.012,0.2588235294) node[left,font=\small] {$ -2 $};
\draw (0,0.4941176471) -- (-.012,0.4941176471) node[left,font=\small] {$ 0 $};
\draw (0,0.7294117647) -- (-.012,0.7294117647) node[left,font=\small] {$ 2 $};
\draw (0,0.9647058824) -- (-.012,0.9647058824) node[left,font=\small] {$ 4 $};
\node[below,font=\small] at (.5,-.09) {重复试验编号};
\node[rotate=90,font=\small] at (-.13,.5) {以标准误为单位的均值区间};
\end{tikzpicture}\end{center}

固定真实均值为零，独立重复模拟 40 次已知标准误的正态均值区间。实线覆盖零，粗虚线未覆盖。本次固定随机种子的模拟中有 0 个区间未覆盖。模拟比例无需恰好等于 $5\%$，因为 40 次有限重复本身也有随机波动。

每条竖线是不同样本产生的随机区间，水平零线是同一个固定参数。$95\%$ 描述的是长期构造程序的覆盖率。给定其中一条已经算出的区间，它是否包含零已是确定事实；若要讨论参数给定数据后的概率，需要另建贝叶斯模型。参见习题~\ref{q:M-P10-5}。

\clearpage\subsection*{图解七：样本容量与检验功效}

\begin{center}% 使用共享样式已加载的 TikZ，保留原图数据。
\begin{tikzpicture}[x=.87\linewidth,y=.69\linewidth]
% 直接设置数据坐标单位，避免先形成大尺寸再缩放造成溢出。
\begin{scope}[x=0.435\linewidth,y=0.676470588261\linewidth,shift={(1,0)}]
\clip (-1,0) rectangle (1,1.02);
\draw[black,dotted,thick] plot coordinates {(-1.00000,0.885379) (-0.99000,0.879139) (-0.98000,0.872663) (-0.97000,0.865950) (-0.96000,0.858997) (-0.95000,0.851804) (-0.94000,0.844369) (-0.93000,0.836693) (-0.92000,0.828775) (-0.91000,0.820615) (-0.90000,0.812215) (-0.89000,0.803577) (-0.88000,0.794702) (-0.87000,0.785593) (-0.86000,0.776253) (-0.85000,0.766686) (-0.84000,0.756897) (-0.83000,0.746890) (-0.82000,0.736670) (-0.81000,0.726243) (-0.80000,0.715617) (-0.79000,0.704797) (-0.78000,0.693791) (-0.77000,0.682608) (-0.76000,0.671256) (-0.75000,0.659744) (-0.74000,0.648081) (-0.73000,0.636277) (-0.72000,0.624342) (-0.71000,0.612288) (-0.70000,0.600124) (-0.69000,0.587863) (-0.68000,0.575516) (-0.67000,0.563094) (-0.66000,0.550611) (-0.65000,0.538077) (-0.64000,0.525506) (-0.63000,0.512911) (-0.62000,0.500303) (-0.61000,0.487695) (-0.60000,0.475101) (-0.59000,0.462532) (-0.58000,0.450002) (-0.57000,0.437523) (-0.56000,0.425107) (-0.55000,0.412767) (-0.54000,0.400514) (-0.53000,0.388360) (-0.52000,0.376317) (-0.51000,0.364396) (-0.50000,0.352608) (-0.49000,0.340963) (-0.48000,0.329472) (-0.47000,0.318143) (-0.46000,0.306988) (-0.45000,0.296013) (-0.44000,0.285229) (-0.43000,0.274642) (-0.42000,0.264261) (-0.41000,0.254092) (-0.40000,0.244141) (-0.39000,0.234416) (-0.38000,0.224921) (-0.37000,0.215661) (-0.36000,0.206641) (-0.35000,0.197865) (-0.34000,0.189337) (-0.33000,0.181060) (-0.32000,0.173036) (-0.31000,0.165267) (-0.30000,0.157756) (-0.29000,0.150504) (-0.28000,0.143511) (-0.27000,0.136779) (-0.26000,0.130308) (-0.25000,0.124097) (-0.24000,0.118147) (-0.23000,0.112457) (-0.22000,0.107025) (-0.21000,0.101852) (-0.20000,0.096935) (-0.19000,0.092274) (-0.18000,0.087867) (-0.17000,0.083711) (-0.16000,0.079806) (-0.15000,0.076150) (-0.14000,0.072740) (-0.13000,0.069575) (-0.12000,0.066654) (-0.11000,0.063973) (-0.10000,0.061533) (-0.09000,0.059330) (-0.08000,0.057363) (-0.07000,0.055632) (-0.06000,0.054134) (-0.05000,0.052869) (-0.04000,0.051835) (-0.03000,0.051032) (-0.02000,0.050458) (-0.01000,0.050115) (0.00000,0.050000) (0.01000,0.050115) (0.02000,0.050458) (0.03000,0.051032) (0.04000,0.051835) (0.05000,0.052869) (0.06000,0.054134) (0.07000,0.055632) (0.08000,0.057363) (0.09000,0.059330) (0.10000,0.061533) (0.11000,0.063973) (0.12000,0.066654) (0.13000,0.069575) (0.14000,0.072740) (0.15000,0.076150) (0.16000,0.079806) (0.17000,0.083711) (0.18000,0.087867) (0.19000,0.092274) (0.20000,0.096935) (0.21000,0.101852) (0.22000,0.107025) (0.23000,0.112457) (0.24000,0.118147) (0.25000,0.124097) (0.26000,0.130308) (0.27000,0.136779) (0.28000,0.143511) (0.29000,0.150504) (0.30000,0.157756) (0.31000,0.165267) (0.32000,0.173036) (0.33000,0.181060) (0.34000,0.189337) (0.35000,0.197865) (0.36000,0.206641) (0.37000,0.215661) (0.38000,0.224921) (0.39000,0.234416) (0.40000,0.244141) (0.41000,0.254092) (0.42000,0.264261) (0.43000,0.274642) (0.44000,0.285229) (0.45000,0.296013) (0.46000,0.306988) (0.47000,0.318143) (0.48000,0.329472) (0.49000,0.340963) (0.50000,0.352608) (0.51000,0.364396) (0.52000,0.376317) (0.53000,0.388360) (0.54000,0.400514) (0.55000,0.412767) (0.56000,0.425107) (0.57000,0.437523) (0.58000,0.450002) (0.59000,0.462532) (0.60000,0.475101) (0.61000,0.487695) (0.62000,0.500303) (0.63000,0.512911) (0.64000,0.525506) (0.65000,0.538077) (0.66000,0.550611) (0.67000,0.563094) (0.68000,0.575516) (0.69000,0.587863) (0.70000,0.600124) (0.71000,0.612288) (0.72000,0.624342) (0.73000,0.636277) (0.74000,0.648081) (0.75000,0.659744) (0.76000,0.671256) (0.77000,0.682608) (0.78000,0.693791) (0.79000,0.704797) (0.80000,0.715617) (0.81000,0.726243) (0.82000,0.736670) (0.83000,0.746890) (0.84000,0.756897) (0.85000,0.766686) (0.86000,0.776253) (0.87000,0.785593) (0.88000,0.794702) (0.89000,0.803577) (0.90000,0.812215) (0.91000,0.820615) (0.92000,0.828775) (0.93000,0.836693) (0.94000,0.844369) (0.95000,0.851804) (0.96000,0.858997) (0.97000,0.865950) (0.98000,0.872663) (0.99000,0.879139) (1.00000,0.885379)};
\draw[black,dashed,thick] plot coordinates {(-1.00000,0.999782) (-0.99000,0.999732) (-0.98000,0.999672) (-0.97000,0.999600) (-0.96000,0.999513) (-0.95000,0.999409) (-0.94000,0.999285) (-0.93000,0.999137) (-0.92000,0.998962) (-0.91000,0.998754) (-0.90000,0.998509) (-0.89000,0.998220) (-0.88000,0.997882) (-0.87000,0.997486) (-0.86000,0.997024) (-0.85000,0.996488) (-0.84000,0.995866) (-0.83000,0.995147) (-0.82000,0.994319) (-0.81000,0.993368) (-0.80000,0.992278) (-0.79000,0.991035) (-0.78000,0.989619) (-0.77000,0.988012) (-0.76000,0.986193) (-0.75000,0.984141) (-0.74000,0.981834) (-0.73000,0.979246) (-0.72000,0.976352) (-0.71000,0.973126) (-0.70000,0.969541) (-0.69000,0.965569) (-0.68000,0.961180) (-0.67000,0.956346) (-0.66000,0.951038) (-0.65000,0.945227) (-0.64000,0.938883) (-0.63000,0.931978) (-0.62000,0.924487) (-0.61000,0.916383) (-0.60000,0.907642) (-0.59000,0.898242) (-0.58000,0.888165) (-0.57000,0.877394) (-0.56000,0.865914) (-0.55000,0.853718) (-0.54000,0.840797) (-0.53000,0.827151) (-0.52000,0.812782) (-0.51000,0.797697) (-0.50000,0.781908) (-0.49000,0.765431) (-0.48000,0.748288) (-0.47000,0.730505) (-0.46000,0.712114) (-0.45000,0.693151) (-0.44000,0.673656) (-0.43000,0.653676) (-0.42000,0.633259) (-0.41000,0.612458) (-0.40000,0.591331) (-0.39000,0.569935) (-0.38000,0.548332) (-0.37000,0.526587) (-0.36000,0.504764) (-0.35000,0.482930) (-0.34000,0.461148) (-0.33000,0.439487) (-0.32000,0.418009) (-0.31000,0.396779) (-0.30000,0.375856) (-0.29000,0.355301) (-0.28000,0.335168) (-0.27000,0.315510) (-0.26000,0.296376) (-0.25000,0.277810) (-0.24000,0.259854) (-0.23000,0.242542) (-0.22000,0.225908) (-0.21000,0.209978) (-0.20000,0.194775) (-0.19000,0.180319) (-0.18000,0.166623) (-0.17000,0.153698) (-0.16000,0.141551) (-0.15000,0.130186) (-0.14000,0.119603) (-0.13000,0.109799) (-0.12000,0.100771) (-0.11000,0.092513) (-0.10000,0.085016) (-0.09000,0.078271) (-0.08000,0.072271) (-0.07000,0.067004) (-0.06000,0.062462) (-0.05000,0.058635) (-0.04000,0.055517) (-0.03000,0.053099) (-0.02000,0.051376) (-0.01000,0.050344) (0.00000,0.050000) (0.01000,0.050344) (0.02000,0.051376) (0.03000,0.053099) (0.04000,0.055517) (0.05000,0.058635) (0.06000,0.062462) (0.07000,0.067004) (0.08000,0.072271) (0.09000,0.078271) (0.10000,0.085016) (0.11000,0.092513) (0.12000,0.100771) (0.13000,0.109799) (0.14000,0.119603) (0.15000,0.130186) (0.16000,0.141551) (0.17000,0.153698) (0.18000,0.166623) (0.19000,0.180319) (0.20000,0.194775) (0.21000,0.209978) (0.22000,0.225908) (0.23000,0.242542) (0.24000,0.259854) (0.25000,0.277810) (0.26000,0.296376) (0.27000,0.315510) (0.28000,0.335168) (0.29000,0.355301) (0.30000,0.375856) (0.31000,0.396779) (0.32000,0.418009) (0.33000,0.439487) (0.34000,0.461148) (0.35000,0.482930) (0.36000,0.504764) (0.37000,0.526587) (0.38000,0.548332) (0.39000,0.569935) (0.40000,0.591331) (0.41000,0.612458) (0.42000,0.633259) (0.43000,0.653676) (0.44000,0.673656) (0.45000,0.693151) (0.46000,0.712114) (0.47000,0.730505) (0.48000,0.748288) (0.49000,0.765431) (0.50000,0.781908) (0.51000,0.797697) (0.52000,0.812782) (0.53000,0.827151) (0.54000,0.840797) (0.55000,0.853718) (0.56000,0.865914) (0.57000,0.877394) (0.58000,0.888165) (0.59000,0.898242) (0.60000,0.907642) (0.61000,0.916383) (0.62000,0.924487) (0.63000,0.931978) (0.64000,0.938883) (0.65000,0.945227) (0.66000,0.951038) (0.67000,0.956346) (0.68000,0.961180) (0.69000,0.965569) (0.70000,0.969541) (0.71000,0.973126) (0.72000,0.976352) (0.73000,0.979246) (0.74000,0.981834) (0.75000,0.984141) (0.76000,0.986193) (0.77000,0.988012) (0.78000,0.989619) (0.79000,0.991035) (0.80000,0.992278) (0.81000,0.993368) (0.82000,0.994319) (0.83000,0.995147) (0.84000,0.995866) (0.85000,0.996488) (0.86000,0.997024) (0.87000,0.997486) (0.88000,0.997882) (0.89000,0.998220) (0.90000,0.998509) (0.91000,0.998754) (0.92000,0.998962) (0.93000,0.999137) (0.94000,0.999285) (0.95000,0.999409) (0.96000,0.999513) (0.97000,0.999600) (0.98000,0.999672) (0.99000,0.999732) (1.00000,0.999782)};
\draw[black,thick] plot coordinates {(-1.00000,1.000000) (-0.99000,1.000000) (-0.98000,1.000000) (-0.97000,1.000000) (-0.96000,1.000000) (-0.95000,1.000000) (-0.94000,1.000000) (-0.93000,1.000000) (-0.92000,1.000000) (-0.91000,1.000000) (-0.90000,1.000000) (-0.89000,1.000000) (-0.88000,1.000000) (-0.87000,1.000000) (-0.86000,1.000000) (-0.85000,1.000000) (-0.84000,1.000000) (-0.83000,1.000000) (-0.82000,1.000000) (-0.81000,1.000000) (-0.80000,1.000000) (-0.79000,1.000000) (-0.78000,1.000000) (-0.77000,1.000000) (-0.76000,1.000000) (-0.75000,1.000000) (-0.74000,1.000000) (-0.73000,1.000000) (-0.72000,1.000000) (-0.71000,1.000000) (-0.70000,1.000000) (-0.69000,1.000000) (-0.68000,0.999999) (-0.67000,0.999999) (-0.66000,0.999998) (-0.65000,0.999997) (-0.64000,0.999996) (-0.63000,0.999993) (-0.62000,0.999989) (-0.61000,0.999983) (-0.60000,0.999973) (-0.59000,0.999959) (-0.58000,0.999938) (-0.57000,0.999908) (-0.56000,0.999864) (-0.55000,0.999800) (-0.54000,0.999709) (-0.53000,0.999581) (-0.52000,0.999402) (-0.51000,0.999155) (-0.50000,0.998817) (-0.49000,0.998359) (-0.48000,0.997745) (-0.47000,0.996928) (-0.46000,0.995855) (-0.45000,0.994458) (-0.44000,0.992657) (-0.43000,0.990359) (-0.42000,0.987456) (-0.41000,0.983824) (-0.40000,0.979327) (-0.39000,0.973812) (-0.38000,0.967119) (-0.37000,0.959074) (-0.36000,0.949501) (-0.35000,0.938224) (-0.34000,0.925071) (-0.33000,0.909883) (-0.32000,0.892519) (-0.31000,0.872865) (-0.30000,0.850839) (-0.29000,0.826401) (-0.28000,0.799557) (-0.27000,0.770363) (-0.26000,0.738928) (-0.25000,0.705418) (-0.24000,0.670051) (-0.23000,0.633096) (-0.22000,0.594865) (-0.21000,0.555709) (-0.20000,0.516005) (-0.19000,0.476149) (-0.18000,0.436540) (-0.17000,0.397572) (-0.16000,0.359622) (-0.15000,0.323041) (-0.14000,0.288142) (-0.13000,0.255196) (-0.12000,0.224427) (-0.11000,0.196011) (-0.10000,0.170075) (-0.09000,0.146699) (-0.08000,0.125922) (-0.07000,0.107749) (-0.06000,0.092155) (-0.05000,0.079098) (-0.04000,0.068523) (-0.03000,0.060373) (-0.02000,0.054595) (-0.01000,0.051146) (0.00000,0.050000) (0.01000,0.051146) (0.02000,0.054595) (0.03000,0.060373) (0.04000,0.068523) (0.05000,0.079098) (0.06000,0.092155) (0.07000,0.107749) (0.08000,0.125922) (0.09000,0.146699) (0.10000,0.170075) (0.11000,0.196011) (0.12000,0.224427) (0.13000,0.255196) (0.14000,0.288142) (0.15000,0.323041) (0.16000,0.359622) (0.17000,0.397572) (0.18000,0.436540) (0.19000,0.476149) (0.20000,0.516005) (0.21000,0.555709) (0.22000,0.594865) (0.23000,0.633096) (0.24000,0.670051) (0.25000,0.705418) (0.26000,0.738928) (0.27000,0.770363) (0.28000,0.799557) (0.29000,0.826401) (0.30000,0.850839) (0.31000,0.872865) (0.32000,0.892519) (0.33000,0.909883) (0.34000,0.925071) (0.35000,0.938224) (0.36000,0.949501) (0.37000,0.959074) (0.38000,0.967119) (0.39000,0.973812) (0.40000,0.979327) (0.41000,0.983824) (0.42000,0.987456) (0.43000,0.990359) (0.44000,0.992657) (0.45000,0.994458) (0.46000,0.995855) (0.47000,0.996928) (0.48000,0.997745) (0.49000,0.998359) (0.50000,0.998817) (0.51000,0.999155) (0.52000,0.999402) (0.53000,0.999581) (0.54000,0.999709) (0.55000,0.999800) (0.56000,0.999864) (0.57000,0.999908) (0.58000,0.999938) (0.59000,0.999959) (0.60000,0.999973) (0.61000,0.999983) (0.62000,0.999989) (0.63000,0.999993) (0.64000,0.999996) (0.65000,0.999997) (0.66000,0.999998) (0.67000,0.999999) (0.68000,0.999999) (0.69000,1.000000) (0.70000,1.000000) (0.71000,1.000000) (0.72000,1.000000) (0.73000,1.000000) (0.74000,1.000000) (0.75000,1.000000) (0.76000,1.000000) (0.77000,1.000000) (0.78000,1.000000) (0.79000,1.000000) (0.80000,1.000000) (0.81000,1.000000) (0.82000,1.000000) (0.83000,1.000000) (0.84000,1.000000) (0.85000,1.000000) (0.86000,1.000000) (0.87000,1.000000) (0.88000,1.000000) (0.89000,1.000000) (0.90000,1.000000) (0.91000,1.000000) (0.92000,1.000000) (0.93000,1.000000) (0.94000,1.000000) (0.95000,1.000000) (0.96000,1.000000) (0.97000,1.000000) (0.98000,1.000000) (0.99000,1.000000) (1.00000,1.000000)};
\end{scope}
\draw[->] (0,0) -- (1.02,0);
\draw[->] (0,0) -- (0,1.02);
\draw (0,0) -- (0,-.012) node[below,font=\small] {$ -1 $};
\draw (0.25,0) -- (0.25,-.012) node[below,font=\small] {$ -0.5 $};
\draw (0.5,0) -- (0.5,-.012) node[below,font=\small] {$ 0 $};
\draw (0.75,0) -- (0.75,-.012) node[below,font=\small] {$ 0.5 $};
\draw (1,0) -- (1,-.012) node[below,font=\small] {$ 1 $};
\draw (0,0) -- (-.012,0) node[left,font=\small] {$ 0 $};
\draw (0,0.1960784314) -- (-.012,0.1960784314) node[left,font=\small] {$ 0.2 $};
\draw (0,0.3921568627) -- (-.012,0.3921568627) node[left,font=\small] {$ 0.4 $};
\draw (0,0.5882352941) -- (-.012,0.5882352941) node[left,font=\small] {$ 0.6 $};
\draw (0,0.7843137255) -- (-.012,0.7843137255) node[left,font=\small] {$ 0.8 $};
\draw (0,0.9803921569) -- (-.012,0.9803921569) node[left,font=\small] {$ 1 $};
\node[below,font=\small] at (.5,-.09) {真实均值 $\mu$};
\node[rotate=90,font=\small] at (-.13,.5) {拒绝概率};
\node[anchor=south east,fill=white,inner sep=3pt,font=\footnotesize] at (.98,.03) {%
\begin{tabular}{@{}ll@{}}
\tikz[baseline=-.5ex]{\draw[black,dotted,thick] (0,0) -- (.5,0);} & $n=10$\\
\tikz[baseline=-.5ex]{\draw[black,dashed,thick] (0,0) -- (.5,0);} & $n=30$\\
\tikz[baseline=-.5ex]{\draw[black,thick] (0,0) -- (.5,0);} & $n=100$\\
\end{tabular}};
\end{tikzpicture}\end{center}

检验 $H_0:\mu=0$，独立正态观测方差为一，双侧水平为 $0.05$。在 $\mu=0$，三条曲线都等于第一类错误概率 $0.05$；离开零假设后，容量较大的检验更快获得高功效。

功效取决于真实差异、噪声、容量和检验规则。没有显著结果可能只是样本精度不足；反过来，大样本可以使很小的差异达到显著，所以实际意义仍应结合效应大小判断。图中的双侧曲线关于零对称，是该正态模型的对称性所致。参见习题~\ref{q:B14.1}、\ref{q:M-F-11}。

\clearpage\subsection*{图解八：拟合直线与两种不确定性带}

\begin{center}% 使用共享样式已加载的 TikZ，保留原图数据。
\begin{tikzpicture}[x=.87\linewidth,y=.69\linewidth]
% 直接设置数据坐标单位，避免先形成大尺寸再缩放造成溢出。
\begin{scope}[x=0.131818181805\linewidth,y=0.084146341455\linewidth,shift={(2.8,2.8)}]
\clip (-2.8,-2.8) rectangle (3.8,5.4);
\draw[only marks,mark=*,black] plot coordinates {(-2.00000,-0.500000) (-1.00000,-0.200000) (0.00000,1.200000) (1.00000,1.500000) (2.00000,2.900000) (3.00000,2.800000)};
\draw[black,thick] plot coordinates {(-2.50000,-0.953810) (-2.43939,-0.908615) (-2.37879,-0.863420) (-2.31818,-0.818225) (-2.25758,-0.773030) (-2.19697,-0.727835) (-2.13636,-0.682641) (-2.07576,-0.637446) (-2.01515,-0.592251) (-1.95455,-0.547056) (-1.89394,-0.501861) (-1.83333,-0.456667) (-1.77273,-0.411472) (-1.71212,-0.366277) (-1.65152,-0.321082) (-1.59091,-0.275887) (-1.53030,-0.230693) (-1.46970,-0.185498) (-1.40909,-0.140303) (-1.34848,-0.095108) (-1.28788,-0.049913) (-1.22727,-0.004719) (-1.16667,0.040476) (-1.10606,0.085671) (-1.04545,0.130866) (-0.98485,0.176061) (-0.92424,0.221255) (-0.86364,0.266450) (-0.80303,0.311645) (-0.74242,0.356840) (-0.68182,0.402035) (-0.62121,0.447229) (-0.56061,0.492424) (-0.50000,0.537619) (-0.43939,0.582814) (-0.37879,0.628009) (-0.31818,0.673203) (-0.25758,0.718398) (-0.19697,0.763593) (-0.13636,0.808788) (-0.07576,0.853983) (-0.01515,0.899177) (0.04545,0.944372) (0.10606,0.989567) (0.16667,1.034762) (0.22727,1.079957) (0.28788,1.125152) (0.34848,1.170346) (0.40909,1.215541) (0.46970,1.260736) (0.53030,1.305931) (0.59091,1.351126) (0.65152,1.396320) (0.71212,1.441515) (0.77273,1.486710) (0.83333,1.531905) (0.89394,1.577100) (0.95455,1.622294) (1.01515,1.667489) (1.07576,1.712684) (1.13636,1.757879) (1.19697,1.803074) (1.25758,1.848268) (1.31818,1.893463) (1.37879,1.938658) (1.43939,1.983853) (1.50000,2.029048) (1.56061,2.074242) (1.62121,2.119437) (1.68182,2.164632) (1.74242,2.209827) (1.80303,2.255022) (1.86364,2.300216) (1.92424,2.345411) (1.98485,2.390606) (2.04545,2.435801) (2.10606,2.480996) (2.16667,2.526190) (2.22727,2.571385) (2.28788,2.616580) (2.34848,2.661775) (2.40909,2.706970) (2.46970,2.752165) (2.53030,2.797359) (2.59091,2.842554) (2.65152,2.887749) (2.71212,2.932944) (2.77273,2.978139) (2.83333,3.023333) (2.89394,3.068528) (2.95455,3.113723) (3.01515,3.158918) (3.07576,3.204113) (3.13636,3.249307) (3.19697,3.294502) (3.25758,3.339697) (3.31818,3.384892) (3.37879,3.430087) (3.43939,3.475281) (3.50000,3.520476)};
\draw[black,dashed] plot coordinates {(-2.50000,-1.853465) (-2.43939,-1.794578) (-2.37879,-1.735762) (-2.31818,-1.677021) (-2.25758,-1.618358) (-2.19697,-1.559777) (-2.13636,-1.501281) (-2.07576,-1.442876) (-2.01515,-1.384565) (-1.95455,-1.326353) (-1.89394,-1.268246) (-1.83333,-1.210249) (-1.77273,-1.152368) (-1.71212,-1.094607) (-1.65152,-1.036975) (-1.59091,-0.979478) (-1.53030,-0.922122) (-1.46970,-0.864916) (-1.40909,-0.807868) (-1.34848,-0.750986) (-1.28788,-0.694279) (-1.22727,-0.637757) (-1.16667,-0.581429) (-1.10606,-0.525308) (-1.04545,-0.469403) (-0.98485,-0.413726) (-0.92424,-0.358290) (-0.86364,-0.303109) (-0.80303,-0.248195) (-0.74242,-0.193562) (-0.68182,-0.139226) (-0.62121,-0.085202) (-0.56061,-0.031505) (-0.50000,0.021848) (-0.43939,0.074841) (-0.37879,0.127457) (-0.31818,0.179678) (-0.25758,0.231489) (-0.19697,0.282871) (-0.13636,0.333809) (-0.07576,0.384286) (-0.01515,0.434285) (0.04545,0.483793) (0.10606,0.532795) (0.16667,0.581279) (0.22727,0.629233) (0.28788,0.676647) (0.34848,0.723514) (0.40909,0.769827) (0.46970,0.815582) (0.53030,0.860776) (0.59091,0.905411) (0.65152,0.949488) (0.71212,0.993011) (0.77273,1.035986) (0.83333,1.078422) (0.89394,1.120328) (0.95455,1.161715) (1.01515,1.202597) (1.07576,1.242987) (1.13636,1.282900) (1.19697,1.322352) (1.25758,1.361359) (1.31818,1.399938) (1.37879,1.438106) (1.43939,1.475880) (1.50000,1.513276) (1.56061,1.550313) (1.62121,1.587006) (1.68182,1.623371) (1.74242,1.659425) (1.80303,1.695182) (1.86364,1.730657) (1.92424,1.765865) (1.98485,1.800819) (2.04545,1.835532) (2.10606,1.870017) (2.16667,1.904285) (2.22727,1.938347) (2.28788,1.972215) (2.34848,2.005897) (2.40909,2.039405) (2.46970,2.072746) (2.53030,2.105930) (2.59091,2.138964) (2.65152,2.171856) (2.71212,2.204613) (2.77273,2.237243) (2.83333,2.269751) (2.89394,2.302143) (2.95455,2.334426) (3.01515,2.366604) (3.07576,2.398683) (3.13636,2.430667) (3.19697,2.462561) (3.25758,2.494369) (3.31818,2.526096) (3.37879,2.557744) (3.43939,2.589318) (3.50000,2.620821)};
\draw[black,dashed] plot coordinates {(-2.50000,-0.054154) (-2.43939,-0.022651) (-2.37879,0.008923) (-2.31818,0.040571) (-2.25758,0.072297) (-2.19697,0.104106) (-2.13636,0.136000) (-2.07576,0.167984) (-2.01515,0.200063) (-1.95455,0.232241) (-1.89394,0.264523) (-1.83333,0.296916) (-1.77273,0.329424) (-1.71212,0.362053) (-1.65152,0.394811) (-1.59091,0.427703) (-1.53030,0.460737) (-1.46970,0.493921) (-1.40909,0.527262) (-1.34848,0.560770) (-1.28788,0.594452) (-1.22727,0.628320) (-1.16667,0.662382) (-1.10606,0.696650) (-1.04545,0.731134) (-0.98485,0.765847) (-0.92424,0.800801) (-0.86364,0.836009) (-0.80303,0.871485) (-0.74242,0.907242) (-0.68182,0.943296) (-0.62121,0.979661) (-0.56061,1.016354) (-0.50000,1.053390) (-0.43939,1.090787) (-0.37879,1.128561) (-0.31818,1.166729) (-0.25758,1.205308) (-0.19697,1.244315) (-0.13636,1.283767) (-0.07576,1.323680) (-0.01515,1.364070) (0.04545,1.404951) (0.10606,1.446339) (0.16667,1.488245) (0.22727,1.530680) (0.28788,1.573656) (0.34848,1.617179) (0.40909,1.661256) (0.46970,1.705890) (0.53030,1.751085) (0.59091,1.796840) (0.65152,1.843153) (0.71212,1.890019) (0.77273,1.937434) (0.83333,1.985388) (0.89394,2.033871) (0.95455,2.082873) (1.01515,2.132381) (1.07576,2.182381) (1.13636,2.232858) (1.19697,2.283795) (1.25758,2.335178) (1.31818,2.386988) (1.37879,2.439210) (1.43939,2.491826) (1.50000,2.544819) (1.56061,2.598172) (1.62121,2.651869) (1.68182,2.705893) (1.74242,2.760229) (1.80303,2.814861) (1.86364,2.869776) (1.92424,2.924957) (1.98485,2.980393) (2.04545,3.036069) (2.10606,3.091974) (2.16667,3.148096) (2.22727,3.204423) (2.28788,3.260946) (2.34848,3.317653) (2.40909,3.374535) (2.46970,3.431583) (2.53030,3.488789) (2.59091,3.546145) (2.65152,3.603642) (2.71212,3.661274) (2.77273,3.719034) (2.83333,3.776916) (2.89394,3.834913) (2.95455,3.893020) (3.01515,3.951231) (3.07576,4.009542) (3.13636,4.067948) (3.19697,4.126443) (3.25758,4.185025) (3.31818,4.243688) (3.37879,4.302429) (3.43939,4.361245) (3.50000,4.420132)};
\draw[gray,dotted,thick] plot coordinates {(-2.50000,-2.367309) (-2.43939,-2.313439) (-2.37879,-2.259694) (-2.31818,-2.206076) (-2.25758,-2.152588) (-2.19697,-2.099231) (-2.13636,-2.046009) (-2.07576,-1.992923) (-2.01515,-1.939975) (-1.95455,-1.887170) (-1.89394,-1.834508) (-1.83333,-1.781992) (-1.77273,-1.729625) (-1.71212,-1.677408) (-1.65152,-1.625346) (-1.59091,-1.573439) (-1.53030,-1.521690) (-1.46970,-1.470103) (-1.40909,-1.418678) (-1.34848,-1.367419) (-1.28788,-1.316328) (-1.22727,-1.265408) (-1.16667,-1.214660) (-1.10606,-1.164087) (-1.04545,-1.113691) (-0.98485,-1.063475) (-0.92424,-1.013440) (-0.86364,-0.963589) (-0.80303,-0.913924) (-0.74242,-0.864447) (-0.68182,-0.815160) (-0.62121,-0.766065) (-0.56061,-0.717163) (-0.50000,-0.668457) (-0.43939,-0.619948) (-0.37879,-0.571638) (-0.31818,-0.523528) (-0.25758,-0.475620) (-0.19697,-0.427916) (-0.13636,-0.380416) (-0.07576,-0.333121) (-0.01515,-0.286033) (0.04545,-0.239153) (0.10606,-0.192482) (0.16667,-0.146020) (0.22727,-0.099769) (0.28788,-0.053728) (0.34848,-0.007898) (0.40909,0.037721) (0.46970,0.083127) (0.53030,0.128322) (0.59091,0.173305) (0.65152,0.218076) (0.71212,0.262636) (0.77273,0.306985) (0.83333,0.351123) (0.89394,0.395050) (0.95455,0.438769) (1.01515,0.482278) (1.07576,0.525580) (1.13636,0.568675) (1.19697,0.611565) (1.25758,0.654250) (1.31818,0.696731) (1.37879,0.739011) (1.43939,0.781091) (1.50000,0.822971) (1.56061,0.864655) (1.62121,0.906143) (1.68182,0.947437) (1.74242,0.988540) (1.80303,1.029452) (1.86364,1.070177) (1.92424,1.110716) (1.98485,1.151071) (2.04545,1.191244) (2.10606,1.231238) (2.16667,1.271054) (2.22727,1.310696) (2.28788,1.350165) (2.34848,1.389464) (2.40909,1.428594) (2.46970,1.467560) (2.53030,1.506362) (2.59091,1.545003) (2.65152,1.583486) (2.71212,1.621812) (2.77273,1.659986) (2.83333,1.698008) (2.89394,1.735882) (2.95455,1.773610) (3.01515,1.811193) (3.07576,1.848636) (3.13636,1.885939) (3.19697,1.923107) (3.25758,1.960139) (3.31818,1.997041) (3.37879,2.033812) (3.43939,2.070457) (3.50000,2.106977)};
\draw[gray,dotted,thick] plot coordinates {(-2.50000,0.459690) (-2.43939,0.496210) (-2.37879,0.532854) (-2.31818,0.569626) (-2.25758,0.606527) (-2.19697,0.643560) (-2.13636,0.680727) (-2.07576,0.718031) (-2.01515,0.755473) (-1.95455,0.793057) (-1.89394,0.830785) (-1.83333,0.868658) (-1.77273,0.906681) (-1.71212,0.944854) (-1.65152,0.983181) (-1.59091,1.021664) (-1.53030,1.060305) (-1.46970,1.099107) (-1.40909,1.138072) (-1.34848,1.177203) (-1.28788,1.216502) (-1.22727,1.255971) (-1.16667,1.295612) (-1.10606,1.335429) (-1.04545,1.375423) (-0.98485,1.415596) (-0.92424,1.455951) (-0.86364,1.496490) (-0.80303,1.537215) (-0.74242,1.578127) (-0.68182,1.619230) (-0.62121,1.660524) (-0.56061,1.702012) (-0.50000,1.743695) (-0.43939,1.785576) (-0.37879,1.827655) (-0.31818,1.869935) (-0.25758,1.912417) (-0.19697,1.955102) (-0.13636,1.997991) (-0.07576,2.041086) (-0.01515,2.084388) (0.04545,2.127898) (0.10606,2.171616) (0.16667,2.215544) (0.22727,2.259682) (0.28788,2.304031) (0.34848,2.348590) (0.40909,2.393362) (0.46970,2.438345) (0.53030,2.483539) (0.59091,2.528946) (0.65152,2.574564) (0.71212,2.620394) (0.77273,2.666435) (0.83333,2.712687) (0.89394,2.759149) (0.95455,2.805820) (1.01515,2.852700) (1.07576,2.899788) (1.13636,2.947082) (1.19697,2.994582) (1.25758,3.042287) (1.31818,3.090195) (1.37879,3.138305) (1.43939,3.186615) (1.50000,3.235124) (1.56061,3.283830) (1.62121,3.332732) (1.68182,3.381827) (1.74242,3.431114) (1.80303,3.480591) (1.86364,3.530256) (1.92424,3.580107) (1.98485,3.630142) (2.04545,3.680358) (2.10606,3.730754) (2.16667,3.781327) (2.22727,3.832075) (2.28788,3.882995) (2.34848,3.934086) (2.40909,3.985345) (2.46970,4.036769) (2.53030,4.088357) (2.59091,4.140105) (2.65152,4.192012) (2.71212,4.244075) (2.77273,4.296291) (2.83333,4.348658) (2.89394,4.401174) (2.95455,4.453836) (3.01515,4.506642) (3.07576,4.559589) (3.13636,4.612675) (3.19697,4.665898) (3.25758,4.719254) (3.31818,4.772743) (3.37879,4.826361) (3.43939,4.880106) (3.50000,4.933975)};
\end{scope}
\draw[->] (0,0) -- (1.02,0);
\draw[->] (0,0) -- (0,1.02);
\draw (0.1212121212,0) -- (0.1212121212,-.012) node[below,font=\small] {$ -2 $};
\draw (0.2727272727,0) -- (0.2727272727,-.012) node[below,font=\small] {$ -1 $};
\draw (0.4242424242,0) -- (0.4242424242,-.012) node[below,font=\small] {$ 0 $};
\draw (0.5757575758,0) -- (0.5757575758,-.012) node[below,font=\small] {$ 1 $};
\draw (0.7272727273,0) -- (0.7272727273,-.012) node[below,font=\small] {$ 2 $};
\draw (0.8787878788,0) -- (0.8787878788,-.012) node[below,font=\small] {$ 3 $};
\draw (0,0.09756097561) -- (-.012,0.09756097561) node[left,font=\small] {$ -2 $};
\draw (0,0.3414634146) -- (-.012,0.3414634146) node[left,font=\small] {$ 0 $};
\draw (0,0.5853658537) -- (-.012,0.5853658537) node[left,font=\small] {$ 2 $};
\draw (0,0.8292682927) -- (-.012,0.8292682927) node[left,font=\small] {$ 4 $};
\node[below,font=\small] at (.5,-.09) {$x$};
\node[rotate=90,font=\small] at (-.13,.5) {$y$};
\node[anchor=north west,fill=white,inner sep=3pt,font=\footnotesize] at (.03,.98) {%
\begin{tabular}{@{}ll@{}}
\tikz[baseline=-.5ex]{\fill (0.25,0) circle[radius=1.5pt];} & 观测\\
\tikz[baseline=-.5ex]{\draw[black,thick] (0,0) -- (.5,0);} & 拟合均值\\
\tikz[baseline=-.5ex]{\draw[black,dashed] (0,0) -- (.5,0);} & 均值区间边界\\
\tikz[baseline=-.5ex]{\draw[gray,dotted,thick] (0,0) -- (.5,0);} & 预测区间边界\\
\end{tabular}};
\end{tikzpicture}\end{center}

图中使用六个编订示例点，在独立正态同方差模型下绘出逐点 $95\%$ 均值响应区间和新个体预测区间。点线预测带较宽，因为新个体还有自身误差；离样本横坐标中心越远，两类区间都变宽。

这里是每个固定 $x$ 各自的逐点区间，不是整条曲线同时覆盖率为 $95\%$ 的置信带。若同时查看许多位置，不能把逐点覆盖解释为整体覆盖。参见习题~\ref{q:B16.4}。

\clearpage
\section*{题源、许可与编订记录}
\addcontentsline{toc}{section}{题源、许可与编订记录}
\subsection*{中文公开题源}
\textbf{Y：闫一博（yyb100001），《何书元〈概率引论〉答案》。}本册选取其公开 \texttt{main.tex} 中的数学问题与解答，保留原题号作为定位键，重排题面与解答并校订。项目署名为“Solutions-of-Introduction-to-Probability © 2024 by yyb”，公开声明采用 CC BY 4.0。其问题背景来源于何书元《概率引论》；本册未冒称这些教材题目为公开解答作者或编订者的原创。项目链接：\href{https://github.com/yyb100001/Solutions-of-Introduction-to-Probability}{公开仓库及许可说明}。许可链接：\href{https://creativecommons.org/licenses/by/4.0/}{知识共享署名 4.0 国际许可}。

\textbf{A：ayhe123，《概率论讲义及习题解答》。}本册选取公开 \texttt{lecturenote.tex} 的习题与证明，保留原题号，删除依赖未收入讲义的外部交叉引用，并补齐相关推导。项目公开说明其笔记与仿真采用 CC BY 4.0。该讲义基于何书元《概率引论》及相应课程。项目链接：\href{https://github.com/ayhe123/probabilityTheory-lecturenote}{公开仓库及许可说明}；许可同上。中文题源均为学习者编写的公开解答，不把它们等同于教材官方答案。

\subsection*{MIT 开放课程题源}
\textbf{M：Jeremy Orloff、Jennifer French Kamrin，MIT OpenCourseWare 18.05，2022 年春季。}课程原名为 \textit{Introduction to Probability and Statistics}。本册将所选题目与答案翻译、重排为中文，补充解释并校订数值与条件；删除了考试分数、交卷说明、软件选择题的非必要细节及依赖未收录图形的题目。涉及原资料的情境或数据，仅作为该课教学材料，不表述为当前现实统计事实。

课程主页：\href{https://ocw.mit.edu/courses/18-05-introduction-to-probability-and-statistics-spring-2022/}{MIT OCW 18.05}。本册所用课程材料在当前使用条款下采用 \href{https://creativecommons.org/licenses/by-nc-sa/4.0/}{CC BY-NC-SA 4.0（署名—非商业性使用—相同方式共享）}；\href{https://ocw.mit.edu/pages/privacy-and-terms-of-use/}{原站使用条款}提供许可和署名要求。本册不是 MIT 的官方译本，MIT 与原作者不对本册校订和新增解答提供背书。

\textbf{B：编订补充。}为衔接本科数理统计课程而新增的标准模型题、证明题和综合题，均在逐题来源处标为“编订补充”。标注“编订延伸”的多问问题，其原始部分来自列明的课程题，新增部分由编订者给出。方法说明、图解和数值表也是本册的新增整理，不声称在原资料中原样出现。

本训练编 MIT 材料的中文改编与原训练编新增编订内容按 CC BY-NC-SA 4.0 提供；中文公开项目的原有 CC BY 4.0 署名和许可信息保留。各部分原有权利归其相应权利人。再使用时应保留来源、许可链接与修改说明，不据本册新增编排改变原材料自身的许可。检索与版本核对日期：2026 年 10 月 9 日。

\subsection*{课程文件定位表}
下列缩写与逐题索引配合使用。标题中的“题”对应原文件的题号；本册重排题号只用于内部学习顺序。


\noindent\textbf{C10}：\href{https://ocw.mit.edu/courses/18-05-introduction-to-probability-and-statistics-spring-2022/mit18_05_s22_class10_pset_sol.pdf}{第 10 讲课堂题及解答}。\par

\noindent\textbf{C22}：\href{https://ocw.mit.edu/courses/18-05-introduction-to-probability-and-statistics-spring-2022/mit18_05_s22_class22_pset_sol.pdf}{第 22 讲课堂题及解答}。\par

\noindent\textbf{C26}：\href{https://ocw.mit.edu/courses/18-05-introduction-to-probability-and-statistics-spring-2022/mit18_05_s22_class26_pset_sol.pdf}{第 26 讲课堂题及解答}。\par

\noindent\textbf{P7}：\href{https://ocw.mit.edu/courses/18-05-introduction-to-probability-and-statistics-spring-2022/mit18_05_s22_pset07_sol.pdf}{习题集 7及解答}。\par

\noindent\textbf{P8}：\href{https://ocw.mit.edu/courses/18-05-introduction-to-probability-and-statistics-spring-2022/mit18_05_s22_pset08_sol.pdf}{习题集 8及解答}。\par

\noindent\textbf{P9}：\href{https://ocw.mit.edu/courses/18-05-introduction-to-probability-and-statistics-spring-2022/mit18_05_s22_pset09_sol.pdf}{习题集 9及解答}。\par

\noindent\textbf{P10}：\href{https://ocw.mit.edu/courses/18-05-introduction-to-probability-and-statistics-spring-2022/mit18_05_s22_pset10_sol.pdf}{习题集 10及解答}。\par

\noindent\textbf{P11}：\href{https://ocw.mit.edu/courses/18-05-introduction-to-probability-and-statistics-spring-2022/mit18_05_s22_pset11_sol.pdf}{习题集 11及解答}。\par

\noindent\textbf{E2}：\href{https://ocw.mit.edu/courses/18-05-introduction-to-probability-and-statistics-spring-2022/mit18_05_s22_exam02_sol.pdf}{第二次考试及解答}。\par

\noindent\textbf{F}：\href{https://ocw.mit.edu/courses/18-05-introduction-to-probability-and-statistics-spring-2022/mit18_05_s22_exam_final_sol.pdf}{期末考试及解答}。\par

\noindent\textbf{PF}：\href{https://ocw.mit.edu/courses/18-05-introduction-to-probability-and-statistics-spring-2022/mit18_05_s22_prac_exam_final_sol.pdf}{模拟期末考试及解答}。\par

\subsection*{主要勘误与条件补全}
以下是影响结论或推导的代表性修正。索引中的原题号用于返回公开题源核对；本册的译写和校订责任由编订承担。
\begin{enumerate}
\item Y1.9：恰好一副手套后，另外两只应选自两副；修正组合数的文字解释。
\item Y1.13：圆心须远离格线至少半径；补充方格边长不超过硬币直径的情形。
\item Y2.23、Y2.25：补明确给定真实状态后的条件独立，以未舍入先验直接进行多次更新。
\item Y3.6：计算 $P(X\ge3)$ 的补事件必须包含零次成功项。
\item Y3.20、Y3.26、Y3.28、Y3.29、Y4.25、Y4.36 等：用分布函数或密度变换替代把连续单点概率等同于密度的写法。
\item Y3.21：参数 $n=19,p=0.7$ 时，二项众数为 $13$ 与 $14$ 两个。
\item A4.19：混合密度的求导结果必须保留有限求和。
\item A4.23：变换后的支撑为 $0<w<z$，边缘积分贡献额外的 $z$，正确密度为 $z^2e^{-z}/2$。
\item A5.14：按原解的积分模型校订联合密度为 $3/(2x^3y^2)$，原题面幂次在该区域不能归一化。
\item A5.35：圆周均匀与圆内均匀的坐标方差分别为 $R^2/2$ 与 $R^2/4$。
\item A6.15：补全伽马—泊松条件分布的解答，并将不恰当的连续金额泊松表述改为计数模型。
\item A6.17：第三小指数寿命密度中应为 $e^{-5\lambda y}$，不是 $e^{-2\lambda y}$。
\item A6.20、A6.24：校正密度与生成函数证明的变量和下标。
\item A7.5：切比雪夫给出的是下界，不能将最终不等号写成等号。
\item A7.20、A7.21、A7.22、A7.25：补全原资料缺失或未结束的极限定理论证。
\item C22、P10：统一下侧分位数约定，复算方差、标准差和均值区间，校正舍入差异。
\item F11：中央区间是不拒绝区域，外侧才是拒绝域。
\item PF30：正态统计量为 $10$ 时，双侧尾概率极小但严格为正，不写作精确的零。
\end{enumerate}

\clearpage
\section*{逐题主题与来源索引}
\addcontentsline{toc}{section}{逐题主题与来源索引}
右栏“题／解”依次给出题目页和详解页，均可点击跳转。Y、A 后的数字为公开中文资料原题号；M 后接上表文件缩写与原题位置；B 表示编订补充。题源为选编与校订，不保证与原资料逐字相同。
\begin{longtable}{@{}p{13mm}p{72mm}r@{}}
\toprule 本册题号 & 主题与来源键 & 题／解\\\midrule\endfirsthead
\toprule 本册题号 & 主题与来源键 & 题／解\\\midrule\endhead


\ref{q:Y1.1} & 事件分解与德摩根律\newline {\small Y1.1} & \pageref{q:Y1.1}/\pageref{a:Y1.1}\\[2pt]

\ref{q:Y1.2} & 用集合语言表示事件\newline {\small Y1.2} & \pageref{q:Y1.2}/\pageref{a:Y1.2}\\[2pt]

\ref{q:Y1.3} & 有放回抽球与样本空间\newline {\small Y1.3} & \pageref{q:Y1.3}/\pageref{a:Y1.3}\\[2pt]

\ref{q:Y1.4} & 将事件列化为互斥事件列\newline {\small Y1.4} & \pageref{q:Y1.4}/\pageref{a:Y1.4}\\[2pt]

\ref{q:Y1.5} & 两阶段抽样的对称性\newline {\small Y1.5} & \pageref{q:Y1.5}/\pageref{a:Y1.5}\\[2pt]

\ref{q:Y1.7} & 无放回试钥匙\newline {\small Y1.7} & \pageref{q:Y1.7}/\pageref{a:Y1.7}\\[2pt]

\ref{q:Y1.8} & 生日相同的概率\newline {\small Y1.8} & \pageref{q:Y1.8}/\pageref{a:Y1.8}\\[2pt]

\ref{q:Y1.9} & 恰好得到一副手套\newline {\small Y1.9} & \pageref{q:Y1.9}/\pageref{a:Y1.9}\\[2pt]

\ref{q:Y1.10} & 两次抽样的交集\newline {\small Y1.10} & \pageref{q:Y1.10}/\pageref{a:Y1.10}\\[2pt]

\ref{q:Y1.11} & 单位正方形中的圆域\newline {\small Y1.11} & \pageref{q:Y1.11}/\pageref{a:Y1.11}\\[2pt]

\ref{q:Y1.12} & 到达时差与等待\newline {\small Y1.12} & \pageref{q:Y1.12}/\pageref{a:Y1.12}\\[2pt]

\ref{q:Y1.13} & 硬币落格与边界区域\newline {\small Y1.13} & \pageref{q:Y1.13}/\pageref{a:Y1.13}\\[2pt]

\ref{q:Y1.14} & 两次服务是否相遇\newline {\small Y1.14} & \pageref{q:Y1.14}/\pageref{a:Y1.14}\\[2pt]

\ref{q:Y1.15} & 立方体内的球域\newline {\small Y1.15} & \pageref{q:Y1.15}/\pageref{a:Y1.15}\\[2pt]

\ref{q:Y1.17} & 概率的次可加性\newline {\small Y1.17} & \pageref{q:Y1.17}/\pageref{a:Y1.17}\\[2pt]

\ref{q:Y1.18} & 可数个概率一事件\newline {\small Y1.18} & \pageref{q:Y1.18}/\pageref{a:Y1.18}\\[2pt]

\ref{q:Y1.19} & 四个互斥区域的分割\newline {\small Y1.19} & \pageref{q:Y1.19}/\pageref{a:Y1.19}\\[2pt]

\ref{q:Y1.20} & 递增事件列的极限\newline {\small Y1.20} & \pageref{q:Y1.20}/\pageref{a:Y1.20}\\[2pt]

\ref{q:Y2.1} & 至少一次求职成功\newline {\small Y2.1} & \pageref{q:Y2.1}/\pageref{a:Y2.1}\\[2pt]

\ref{q:Y2.2} & 自习室占用问题\newline {\small Y2.2} & \pageref{q:Y2.2}/\pageref{a:Y2.2}\\[2pt]

\ref{q:Y2.3} & 无人拿到自己帽子\newline {\small Y2.3} & \pageref{q:Y2.3}/\pageref{a:Y2.3}\\[2pt]

\ref{q:Y2.4} & 有放回点名的覆盖概率\newline {\small Y2.4} & \pageref{q:Y2.4}/\pageref{a:Y2.4}\\[2pt]

\ref{q:Y2.5} & 信封匹配与容斥\newline {\small Y2.5} & \pageref{q:Y2.5}/\pageref{a:Y2.5}\\[2pt]

\ref{q:Y2.6} & 独立任务的共同成功\newline {\small Y2.6} & \pageref{q:Y2.6}/\pageref{a:Y2.6}\\[2pt]

\ref{q:Y2.7} & 发现信号的试验次数\newline {\small Y2.7} & \pageref{q:Y2.7}/\pageref{a:Y2.7}\\[2pt]

\ref{q:Y2.9} & 保证成功概率的人员数\newline {\small Y2.9} & \pageref{q:Y2.9}/\pageref{a:Y2.9}\\[2pt]

\ref{q:Y2.10} & 德梅雷的骰子问题\newline {\small Y2.10} & \pageref{q:Y2.10}/\pageref{a:Y2.10}\\[2pt]

\ref{q:Y2.11} & 条件成绩概率\newline {\small Y2.11} & \pageref{q:Y2.11}/\pageref{a:Y2.11}\\[2pt]

\ref{q:Y2.12} & 连续受损的生存概率\newline {\small Y2.12} & \pageref{q:Y2.12}/\pageref{a:Y2.12}\\[2pt]

\ref{q:Y2.13} & 已知一枚骰子的点数\newline {\small Y2.13} & \pageref{q:Y2.13}/\pageref{a:Y2.13}\\[2pt]

\ref{q:Y2.15} & 两孩家庭中的条件事件\newline {\small Y2.15} & \pageref{q:Y2.15}/\pageref{a:Y2.15}\\[2pt]

\ref{q:Y2.16} & 通知渠道的全概率\newline {\small Y2.16} & \pageref{q:Y2.16}/\pageref{a:Y2.16}\\[2pt]

\ref{q:Y2.17} & 隐私调查的随机化回答\newline {\small Y2.17} & \pageref{q:Y2.17}/\pageref{a:Y2.17}\\[2pt]

\ref{q:Y2.19} & 比分差达到二的比赛\newline {\small Y2.19} & \pageref{q:Y2.19}/\pageref{a:Y2.19}\\[2pt]

\ref{q:Y2.20} & 加强抽样的波利亚球袋\newline {\small Y2.20} & \pageref{q:Y2.20}/\pageref{a:Y2.20}\\[2pt]

\ref{q:Y2.21} & 已知正点反推交通方式\newline {\small Y2.21} & \pageref{q:Y2.21}/\pageref{a:Y2.21}\\[2pt]

\ref{q:Y2.22} & 观察合格品后的状态\newline {\small Y2.22} & \pageref{q:Y2.22}/\pageref{a:Y2.22}\\[2pt]

\ref{q:Y2.23} & 低先验概率下的连续检测\newline {\small Y2.23} & \pageref{q:Y2.23}/\pageref{a:Y2.23}\\[2pt]

\ref{q:Y2.25} & 目击证据与先验比例\newline {\small Y2.25} & \pageref{q:Y2.25}/\pageref{a:Y2.25}\\[2pt]

\ref{q:Y2.26} & 无穷多次成功的条件\newline {\small Y2.26} & \pageref{q:Y2.26}/\pageref{a:Y2.26}\\[2pt]

\ref{q:Y2.27} & 最终不再失败的条件\newline {\small Y2.27} & \pageref{q:Y2.27}/\pageref{a:Y2.27}\\[2pt]

\ref{q:Y2.28} & 加红球后的黑球出现次数\newline {\small Y2.28} & \pageref{q:Y2.28}/\pageref{a:Y2.28}\\[2pt]

\ref{q:Y2.29} & 黑红同时增长的抽球\newline {\small Y2.29} & \pageref{q:Y2.29}/\pageref{a:Y2.29}\\[2pt]

\ref{q:Y2.30} & 全概率公式的证明\newline {\small Y2.30} & \pageref{q:Y2.30}/\pageref{a:Y2.30}\\[2pt]

\ref{q:Y2.31} & 贝叶斯公式的证明\newline {\small Y2.31} & \pageref{q:Y2.31}/\pageref{a:Y2.31}\\[2pt]

\ref{q:Y2.32} & 串联与并联可靠性\newline {\small Y2.32} & \pageref{q:Y2.32}/\pageref{a:Y2.32}\\[2pt]

\ref{q:Y2.33} & 比赛局数与胜率\newline {\small Y2.33} & \pageref{q:Y2.33}/\pageref{a:Y2.33}\\[2pt]

\ref{q:Y2.34} & 集体迟到的提前量\newline {\small Y2.34} & \pageref{q:Y2.34}/\pageref{a:Y2.34}\\[2pt]

\ref{q:Y3.1} & 分布函数的无穷远极限\newline {\small Y3.1} & \pageref{q:Y3.1}/\pageref{a:Y3.1}\\[2pt]

\ref{q:Y3.2} & 指数型密度的归一化\newline {\small Y3.2} & \pageref{q:Y3.2}/\pageref{a:Y3.2}\\[2pt]

\ref{q:Y3.3} & 二次指数可积的条件\newline {\small Y3.3} & \pageref{q:Y3.3}/\pageref{a:Y3.3}\\[2pt]

\ref{q:Y3.4} & 最大值与最小值恒等式\newline {\small Y3.4} & \pageref{q:Y3.4}/\pageref{a:Y3.4}\\[2pt]

\ref{q:Y3.5} & 分段密度的归一化\newline {\small Y3.5} & \pageref{q:Y3.5}/\pageref{a:Y3.5}\\[2pt]

\ref{q:Y3.6} & 二项计数的尾概率\newline {\small Y3.6} & \pageref{q:Y3.6}/\pageref{a:Y3.6}\\[2pt]

\ref{q:Y3.7} & 两组二项计数的比较\newline {\small Y3.7} & \pageref{q:Y3.7}/\pageref{a:Y3.7}\\[2pt]

\ref{q:Y3.8} & 首次与第多次成功\newline {\small Y3.8} & \pageref{q:Y3.8}/\pageref{a:Y3.8}\\[2pt]

\ref{q:Y3.9} & 泊松分裂与独立性\newline {\small Y3.9} & \pageref{q:Y3.9}/\pageref{a:Y3.9}\\[2pt]

\ref{q:Y3.13} & 周期信号的均匀等待\newline {\small Y3.13} & \pageref{q:Y3.13}/\pageref{a:Y3.13}\\[2pt]

\ref{q:Y3.14} & 恒定风险率与指数寿命\newline {\small Y3.14} & \pageref{q:Y3.14}/\pageref{a:Y3.14}\\[2pt]

\ref{q:Y3.15} & 正态尺寸的合格概率\newline {\small Y3.15} & \pageref{q:Y3.15}/\pageref{a:Y3.15}\\[2pt]

\ref{q:Y3.16} & 泊松变量的平方根\newline {\small Y3.16} & \pageref{q:Y3.16}/\pageref{a:Y3.16}\\[2pt]

\ref{q:Y3.17} & 均匀电流的平方变换\newline {\small Y3.17} & \pageref{q:Y3.17}/\pageref{a:Y3.17}\\[2pt]

\ref{q:Y3.18} & 指数变量的仿射变换\newline {\small Y3.18} & \pageref{q:Y3.18}/\pageref{a:Y3.18}\\[2pt]

\ref{q:Y3.19} & 半柯西变量的对数\newline {\small Y3.19} & \pageref{q:Y3.19}/\pageref{a:Y3.19}\\[2pt]

\ref{q:Y3.20} & 倒数、绝对值与正切变换\newline {\small Y3.20} & \pageref{q:Y3.20}/\pageref{a:Y3.20}\\[2pt]

\ref{q:Y3.21} & 二项众数与参数估计\newline {\small Y3.21} & \pageref{q:Y3.21}/\pageref{a:Y3.21}\\[2pt]

\ref{q:Y3.22} & 迟到人数的众数和估计\newline {\small Y3.22} & \pageref{q:Y3.22}/\pageref{a:Y3.22}\\[2pt]

\ref{q:Y3.23} & 泊松众数与似然\newline {\small Y3.23} & \pageref{q:Y3.23}/\pageref{a:Y3.23}\\[2pt]

\ref{q:Y3.24} & 计数结果与参数的区别\newline {\small Y3.24} & \pageref{q:Y3.24}/\pageref{a:Y3.24}\\[2pt]

\ref{q:Y3.25} & 骰子样本的最大最小值\newline {\small Y3.25} & \pageref{q:Y3.25}/\pageref{a:Y3.25}\\[2pt]

\ref{q:Y3.26} & 圆周投影的反正弦密度\newline {\small Y3.26} & \pageref{q:Y3.26}/\pageref{a:Y3.26}\\[2pt]

\ref{q:Y3.28} & 正弦变换的两条分支\newline {\small Y3.28} & \pageref{q:Y3.28}/\pageref{a:Y3.28}\\[2pt]

\ref{q:Y3.29} & 逆分布函数生成正态数\newline {\small Y3.29} & \pageref{q:Y3.29}/\pageref{a:Y3.29}\\[2pt]

\ref{q:Y3.31} & 分段单调变换\newline {\small Y3.31} & \pageref{q:Y3.31}/\pageref{a:Y3.31}\\[2pt]

\ref{q:Y4.3} & 联合概率表与边缘分布\newline {\small Y4.3} & \pageref{q:Y4.3}/\pageref{a:Y4.3}\\[2pt]

\ref{q:Y4.6} & 同参数二项变量之和\newline {\small Y4.6} & \pageref{q:Y4.6}/\pageref{a:Y4.6}\\[2pt]

\ref{q:Y4.8} & 函数依赖的离散向量\newline {\small Y4.8} & \pageref{q:Y4.8}/\pageref{a:Y4.8}\\[2pt]

\ref{q:Y4.9} & 最小值的尾概率\newline {\small Y4.9} & \pageref{q:Y4.9}/\pageref{a:Y4.9}\\[2pt]

\ref{q:Y4.15} & 三角支撑上的边缘密度\newline {\small Y4.15} & \pageref{q:Y4.15}/\pageref{a:Y4.15}\\[2pt]

\ref{q:Y4.16} & 随机家庭规模与男孩数\newline {\small Y4.16} & \pageref{q:Y4.16}/\pageref{a:Y4.16}\\[2pt]

\ref{q:Y4.22} & 连续变量的离散随机缩放\newline {\small Y4.22} & \pageref{q:Y4.22}/\pageref{a:Y4.22}\\[2pt]

\ref{q:Y4.25} & 博克斯—穆勒正态生成\newline {\small Y4.25} & \pageref{q:Y4.25}/\pageref{a:Y4.25}\\[2pt]

\ref{q:Y4.30} & 指数变量的极值与比值\newline {\small Y4.30} & \pageref{q:Y4.30}/\pageref{a:Y4.30}\\[2pt]

\ref{q:Y4.34} & 平方和与平方比的变换\newline {\small Y4.34} & \pageref{q:Y4.34}/\pageref{a:Y4.34}\\[2pt]

\ref{q:Y4.36} & 线性变换的密度公式\newline {\small Y4.36} & \pageref{q:Y4.36}/\pageref{a:Y4.36}\\[2pt]

\ref{q:A4.4} & 常数与任意变量独立\newline {\small A4.4} & \pageref{q:A4.4}/\pageref{a:A4.4}\\[2pt]

\ref{q:A4.5} & 联合分布函数的右连续性\newline {\small A4.5} & \pageref{q:A4.5}/\pageref{a:A4.5}\\[2pt]

\ref{q:A4.11} & 曲边区域上的联合密度\newline {\small A4.11} & \pageref{q:A4.11}/\pageref{a:A4.11}\\[2pt]

\ref{q:A4.14} & 指数寿命的比较\newline {\small A4.14} & \pageref{q:A4.14}/\pageref{a:A4.14}\\[2pt]

\ref{q:A4.17} & 不独立变量的平方独立\newline {\small A4.17} & \pageref{q:A4.17}/\pageref{a:A4.17}\\[2pt]

\ref{q:A4.18} & 曲线下均匀点的投影\newline {\small A4.18} & \pageref{q:A4.18}/\pageref{a:A4.18}\\[2pt]

\ref{q:A4.19} & 离散与连续变量之差\newline {\small A4.19} & \pageref{q:A4.19}/\pageref{a:A4.19}\\[2pt]

\ref{q:A4.23} & 求和变换后的边缘积分\newline {\small A4.23} & \pageref{q:A4.23}/\pageref{a:A4.23}\\[2pt]

\ref{q:A4.27} & 独立正态的线性组合\newline {\small A4.27} & \pageref{q:A4.27}/\pageref{a:A4.27}\\[2pt]

\ref{q:A4.28} & 二维正态的条件分布\newline {\small A4.28} & \pageref{q:A4.28}/\pageref{a:A4.28}\\[2pt]

\ref{q:A4.32} & 尾分布幂次与竞争概率\newline {\small A4.32} & \pageref{q:A4.32}/\pageref{a:A4.32}\\[2pt]

\ref{q:A5.5} & 广告信息之间的等待\newline {\small A5.5} & \pageref{q:A5.5}/\pageref{a:A5.5}\\[2pt]

\ref{q:A5.8} & 标准正态半径的期望\newline {\small A5.8} & \pageref{q:A5.8}/\pageref{a:A5.8}\\[2pt]

\ref{q:A5.9} & 交换对称与比例期望\newline {\small A5.9} & \pageref{q:A5.9}/\pageref{a:A5.9}\\[2pt]

\ref{q:A5.12} & 离散极值的尾和公式\newline {\small A5.12} & \pageref{q:A5.12}/\pageref{a:A5.12}\\[2pt]

\ref{q:A5.14} & 曲边密度的归一化勘误\newline {\small A5.14} & \pageref{q:A5.14}/\pageref{a:A5.14}\\[2pt]

\ref{q:A5.17} & 全部指数寿命结束的等待\newline {\small A5.17} & \pageref{q:A5.17}/\pageref{a:A5.17}\\[2pt]

\ref{q:A5.20} & 样本均值与方差的无偏性\newline {\small A5.20} & \pageref{q:A5.20}/\pageref{a:A5.20}\\[2pt]

\ref{q:A5.22} & 最优逆方差加权\newline {\small A5.22} & \pageref{q:A5.22}/\pageref{a:A5.22}\\[2pt]

\ref{q:A5.25} & 非负变量的尾积分\newline {\small A5.25} & \pageref{q:A5.25}/\pageref{a:A5.25}\\[2pt]

\ref{q:A5.26} & 和与差的相关系数\newline {\small A5.26} & \pageref{q:A5.26}/\pageref{a:A5.26}\\[2pt]

\ref{q:A5.28} & 正方形密度的协方差\newline {\small A5.28} & \pageref{q:A5.28}/\pageref{a:A5.28}\\[2pt]

\ref{q:A5.30} & 正态尺寸差与装配\newline {\small A5.30} & \pageref{q:A5.30}/\pageref{a:A5.30}\\[2pt]

\ref{q:A5.32} & 不相关但不独立的反例\newline {\small A5.32} & \pageref{q:A5.32}/\pageref{a:A5.32}\\[2pt]

\ref{q:A5.33} & 中心正态的高阶矩\newline {\small A5.33} & \pageref{q:A5.33}/\pageref{a:A5.33}\\[2pt]

\ref{q:A5.35} & 圆周与圆内均匀模型\newline {\small A5.35} & \pageref{q:A5.35}/\pageref{a:A5.35}\\[2pt]

\ref{q:A6.1} & 随机组数下的条件二项\newline {\small A6.1} & \pageref{q:A6.1}/\pageref{a:A6.1}\\[2pt]

\ref{q:A6.3} & 几何个指数等待之和\newline {\small A6.3} & \pageref{q:A6.3}/\pageref{a:A6.3}\\[2pt]

\ref{q:A6.4} & 复合泊松等待的均值方差\newline {\small A6.4} & \pageref{q:A6.4}/\pageref{a:A6.4}\\[2pt]

\ref{q:A6.9} & 条件指数与条件伽马\newline {\small A6.9} & \pageref{q:A6.9}/\pageref{a:A6.9}\\[2pt]

\ref{q:A6.14} & 威布尔寿命的剩余等待\newline {\small A6.14} & \pageref{q:A6.14}/\pageref{a:A6.14}\\[2pt]

\ref{q:A6.15} & 泊松观察后的伽马后验\newline {\small A6.15} & \pageref{q:A6.15}/\pageref{a:A6.15}\\[2pt]

\ref{q:A6.17} & 七个接口的第三次失效\newline {\small A6.17} & \pageref{q:A6.17}/\pageref{a:A6.17}\\[2pt]

\ref{q:A6.20} & 指数顺序统计量的间隔\newline {\small A6.20} & \pageref{q:A6.20}/\pageref{a:A6.20}\\[2pt]

\ref{q:A6.21} & 仿射变换的概率母函数\newline {\small A6.21} & \pageref{q:A6.21}/\pageref{a:A6.21}\\[2pt]

\ref{q:A6.24} & 两组硬币成功数之差\newline {\small A6.24} & \pageref{q:A6.24}/\pageref{a:A6.24}\\[2pt]

\ref{q:A6.29} & 均匀和的特征函数极限\newline {\small A6.29} & \pageref{q:A6.29}/\pageref{a:A6.29}\\[2pt]

\ref{q:A6.30} & 指数和的特征函数极限\newline {\small A6.30} & \pageref{q:A6.30}/\pageref{a:A6.30}\\[2pt]

\ref{q:A6.34} & 二项分布的泊松极限\newline {\small A6.34} & \pageref{q:A6.34}/\pageref{a:A6.34}\\[2pt]

\ref{q:A6.35} & 给定总和后的部分和\newline {\small A6.35} & \pageref{q:A6.35}/\pageref{a:A6.35}\\[2pt]

\ref{q:A7.2} & 截尾期望的概率上界\newline {\small A7.2} & \pageref{q:A7.2}/\pageref{a:A7.2}\\[2pt]

\ref{q:A7.3} & 柯西不等式与尾概率\newline {\small A7.3} & \pageref{q:A7.3}/\pageref{a:A7.3}\\[2pt]

\ref{q:A7.4} & 最小值的概率界\newline {\small A7.4} & \pageref{q:A7.4}/\pageref{a:A7.4}\\[2pt]

\ref{q:A7.5} & 伽马矩与切比雪夫界\newline {\small A7.5} & \pageref{q:A7.5}/\pageref{a:A7.5}\\[2pt]

\ref{q:A7.10} & 三角函数样本均值之比\newline {\small A7.10} & \pageref{q:A7.10}/\pageref{a:A7.10}\\[2pt]

\ref{q:A7.14} & 比例调查的样本量\newline {\small A7.14} & \pageref{q:A7.14}/\pageref{a:A7.14}\\[2pt]

\ref{q:A7.15} & 两组样本均值差的近似\newline {\small A7.15} & \pageref{q:A7.15}/\pageref{a:A7.15}\\[2pt]

\ref{q:A7.20} & 泊松概率和的二分极限\newline {\small A7.20} & \pageref{q:A7.20}/\pageref{a:A7.20}\\[2pt]

\ref{q:A7.21} & 缩放收敛与紧性\newline {\small A7.21} & \pageref{q:A7.21}/\pageref{a:A7.21}\\[2pt]

\ref{q:A7.22} & 不同参数二项和的极限\newline {\small A7.22} & \pageref{q:A7.22}/\pageref{a:A7.22}\\[2pt]

\ref{q:A7.25} & 小扰动不改变分布极限\newline {\small A7.25} & \pageref{q:A7.25}/\pageref{a:A7.25}\\[2pt]

\ref{q:M-F-C4} & 统计量与未知参数\newline {\small M-F-C4} & \pageref{q:M-F-C4}/\pageref{a:M-F-C4}\\[2pt]

\ref{q:B8.1} & 样本均值的均值与方差\newline {\small B8.1} & \pageref{q:B8.1}/\pageref{a:B8.1}\\[2pt]

\ref{q:B8.2} & 样本方差的无偏性\newline {\small B8.2} & \pageref{q:B8.2}/\pageref{a:B8.2}\\[2pt]

\ref{q:B8.3} & 正态样本均值和残差的独立性\newline {\small B8.3} & \pageref{q:B8.3}/\pageref{a:B8.3}\\[2pt]

\ref{q:B8.4} & 卡方分布及自由度\newline {\small B8.4} & \pageref{q:B8.4}/\pageref{a:B8.4}\\[2pt]

\ref{q:B8.5} & 学生 t 统计量\newline {\small B8.5} & \pageref{q:B8.5}/\pageref{a:B8.5}\\[2pt]

\ref{q:B8.6} & 方差比与 F 分布\newline {\small B8.6} & \pageref{q:B8.6}/\pageref{a:B8.6}\\[2pt]

\ref{q:B8.7} & 顺序统计量的分布\newline {\small B8.7} & \pageref{q:B8.7}/\pageref{a:B8.7}\\[2pt]

\ref{q:B8.8} & 均匀样本最大值\newline {\small B8.8} & \pageref{q:B8.8}/\pageref{a:B8.8}\\[2pt]

\ref{q:B8.9} & 指数样本的最小值与和\newline {\small B8.9} & \pageref{q:B8.9}/\pageref{a:B8.9}\\[2pt]

\ref{q:M-F-C13} & 中心极限定理的正确标准化\newline {\small M-F-C13} & \pageref{q:M-F-C13}/\pageref{a:M-F-C13}\\[2pt]

\ref{q:B9.1} & 伯努利矩估计\newline {\small B9.1} & \pageref{q:B9.1}/\pageref{a:B9.1}\\[2pt]

\ref{q:B9.2} & 均匀分布的矩估计与极值估计\newline {\small B9.2} & \pageref{q:B9.2}/\pageref{a:B9.2}\\[2pt]

\ref{q:B9.3} & 伽马分布的两个矩参数\newline {\small B9.3} & \pageref{q:B9.3}/\pageref{a:B9.3}\\[2pt]

\ref{q:B9.4} & 正态分布的矩估计\newline {\small B9.4} & \pageref{q:B9.4}/\pageref{a:B9.4}\\[2pt]

\ref{q:B9.5} & 偏差与均方误差分解\newline {\small B9.5} & \pageref{q:B9.5}/\pageref{a:B9.5}\\[2pt]

\ref{q:B9.6} & 样本方差的最小均方误差缩放\newline {\small B9.6} & \pageref{q:B9.6}/\pageref{a:B9.6}\\[2pt]

\ref{q:B9.7} & 一致性与渐近正态性\newline {\small B9.7} & \pageref{q:B9.7}/\pageref{a:B9.7}\\[2pt]

\ref{q:B9.8} & 样本分位数的直观比较\newline {\small B9.8} & \pageref{q:B9.8}/\pageref{a:B9.8}\\[2pt]

\ref{q:M-C10-1} & 离散参数空间中的硬币似然\newline {\small M-C10-1} & \pageref{q:M-C10-1}/\pageref{a:M-C10-1}\\[2pt]

\ref{q:M-C10-2} & 指数寿命的极大似然估计\newline {\small M-C10-2} & \pageref{q:M-C10-2}/\pageref{a:M-C10-2}\\[2pt]

\ref{q:M-E2-II1} & 已知形状的伽马率估计\newline {\small M-E2-II1} & \pageref{q:M-E2-II1}/\pageref{a:M-E2-II1}\\[2pt]

\ref{q:B10.1} & 正态模型的联合极大似然\newline {\small B10.1} & \pageref{q:B10.1}/\pageref{a:B10.1}\\[2pt]

\ref{q:B10.2} & 均匀端点与支撑集\newline {\small B10.2} & \pageref{q:B10.2}/\pageref{a:B10.2}\\[2pt]

\ref{q:B10.3} & 几何分布的估计与边界\newline {\small B10.3} & \pageref{q:B10.3}/\pageref{a:B10.3}\\[2pt]

\ref{q:B10.4} & 带删失的指数寿命\newline {\small B10.4} & \pageref{q:B10.4}/\pageref{a:B10.4}\\[2pt]

\ref{q:B10.5} & 极大似然的不变性\newline {\small B10.5} & \pageref{q:B10.5}/\pageref{a:B10.5}\\[2pt]

\ref{q:B10.6} & 指数率估计的有限样本偏差\newline {\small B10.6} & \pageref{q:B10.6}/\pageref{a:B10.6}\\[2pt]

\ref{q:B11.1} & 伯努利样本的充分统计量\newline {\small B11.1} & \pageref{q:B11.1}/\pageref{a:B11.1}\\[2pt]

\ref{q:B11.2} & 泊松样本的充分性与条件二项分布\newline {\small B11.2} & \pageref{q:B11.2}/\pageref{a:B11.2}\\[2pt]

\ref{q:B11.3} & 拉奥—布莱克改进\newline {\small B11.3} & \pageref{q:B11.3}/\pageref{a:B11.3}\\[2pt]

\ref{q:B11.4} & 二项族的完全性\newline {\small B11.4} & \pageref{q:B11.4}/\pageref{a:B11.4}\\[2pt]

\ref{q:B11.5} & 估计平方概率\newline {\small B11.5} & \pageref{q:B11.5}/\pageref{a:B11.5}\\[2pt]

\ref{q:B11.6} & 伯努利模型的信息量与下界\newline {\small B11.6} & \pageref{q:B11.6}/\pageref{a:B11.6}\\[2pt]

\ref{q:B11.7} & 信息下界并非总能达到\newline {\small B11.7} & \pageref{q:B11.7}/\pageref{a:B11.7}\\[2pt]

\ref{q:B11.8} & 非正则模型中的错误下界\newline {\small B11.8} & \pageref{q:B11.8}/\pageref{a:B11.8}\\[2pt]

\ref{q:M-E2-II2} & 三类别观察的贝叶斯更新\newline {\small M-E2-II2} & \pageref{q:M-E2-II2}/\pageref{a:M-E2-II2}\\[2pt]

\ref{q:M-E2-II3} & 先验赔率与后验赔率\newline {\small M-E2-II3} & \pageref{q:M-E2-II3}/\pageref{a:M-E2-II3}\\[2pt]

\ref{q:M-F-C14} & 贝塔与二项共轭\newline {\small M-F-C14} & \pageref{q:M-F-C14}/\pageref{a:M-F-C14}\\[2pt]

\ref{q:B12.1} & 平方损失下的贝叶斯估计\newline {\small B12.1} & \pageref{q:B12.1}/\pageref{a:B12.1}\\[2pt]

\ref{q:B12.2} & 伽马与泊松共轭的多次观察\newline {\small B12.2} & \pageref{q:B12.2}/\pageref{a:B12.2}\\[2pt]

\ref{q:B12.3} & 正态均值的共轭更新\newline {\small B12.3} & \pageref{q:B12.3}/\pageref{a:B12.3}\\[2pt]

\ref{q:M-P7-1} & 主持人规则改变三门问题\newline {\small M-P7-1} & \pageref{q:M-P7-1}/\pageref{a:M-P7-1}\\[2pt]

\ref{q:M-P7-2} & 未知面数骰子的预测\newline {\small M-P7-2} & \pageref{q:M-P7-2}/\pageref{a:M-P7-2}\\[2pt]

\ref{q:M-P8-1fg} & 显著性与后验概率的区别\newline {\small M-P8-1fg} & \pageref{q:M-P8-1fg}/\pageref{a:M-P8-1fg}\\[2pt]

\ref{q:M-C22-1} & 小样本正态总体的三类区间\newline {\small M-C22-1} & \pageref{q:M-C22-1}/\pageref{a:M-C22-1}\\[2pt]

\ref{q:M-P10-1} & 已知标准差的置信水平比较\newline {\small M-P10-1} & \pageref{q:M-P10-1}/\pageref{a:M-P10-1}\\[2pt]

\ref{q:M-P10-2} & 灌装量的正态均值区间\newline {\small M-P10-2} & \pageref{q:M-P10-2}/\pageref{a:M-P10-2}\\[2pt]

\ref{q:M-P10-3} & 正态方差的区间\newline {\small M-P10-3} & \pageref{q:M-P10-3}/\pageref{a:M-P10-3}\\[2pt]

\ref{q:M-P11-1} & 身高区间与宽度要求\newline {\small M-P11-1} & \pageref{q:M-P11-1}/\pageref{a:M-P11-1}\\[2pt]

\ref{q:M-P11-2} & 顺序统计量构造端点区间\newline {\small M-P11-2} & \pageref{q:M-P11-2}/\pageref{a:M-P11-2}\\[2pt]

\ref{q:M-C22-3} & 比例误差与样本量\newline {\small M-C22-3} & \pageref{q:M-C22-3}/\pageref{a:M-C22-3}\\[2pt]

\ref{q:B13.1} & 指数率的精确区间\newline {\small B13.1} & \pageref{q:B13.1}/\pageref{a:B13.1}\\[2pt]

\ref{q:B13.2} & 均匀端点的单侧区间\newline {\small B13.2} & \pageref{q:B13.2}/\pageref{a:B13.2}\\[2pt]

\ref{q:B13.3} & 两个均值差的合并方差区间\newline {\small B13.3} & \pageref{q:B13.3}/\pageref{a:B13.3}\\[2pt]

\ref{q:B13.4} & 韦尔奇近似区间\newline {\small B13.4} & \pageref{q:B13.4}/\pageref{a:B13.4}\\[2pt]

\ref{q:B13.5} & 配对差的区间\newline {\small B13.5} & \pageref{q:B13.5}/\pageref{a:B13.5}\\[2pt]

\ref{q:M-P10-5} & 置信区间与后验区间不能互换\newline {\small M-P10-5} & \pageref{q:M-P10-5}/\pageref{a:M-P10-5}\\[2pt]

\ref{q:B13.6} & 自助法百分位区间\newline {\small B13.6} & \pageref{q:B13.6}/\pageref{a:B13.6}\\[2pt]

\ref{q:M-P8-1ab} & 硬币的精确双侧检验\newline {\small M-P8-1ab} & \pageref{q:M-P8-1ab}/\pageref{a:M-P8-1ab}\\[2pt]

\ref{q:M-P8-1cd} & 离散拒绝域与功效\newline {\small M-P8-1cd} & \pageref{q:M-P8-1cd}/\pageref{a:M-P8-1cd}\\[2pt]

\ref{q:M-P8-2} & 第一类与第二类错误\newline {\small M-P8-2} & \pageref{q:M-P8-2}/\pageref{a:M-P8-2}\\[2pt]

\ref{q:M-P8-3} & 复合零假设与功效规划\newline {\small M-P8-3} & \pageref{q:M-P8-3}/\pageref{a:M-P8-3}\\[2pt]

\ref{q:M-P8-4} & 直接计算两类错误\newline {\small M-P8-4} & \pageref{q:M-P8-4}/\pageref{a:M-P8-4}\\[2pt]

\ref{q:M-P9-1} & 单侧与双侧 t 检验\newline {\small M-P9-1} & \pageref{q:M-P9-1}/\pageref{a:M-P9-1}\\[2pt]

\ref{q:M-P9-2} & 停止规则改变抽样分布\newline {\small M-P9-2} & \pageref{q:M-P9-2}/\pageref{a:M-P9-2}\\[2pt]

\ref{q:M-E2-II4} & 正态右尾检验的分位数\newline {\small M-E2-II4} & \pageref{q:M-E2-II4}/\pageref{a:M-E2-II4}\\[2pt]

\ref{q:M-F-C10} & 检验与置信区间对偶\newline {\small M-F-C10} & \pageref{q:M-F-C10}/\pageref{a:M-F-C10}\\[2pt]

\ref{q:B14.1} & 奈曼—皮尔逊检验\newline {\small B14.1} & \pageref{q:B14.1}/\pageref{a:B14.1}\\[2pt]

\ref{q:B14.2} & 正态均值的似然比检验\newline {\small B14.2} & \pageref{q:B14.2}/\pageref{a:B14.2}\\[2pt]

\ref{q:B14.3} & 多重检验与邦费罗尼控制\newline {\small B14.3} & \pageref{q:B14.3}/\pageref{a:B14.3}\\[2pt]

\ref{q:B14.4} & 不拒绝不等于等效\newline {\small B14.4} & \pageref{q:B14.4}/\pageref{a:B14.4}\\[2pt]

\ref{q:M-P9-3} & 单个正态总体的方差检验\newline {\small M-P9-3} & \pageref{q:M-P9-3}/\pageref{a:M-P9-3}\\[2pt]

\ref{q:M-E2-II5} & 两个总体的类别比例比较\newline {\small M-E2-II5} & \pageref{q:M-E2-II5}/\pageref{a:M-E2-II5}\\[2pt]

\ref{q:B15.1} & 多项分类的拟合优度\newline {\small B15.1} & \pageref{q:B15.1}/\pageref{a:B15.1}\\[2pt]

\ref{q:B15.2} & 估计参数后的自由度\newline {\small B15.2} & \pageref{q:B15.2}/\pageref{a:B15.2}\\[2pt]

\ref{q:B15.3} & 列联表独立性检验\newline {\small B15.3} & \pageref{q:B15.3}/\pageref{a:B15.3}\\[2pt]

\ref{q:B15.4} & 小样本的费希尔精确检验\newline {\small B15.4} & \pageref{q:B15.4}/\pageref{a:B15.4}\\[2pt]

\ref{q:M-P9-5} & 两个正态方差的 F 检验\newline {\small M-P9-5} & \pageref{q:M-P9-5}/\pageref{a:M-P9-5}\\[2pt]

\ref{q:M-C26-1a} & 三个点的最小二乘直线\newline {\small M-C26-1a} & \pageref{q:M-C26-1a}/\pageref{a:M-C26-1a}\\[2pt]

\ref{q:M-C26-1b} & 二次多项式仍是线性回归\newline {\small M-C26-1b} & \pageref{q:M-C26-1b}/\pageref{a:M-C26-1b}\\[2pt]

\ref{q:M-C26-1c} & 对数尺度的指数拟合\newline {\small M-C26-1c} & \pageref{q:M-C26-1c}/\pageref{a:M-C26-1c}\\[2pt]

\ref{q:M-C26-2d} & 正态误差下的最小二乘与似然\newline {\small M-C26-2d} & \pageref{q:M-C26-2d}/\pageref{a:M-C26-2d}\\[2pt]

\ref{q:B16.1} & 正规方程与残差正交\newline {\small B16.1} & \pageref{q:B16.1}/\pageref{a:B16.1}\\[2pt]

\ref{q:B16.2} & 高斯—马尔可夫最优性\newline {\small B16.2} & \pageref{q:B16.2}/\pageref{a:B16.2}\\[2pt]

\ref{q:B16.3} & 斜率估计的方差与区间\newline {\small B16.3} & \pageref{q:B16.3}/\pageref{a:B16.3}\\[2pt]

\ref{q:B16.4} & 均值响应区间与预测区间\newline {\small B16.4} & \pageref{q:B16.4}/\pageref{a:B16.4}\\[2pt]

\ref{q:B16.5} & 单因素方差分析\newline {\small B16.5} & \pageref{q:B16.5}/\pageref{a:B16.5}\\[2pt]

\ref{q:B16.6} & 相关系数与回归解释度\newline {\small B16.6} & \pageref{q:B16.6}/\pageref{a:B16.6}\\[2pt]

\ref{q:M-F-5} & 帕累托模型从估计到区间\newline {\small M-F-5} & \pageref{q:M-F-5}/\pageref{a:M-F-5}\\[2pt]

\ref{q:M-PF-26} & 瑞利参数与变换统计量\newline {\small M-PF-26} & \pageref{q:M-PF-26}/\pageref{a:M-PF-26}\\[2pt]

\ref{q:M-F-7} & 定位中的连续更新与预测\newline {\small M-F-7} & \pageref{q:M-F-7}/\pageref{a:M-F-7}\\[2pt]

\ref{q:M-F-8} & 六类别概率的综合检验\newline {\small M-F-8} & \pageref{q:M-F-8}/\pageref{a:M-F-8}\\[2pt]

\ref{q:M-F-9} & 两样本检验的完整流程\newline {\small M-F-9} & \pageref{q:M-F-9}/\pageref{a:M-F-9}\\[2pt]

\ref{q:M-F-11} & 双侧检验在指定备择下的功效\newline {\small M-F-11} & \pageref{q:M-F-11}/\pageref{a:M-F-11}\\[2pt]

\ref{q:M-F-12} & 折线基函数的最小二乘\newline {\small M-F-12} & \pageref{q:M-F-12}/\pageref{a:M-F-12}\\[2pt]

\ref{q:M-F-13} & 自助分位数的读法\newline {\small M-F-13} & \pageref{q:M-F-13}/\pageref{a:M-F-13}\\[2pt]

\ref{q:M-PF-18} & 抽球机制与逆向条件概率\newline {\small M-PF-18} & \pageref{q:M-PF-18}/\pageref{a:M-PF-18}\\[2pt]

\ref{q:M-PF-19} & 停时分布的边界计算\newline {\small M-PF-19} & \pageref{q:M-PF-19}/\pageref{a:M-PF-19}\\[2pt]

\ref{q:M-PF-20} & 分段密度与分位数\newline {\small M-PF-20} & \pageref{q:M-PF-20}/\pageref{a:M-PF-20}\\[2pt]

\ref{q:M-PF-21} & 从分布函数回到密度与期望\newline {\small M-PF-21} & \pageref{q:M-PF-21}/\pageref{a:M-PF-21}\\[2pt]

\ref{q:M-PF-29} & 参数自助法与基本区间\newline {\small M-PF-29} & \pageref{q:M-PF-29}/\pageref{a:M-PF-29}\\[2pt]

\ref{q:M-PF-30} & 极小尾概率与数值稳定性\newline {\small M-PF-30} & \pageref{q:M-PF-30}/\pageref{a:M-PF-30}\\[2pt]

\ref{q:M-PF-31} & 正态测量的后验精度规划\newline {\small M-PF-31} & \pageref{q:M-PF-31}/\pageref{a:M-PF-31}\\[2pt]

\ref{q:B17.1} & 泊松分裂与参数推断\newline {\small B17.1} & \pageref{q:B17.1}/\pageref{a:B17.1}\\[2pt]

\ref{q:B17.2} & 有限总体抽样与修正因子\newline {\small B17.2} & \pageref{q:B17.2}/\pageref{a:B17.2}\\[2pt]

\ref{q:B17.3} & 随机化回答的估计误差\newline {\small B17.3} & \pageref{q:B17.3}/\pageref{a:B17.3}\\[2pt]

\ref{q:B17.4} & 尾和公式与收集全套\newline {\small B17.4} & \pageref{q:B17.4}/\pageref{a:B17.4}\\[2pt]

\ref{q:B17.5} & 随机和的复合泊松结构\newline {\small B17.5} & \pageref{q:B17.5}/\pageref{a:B17.5}\\[2pt]

\ref{q:B17.6} & 正态变量平方和与比值\newline {\small B17.6} & \pageref{q:B17.6}/\pageref{a:B17.6}\\[2pt]

\ref{q:B17.7} & 指数竞争与获胜标记\newline {\small B17.7} & \pageref{q:B17.7}/\pageref{a:B17.7}\\[2pt]

\ref{q:B17.8} & 条件期望作为最小平方预测\newline {\small B17.8} & \pageref{q:B17.8}/\pageref{a:B17.8}\\[2pt]

\ref{q:B17.9} & 混合正态的方差分解\newline {\small B17.9} & \pageref{q:B17.9}/\pageref{a:B17.9}\\[2pt]

\ref{q:B17.10} & 德尔塔方法估计风险函数\newline {\small B17.10} & \pageref{q:B17.10}/\pageref{a:B17.10}\\[2pt]

\ref{q:B17.11} & 威尔逊比例区间的反演\newline {\small B17.11} & \pageref{q:B17.11}/\pageref{a:B17.11}\\[2pt]

\ref{q:B17.12} & 正态模型的充分统计量\newline {\small B17.12} & \pageref{q:B17.12}/\pageref{a:B17.12}\\[2pt]

\ref{q:B17.13} & 最小值模型的非对称估计\newline {\small B17.13} & \pageref{q:B17.13}/\pageref{a:B17.13}\\[2pt]

\ref{q:B17.14} & 贝塔—二项预测的额外变异\newline {\small B17.14} & \pageref{q:B17.14}/\pageref{a:B17.14}\\[2pt]

\ref{q:B17.15} & 正态性、独立性与零相关\newline {\small B17.15} & \pageref{q:B17.15}/\pageref{a:B17.15}\\[2pt]

\ref{q:B17.16} & 支撑集破坏独立性\newline {\small B17.16} & \pageref{q:B17.16}/\pageref{a:B17.16}\\[2pt]

\ref{q:B17.17} & 蒙特卡洛积分与标准误\newline {\small B17.17} & \pageref{q:B17.17}/\pageref{a:B17.17}\\[2pt]

\ref{q:B17.18} & 反向变量的方差缩减\newline {\small B17.18} & \pageref{q:B17.18}/\pageref{a:B17.18}\\[2pt]

\ref{q:B17.19} & 切比雪夫界与正态近似的角色\newline {\small B17.19} & \pageref{q:B17.19}/\pageref{a:B17.19}\\[2pt]

\ref{q:B17.20} & 学生化与斯卢茨基定理\newline {\small B17.20} & \pageref{q:B17.20}/\pageref{a:B17.20}\\[2pt]

\ref{q:B17.21} & 三种收敛之间的辨析\newline {\small B17.21} & \pageref{q:B17.21}/\pageref{a:B17.21}\\[2pt]

\ref{q:B17.22} & 参数空间约束下的正态估计\newline {\small B17.22} & \pageref{q:B17.22}/\pageref{a:B17.22}\\[2pt]

\ref{q:B17.23} & 遗漏变量导致斜率变化\newline {\small B17.23} & \pageref{q:B17.23}/\pageref{a:B17.23}\\[2pt]

\ref{q:B17.24} & 加权最小二乘\newline {\small B17.24} & \pageref{q:B17.24}/\pageref{a:B17.24}\\[2pt]

\bottomrule\end{longtable}

\noindent 全册题目、详解、方法说明与索引组成完整学习资料。引用原题时，应同时给出原项目或课程文件及原题号；本册题号用于检索，不替代原作者署名。



\clearpage\section*{竞赛题源总索引与编订边界}

\addcontentsline{toc}{section}{竞赛题源总索引与编订边界}

本索引只统计85个竞赛题目来源条目。课程训练的264题另有上面的逐题索引，不能并入竞赛真题数量。检索日期为2026年10月10日。表中页码可点击回到相应题目；有些选题包含多个子问，仍按一个原题号计数。

\subsection*{全国大学生数学竞赛：十三个条目}

以下全部为数学类高年级决赛相关题，题号采用所引公开题解的来源元数据。组别、原题号及题面属于公开题解转录层级；每个正文题目均提供固定源码版本的逐题链接。

\begin{longtable}{@{}p{22mm}p{65mm}r@{}}\toprule 届次与题号 & 主题 & 页码\\\midrule\endhead

5届，9题 & 集券等待时间 & \pageref{pr:fin05-hi:9}\\

6届，14题 & 随机需求与期望收益 & \pageref{pr:fin06-hi:14}\\

7届，13题 & 两袋交换与两态链 & \pageref{pr:fin07-hi:13}\\

8届，13题 & 二进制随机级数 & \pageref{pr:fin08-hi:13}\\

9届，13题 & 最小值与最大值 & \pageref{pr:fin09-hi:13}\\

10届，13题 & 加权独立和的极限 & \pageref{pr:fin10-hi:13}\\

11届，12题 & 截断大数定律 & \pageref{pr:fin11-high:12}\\

12届，12题 & 稀疏红球计数 & \pageref{pr:fin12-high:12}\\

13届，12题 & 正态投影与符号回归 & \pageref{pr:fin13-high:12}\\

14届，12题 & 泊松到达与随机积分 & \pageref{pr:fin14-high:12}\\

15届，12题 & 格点随机游走的扩散极限 & \pageref{pr:fin15-high:12}\\

16届，12题 & 指数样本、Gamma与Beta & \pageref{pr:fin16-high:12}\\

17届，12题 & 有限离散分布的熵 & \pageref{pr:fin17-high:12}\\

\bottomrule\end{longtable}

\subsection*{丘成桐大学生数学竞赛：72个条目}

官方入口为\href{https://yau-contest.com/lists-jxxg.html}{相关下载}，原始材料取自\href{https://yau-contest.com/index/reldown?id=9}{2010--2026历年笔试真题压缩包}。本书未收入2010、2011年的混合科目卷或任何面试题。2012--2019年使用各年个人赛概率统计卷，2020--2026年使用各年概率统计笔试卷。没有收入某年份的全部题目时，表中明确列出实际选题号。

\begin{longtable}{@{}p{12mm}p{81mm}r@{}}\toprule 年份 & 原题号与主题 & 页码\\\midrule\endhead

2012 & 原题1：强化抽样与随机乘积 & \pageref{yau:2012-1}\\

2012 & 原题2：条件期望的正值集合与正态投影 & \pageref{yau:2012-2}\\

2012 & 原题3：联合分布的乘积分布混合 & \pageref{yau:2012-3}\\

2012 & 原题4：秩和统计量的均值与方差 & \pageref{yau:2012-4}\\

2012 & 原题5：不相关但不独立 & \pageref{yau:2012-5}\\

2012 & 原题6：指数比例与随机幂 & \pageref{yau:2012-6}\\

2013 & 原题1：尾概率与有界变换 & \pageref{yau:2013-1}\\

2013 & 原题2：两个正态变量在第一象限内的和 & \pageref{yau:2013-2}\\

2013 & 原题3：纪录次数 & \pageref{yau:2013-3}\\

2013 & 原题4：指数模型的平均 p 值 & \pageref{yau:2013-4}\\

2013 & 原题5：均匀位置族的条件覆盖率 & \pageref{yau:2013-5}\\

2013 & 原题6：离散均匀总体上端点的 UMVUE & \pageref{yau:2013-6}\\

2014 & 原题1：正态分布的积分恒等式刻画 & \pageref{yau:2014-1}\\

2014 & 原题2：正交随机级数的一个充分条件 & \pageref{yau:2014-2}\\

2014 & 原题3：改变测度后的条件期望 & \pageref{yau:2014-3}\\

2014 & 原题4：用最大值或均值预测第一个样本 & \pageref{yau:2014-4}\\

2014 & 原题5：受约束多项分布中的估计与检验 & \pageref{yau:2014-5}\\

2015 & 原题1：二次矩的最大值与正态回归 & \pageref{yau:2015-1}\\

2015 & 原题2：连续两次正面的等待时间 & \pageref{yau:2015-2}\\

2015 & 原题3：正态平方级数的可积性 & \pageref{yau:2015-3}\\

2015 & 原题4：指数样本最小值的估计性质 & \pageref{yau:2015-4}\\

2015 & 原题5：已知相对方差的正态模型 & \pageref{yau:2015-5}\\

2016 & 原题1：不立即回头的格点游走 & \pageref{yau:2016-1}\\

2016 & 原题2：随机符号矩阵的算子范数 & \pageref{yau:2016-2}\\

2016 & 原题3：非均匀投球中的二球盒数 & \pageref{yau:2016-3}\\

2016 & 原题4：均值与倒数均值的乘积 & \pageref{yau:2016-4}\\

2016 & 原题5：负漂移下的非负访问次数 & \pageref{yau:2016-5}\\

2017 & 原题1：抽去球直到仅余一种颜色 & \pageref{yau:2017-1}\\

2017 & 原题2：二叉树信号传播的阈值辨析 & \pageref{yau:2017-2}\\

2017 & 原题3：不同硬币参数的无限序列测度 & \pageref{yau:2017-3}\\

2017 & 原题4：随机图最大团的粗上下界 & \pageref{yau:2017-4}\\

2017 & 原题5：纯死亡过程的概率生成函数 & \pageref{yau:2017-5}\\

2018 & 原题1：搜索失败后的最优下一步 & \pageref{yau:2018-1}\\

2018 & 原题2：圆周游走中最后被访问的人 & \pageref{yau:2018-2}\\

2018 & 原题3：多人各带多件礼物的归还数 & \pageref{yau:2018-3}\\

2018 & 原题4：精度矩阵与条件独立 & \pageref{yau:2018-4}\\

2018 & 原题5：正态向量的投影与样本相关系数 & \pageref{yau:2018-5}\\

2019 & 原题1：正随机变量的根上极限 & \pageref{yau:2019-1}\\

2019 & 原题2：反射随机游走的正常返性 & \pageref{yau:2019-2}\\

2019 & 原题4(a)--(e)：完全随机分组的有限总体推断 & \pageref{yau:2019-4}\\

2020 & 原题1：有界零均值变量的指数矩 & \pageref{yau:2020-1}\\

2020 & 原题2：泊松分布的差分恒等式 & \pageref{yau:2020-2}\\

2020 & 原题3：偏置随机游走命中零点 & \pageref{yau:2020-3}\\

2020 & 原题4：独立坐标与旋转不变性的正态刻画 & \pageref{yau:2020-4}\\

2020 & 原题5：平衡分组再随机化的无偏性 & \pageref{yau:2020-5}\\

2020 & 原题6：不等大小组的反例 & \pageref{yau:2020-6}\\

2021 & 原题1：弱收敛与较低阶矩收敛 & \pageref{yau:2021-1}\\

2021 & 原题2：马尔可夫链的访问概率和 & \pageref{yau:2021-2}\\

2021 & 原题3：随机符号和的几乎处处增长率 & \pageref{yau:2021-3}\\

2021 & 原题4：随机零一矩阵行列式的奇偶性 & \pageref{yau:2021-4}\\

2021 & 原题5：保持随机性的协变量平衡算法 & \pageref{yau:2021-5}\\

2021 & 原题6：已知随机化机制下的估计与精确检验 & \pageref{yau:2021-6}\\

2022 & 原题1：高斯分布族对弱极限封闭 & \pageref{yau:2022-1}\\

2022 & 原题2：全变差距离的最大耦合表示 & \pageref{yau:2022-2}\\

2022 & 原题3：骰子模式11与12的等待时间 & \pageref{yau:2022-3}\\

2022 & 原题4：共同重置状态产生的收缩 & \pageref{yau:2022-4}\\

2022 & 原题5：正态线性模型中的标准化均值 & \pageref{yau:2022-5}\\

2022 & 原题6：异方差正态均值的极小极大估计 & \pageref{yau:2022-6}\\

2023 & 原题1：强化预测分布的边缘与映射 & \pageref{yau:2023-1}\\

2023 & 原题2：逐步精细的信息与条件期望极限 & \pageref{yau:2023-2}\\

2023 & 原题3：纪录数的宽于标准差的偏差 & \pageref{yau:2023-3}\\

2023 & 原题4：混合表示是否就是条件概率 & \pageref{yau:2023-4}\\

2023 & 原题5：双边停止抽样下的完全性与无偏估计 & \pageref{yau:2023-5}\\

2023 & 原题6：潜在泊松混合中的比率估计与精确检验 & \pageref{yau:2023-6}\\

2024 & 原题1：公平多方博弈中内部对局数 & \pageref{yau:2024-1}\\

2024 & 原题3：随机幂的小数位极限 & \pageref{yau:2024-3}\\

2024 & 原题5：分组极大值再取最小值的中位数估计 & \pageref{yau:2024-5}\\

2025 & 原题1(1)：单个正态观测的比例型置信区间 & \pageref{yau:2025-1}\\

2026 & 原题1：序数变量的联合分布与优胜概率 & \pageref{yau:2026-1}\\

2026 & 原题2：分块回归与残差化定理 & \pageref{yau:2026-2}\\

2026 & 原题3：偏斜正态密度与均值 & \pageref{yau:2026-3}\\

2026 & 原题4：布朗运动对称区间退出时间的二阶矩 & \pageref{yau:2026-4}\\

\bottomrule\end{longtable}

\subsection*{主要校订与选录限制}

2017年个人赛第2题的信噪比“发散”表述与其定义不相容，本书给出完整递推，结论为阈值一侧趋零、另一侧趋于有限正值，见例~\ref{yau:2017-2}。2018年个人赛第5题仅要求向量非零不足以定义样本相关系数，本书补充中心化向量非零，见例~\ref{yau:2018-5}。2022年第5题在残差自由度为2时，无偏估计可积但通常不平方可积，本书按有限方差定义单独讨论，见例~\ref{yau:2022-5}。

2019年个人赛第4题只选录(a)--(e)的完全随机分组推断；2025年第1题只选录第(1)问的覆盖概率计算。2024年第5题区分有限样本精确中位数无偏与渐近中位数无偏，并保留整数分组约束。其余细节校订包括正态退化方差、有限期望的预先验证、连续条件分布、边界参数、潜在结果和测度核的条件，不改变它们所属的来源题号。

本书的数学解答给出可检查的推导；源文件检查、符号计算与数值核对不等于赛事组委会审定，也不替代逐题的独立数学审稿。阅读进阶题时，应将题后注明的标准定理及其适用条件一起使用。

\subsection*{再编译说明}

本文件为独立的中文 TeX 源稿，仅依赖项目根目录的 researchnotes.sty 与标准 TeX 宏包，不依赖下载包、临时脚本或其他正文文件。从项目根目录或本文件所在目录均可使用 XeLaTeX 编译；交叉引用和目录至少需要两次编译。临时输出按项目约定放入根目录 tmp/概率论与数理统计/，交付 PDF 放在本文件所在目录。共享样式文件保持不变。



